\documentclass[11pt,reqno]{article}
% Typography and page geometry adapted from the CCSD reference manuscript.
\usepackage[a4paper,left=25mm,right=25mm,top=24mm,bottom=25mm,headheight=14pt]{geometry}
\usepackage{amsmath,amssymb,amsthm,mathrsfs,booktabs,longtable,array}
\usepackage[T1]{fontenc}
\usepackage{mathpazo,microtype,titlesec,fancyhdr,enumitem,etoolbox}
\usepackage{xcolor,needspace}
\definecolor{HeadingInk}{HTML}{181818}
\definecolor{ProvedGreen}{HTML}{16803C}
\definecolor{PartialAmber}{HTML}{70542A}
\definecolor{ScopeBackground}{HTML}{F5F5F5}
\linespread{1.045}
\setlength{\parindent}{1.2em}
\setlength{\parskip}{2pt}
\setlength{\emergencystretch}{3em}
\tolerance=2500
\setlength{\tabcolsep}{4pt}
\raggedbottom
\titleformat{\section}[block]{\fontsize{17}{21}\selectfont\sffamily\bfseries\filcenter}{\thesection}{0.65em}{}
\titleformat{name=\section,numberless}{\large\sffamily\bfseries}{}{0pt}{}
\titleformat{\subsection}{\normalsize\sffamily\bfseries}{\thesubsection}{0.65em}{}
\titleformat{\subsubsection}{\normalsize\bfseries}{\thesubsubsection}{0.65em}{}
\titlespacing*{\section}{0pt}{2.0ex plus 1ex minus .2ex}{1.1ex}
\titlespacing*{\subsection}{0pt}{1.6ex plus .5ex}{0.7ex}
\titlespacing*{\subsubsection}{0pt}{1.4ex plus .4ex}{0.6ex}
\pretocmd{\section}{\Needspace{6\baselineskip}}{}{}
\setlist[enumerate]{leftmargin=1.8em,itemsep=4pt,topsep=6pt}
\setlist[itemize]{leftmargin=1.6em,itemsep=3pt,topsep=5pt}
\pagestyle{fancy}
\fancyhf{}
\fancyhead[L]{\small\itshape Hodge classes on abelian sixfolds}
\fancyfoot[C]{\small\thepage}
\renewcommand{\headrulewidth}{0pt}
\AtBeginEnvironment{thebibliography}{\small}
\renewcommand{\maketitle}{%
  \thispagestyle{plain}%
  \begingroup\raggedright
  {\fontsize{22}{26}\selectfont\bfseries
  Hodge Classes on Abelian Sixfolds\\
  Constructions and a Conditional Proof\par}
  \vspace{8pt}
  {\small Research draft\hfill October 8, 2026\par}
  \vspace{12pt}\hrule\vspace{14pt}
  \endgroup}
\renewenvironment{abstract}
  {\small\noindent\textbf{Abstract.}\enspace\ignorespaces}
  {\par\vspace{10pt}}
\newtheoremstyle{emphasizedplain}
  {9pt plus 2pt minus 1pt}{9pt plus 2pt minus 1pt}
  {\itshape}{}{\bfseries}{.}{0.6em}
  {\thmname{#1}\thmnumber{ #2}\thmnote{\textbf{ (#3)}}}
\newtheoremstyle{emphasizeddefinition}
  {9pt plus 2pt minus 1pt}{9pt plus 2pt minus 1pt}
  {\normalfont}{}{\bfseries}{.}{0.6em}
  {\thmname{#1}\thmnumber{ #2}\thmnote{\textbf{ (#3)}}}
\theoremstyle{emphasizedplain}
\usepackage[hidelinks,bookmarksdepth=2]{hyperref}
\hypersetup{pdftitle={Hodge Classes on Abelian Sixfolds: Constructions and a Conditional Proof},
 pdfauthor={Research draft},pdfsubject={Rational Hodge classes and conditional constructions}}
\DeclareRobustCommand{\provedstatus}{%
  \texorpdfstring{\textcolor{ProvedGreen}{\textbf{Proved}}}{Proved}}
\DeclareRobustCommand{\partialstatus}{%
  \texorpdfstring{\textcolor{PartialAmber}{\textbf{Partially proved}}}{Partially proved}}
\renewcommand{\proofname}{\textbf{Proof}}
\newtheorem{theorem}{Theorem}[section]
\newtheorem{lemma}[theorem]{Lemma}
\newtheorem{proposition}[theorem]{Proposition}
\newtheorem{corollary}[theorem]{Corollary}
\theoremstyle{emphasizeddefinition}
\newtheorem{example}[theorem]{Example}
\theoremstyle{emphasizeddefinition}
\newtheorem{remark}[theorem]{Remark}
\newtheorem{prelemma}{Lemma}
\newtheorem{pretheorem}[prelemma]{Theorem}
\newtheorem{preproposition}[prelemma]{Proposition}
\renewcommand{\theprelemma}{P.\arabic{prelemma}}
\DeclareMathOperator{\End}{End}
\DeclareMathOperator{\Hom}{Hom}
\DeclareMathOperator{\Hg}{Hg}
\DeclareMathOperator{\MT}{MT}
\DeclareMathOperator{\Sp}{Sp}
\DeclareMathOperator{\SO}{SO}
\DeclareMathOperator{\GL}{GL}
\DeclareMathOperator{\SL}{SL}
\DeclareMathOperator{\SU}{SU}
\DeclareMathOperator{\Res}{Res}
\DeclareMathOperator{\Tr}{Tr}
\DeclareMathOperator{\Trd}{Trd}
\DeclareMathOperator{\Sym}{Sym}
\DeclareMathOperator{\NS}{NS}
\DeclareMathOperator{\CH}{CH}
\DeclareMathOperator{\rk}{rk}
\newcommand{\Q}{\mathbf Q}
\newcommand{\C}{\mathbf C}
\newcommand{\R}{\mathbf R}
\newcommand{\Z}{\mathbf Z}
\newcommand{\B}{B}
\newcommand{\D}{D}
\newcommand{\dd}{\mathrm{der}}
\newcommand{\status}[1]{%
  \par\nobreak\smallskip\noindent
  {\setlength{\fboxsep}{7pt}%
   \colorbox{ScopeBackground}{%
    \parbox{\dimexpr\linewidth-2\fboxsep\relax}{%
     \small\textbf{\textcolor{HeadingInk}{Scope and status.}} #1}}}%
  \par\nobreak\medskip}
\title{Hodge Classes on Abelian Sixfolds: Constructions and a Conditional Proof}
\author{Research draft}
\date{October 8, 2026}
\begin{document}
\maketitle
\begin{abstract}
We revisit a twenty-one-case research manuscript on complex abelian sixfolds.
We give arguments for divisor generation on all powers in the four Type I
and two Type II cases, compute the Hodge group and the seven-dimensional
middle-degree exceptional contribution in Type III, and analyze the
imaginary-quadratic Type IV case. Split Weil algebraicity yields results on
all powers, including an explicitly constructed quaternionic family.
We give examples, Hodge dimensions, and algebraic generators for all rows.
An odd-rank polarization adjustment resolves the mixed Weil plane in
the full $5+1$ case. A joint Hodge-circle analysis settles the entire
$3+3$ row. Low-dimensional inputs also settle the $3+2+1$,
$3+1+1+1$, and all surface--elliptic product rows. For quartic and sextic CM endomorphism fields we determine the group
alternatives and invariant generators. A real-quadratic polarization
adjustment also proves the entire quartic-center row. A split trace
polarization, together with a rank bound proved using a two-dimensional
rational relation space, settles every simple CM sixfold on all powers.
A quaternionic discriminant criterion
enlarges the proved Type III strata and gives a family outside every
split-subfield construction. Repeated elliptic factors settle a further
complete stratum of the $4+1+1$ row. Explicit dihedral quotients of cyclic
cover products algebraize quartic mixed classes in families in both
$4+2$ and $4+1+1$, including their special parameters. The residual
joint groups and generators in both rows are determined, including
CM fourfolds. Hodge-rank-three CM fourfolds are independent of simple
CM surfaces; every residual coupling has the same quartic mixed model.
Their
remaining algebraicity reduces to one nonzero projected cycle.
For the residual mixed models, a tensor-contraction argument also
shows that this single-cycle condition suffices on all powers.
The weighted-cover extension covers seventeen branch orbits and
adds two marked rational Hermitian types to the previous examples.
Completion of the nonconstant period maps and specialization of their
fixed Hodge subvariations cover the represented one-dimensional PEL
components, including the known nonsplit sextic-center component.
An eighteenfold cyclic-cover construction also proves a nonsplit
sextic-center family, with Weil discriminant $-2$ over $\Q(\sqrt{-3})$.
Using the newly released split-eightfold theorem of OpenAI as an external
input, we prove a stabilization and divisor-contraction theorem for
arbitrary sixfold Weil planes. It completes S7, S8, and S10, including
their all-powers assertions. OpenAI's separate CM theorem also settles
every CM-fourfold subcase in N2 and N3. The specified complete
Kuga--Satake correspondence gives a conditional construction of the
mixed cycle for every remaining non-CM N3 model. A determinant
stabilization in an auxiliary spin representation, together with the
separate degree-four spin-tensor algebraicity assertion and the CM
theorem, gives the residual N2 cycle. All twenty-one rows are therefore
covered relative to these explicit external inputs. This yields a
conditional proof of the rational Hodge conjecture for every complex
abelian sixfold. The new external preprints have not been fully
independently verified here; neither their proofs nor the arguments
proposed here are certified by this manuscript.
\end{abstract}
\clearpage
{\setlength{\parskip}{0pt}\tableofcontents}
\clearpage

\section*{Conventions, scope, and proof infrastructure}
\addcontentsline{toc}{section}{Conventions, scope, and proof infrastructure}
All abelian varieties are over $\C$, and all cycle coefficients are rational.
For an abelian variety $A$ set
\[
 \B^p(A)=H^{2p}(A,\Q)\cap H^{p,p}(A),\qquad
 \D^\bullet(A)=\Q[\NS(A)_\Q]\subseteq\B^\bullet(A).
\]
Our main target is $A$ itself with $\dim A=6$. Statements about every
$A^n$ are optional strengthenings, explicitly distinguished below.
Here $A^n$ denotes a product of $n$ copies, so $\dim A^n=6n$, not $6^n$.
For product rows, the displayed factors are simple; repeated or isogenous
factors are allowed. The twenty-one headings are coarse cases, and an
endomorphism or signature condition must not be added silently.
The Hodge conjecture in codimension $p$ asserts the surjectivity of
$\mathrm{cl}:\CH^p(A)_\Q\longrightarrow\B^p(A)$.
We write $P_A(t)=\sum_p\dim_\Q\B^p(A)t^p$; $t$ counts codimension,
not cohomological degree. The condition $\B^\bullet(A^n)=\D^\bullet(A^n)$
for every $n$ is stronger than the Hodge conjecture for $A$ itself.
We retain the original row labels S1--S11 and N1--N10. A result on a specified
stratum or family is never counted as a result on its entire coarse row.

\subsection*{Reading the status labels}
\provedstatus{} in a row heading means that a sixfold algebraicity argument
covering its whole stated row is supplied, using the cited external
results. It does not mean that every joint Hodge group has been
classified, or that this draft has passed independent review.
There are twenty-one \provedstatus{} rows,
relative to the external results specified below. In particular,
S7, S8, and S10 use the newly released split-eightfold theorem
\cite[Theorem 1.1]{OpenAIEightfold}; the additional CM strata in N2 and
N3 use the separate CM theorem \cite[Theorem 1.1]{OpenAICM}.
The residual non-CM N3 cases also use the specified full
Kuga--Satake algebraicity assertion \cite[Theorem 1.1]{OpenAIKS}.
The residual non-CM N2 cases use the auxiliary-spin proposition
of \cite{OpenAIK3Products} and the CM theorem; its precise hypotheses
and tensor are stated before Theorem~\ref{p:N2auxiliary}.
These are newly released preprints. We use their
stated results as external inputs, not as the outcome of a complete
independent verification of their proofs. Before these new external
inputs were introduced, sixteen whole rows had been established.
\begin{center}
\begin{tabular}{>{\raggedright\arraybackslash}p{0.16\textwidth}>{\raggedright\arraybackslash}p{0.43\textwidth}>{\raggedright\arraybackslash}p{0.29\textwidth}}
\toprule
Rows & Hodge structure / generators & Algebraicity on the sixfold\\
\midrule
S1--S6 & Full groups; divisor generators & Whole rows\\
S7, S8 & Groups and exceptional generators & Whole rows via split eightfolds\\
S10 & Two unitary groups; Weil plane or pairings & Whole row via split eightfolds\\
S9 & Two unitary group alternatives & Whole row, including Weil classes\\
S11 & Hodge rank five or six; Weil plane or pairings & Whole row, on all powers\\
N1, N4--N6 & Mixed generators identified & Whole rows\\
N2 & Residual groups and generators determined & Whole row, conditional on auxiliary spin and CM inputs\\
N3 & Residual groups and generators determined & Whole row, conditional on specified KS correspondence\\
N7--N9 & Explicit examples; no uniform joint-group list & Whole rows by external theorem\\
N10 & Arbitrary elliptic products & Whole row\\
\bottomrule
\end{tabular}
\end{center}

\subsection*{What the twenty-one rows classify}
For a simple sixfold let $F$ be the center of $D=\End^0(A)$, and in Type IV
write $[F:\Q]=2e_0$, $[D:F]=d^2$, and $m=\dim_D H_1(A,\Q)$.
The equality $6=e_0d^2m$ forces $d=1$, since $6$ is squarefree.
Thus the Type IV centers have degrees $2,4,6,12$. In Type I the real field
degree divides $6$. In Types II and III, $6=2em$, where $e$ is the degree of
the real center. Type III with $m=1$ is excluded for simple complex abelian
varieties; see \cite[\S3.2]{Vial}. These restrictions give the eleven original
simple rows. The ten nontrivial partitions of $6$ give the product rows.
They classify numerical endomorphism and dimension data, not all
Mumford--Tate strata, CM types, signatures, or polarization discriminants.

\subsection*{Isogenies, invariants, and descent}
\begin{prelemma}\label{p:isogeny}
Rational Hodge classes, their algebraicity, and the Hodge conjecture on all
powers are invariant under isogeny.
\end{prelemma}
\begin{proof}
If $f:A\to A'$ is an isogeny, its pullback is a cohomological ring
isomorphism over $\Q$. A quasi-inverse isogeny and the degree formulas show
that its inverse preserves the rational span of algebraic cycle classes.
Pullback and pushforward are algebraic operations. Apply the same argument
to $f^n$ for every $n$.
\end{proof}
Set $V=H^1(A,\Q)$. Then $H^\bullet(A^n,\Q)=\bigwedge^\bullet(V^{\oplus n})$,
and
\[
 \B^p(A^n)=\left(\bigwedge^{2p}(V^{\oplus n})\right)^{\Hg(A)}.
\]
Equivalently one can use $\MT(A)$ on the Tate-twisted space
$H^{2p}(A^n,\Q)(p)$. Untwisted positive-degree cohomology must not be
identified with the invariant space of $\MT(A)$: homotheties act on it.
The identity $\Hg(A^n)=\Hg(A)$ uses the diagonal action.
Furthermore
\[
 \End_{\Hg(A)}H^1(A,\Q)=\End^0(A)^{\mathrm{op}}.
\]
The opposite-algebra convention on cohomology does not affect any commutant
or dimension argument below. These facts follow from the equivalence between
abelian varieties up to isogeny and polarizable rational weight-one Hodge
structures; see \cite[\S\S1--2]{Moonen}.

\begin{prelemma}\label{p:fft}
Suppose that the complexified Hodge group acts on $V_\C$ as a product of full
symplectic groups on their standard representations, allowing repeated
copies. Then every Hodge class on every power is generated by divisor
classes. The same conclusion holds for a product of full general linear
groups acting on standard representations and their duals.
\end{prelemma}
\begin{proof}
The first fundamental theorem of invariant theory says that tensor
invariants for the symplectic group are generated by pairwise alternating
contractions; for the general linear group they are generated by evaluations
between a standard representation and its dual. The alternating projector
from tensor powers to exterior powers is equivariant in characteristic zero.
Consequently the exterior invariant algebra is generated by its degree-two
invariants: each matching gives the exterior product of its degree-two
pairings, up to the signs of the alternating projector.
For a product group the construction applies to each independent summand.
The rational degree-two invariant space is precisely $\NS(A^n)_\Q$, by the
Lefschetz $(1,1)$ theorem. Invariants commute with scalar extension for
reductive groups, and the image of each multiplication map on rational
invariants also commutes with scalar extension. Surjectivity over $\C$
therefore implies surjectivity over $\Q$. Products of divisor cycles are
algebraic. This proves the assertion without assuming that individual split
complex pairings are rational. See also \cite[Theorem 2.14]{Moonen}.
\end{proof}

\subsection*{The small-dimensional representation argument}
We use the weight-one representation theorem: in any irreducible complex
constituent of $H^1(A,\C)$, every nontrivial simple Lie-algebra factor acts
through a classical minuscule representation. A Hodge cocharacter is
nontrivial on at most one factor acting in that constituent. All rational
simple factors are generated by the rational conjugates of the Hodge
cocharacter. These statements are explained and proved in
\cite[\S3, especially (3.9)--(3.13)]{Moonen}; they are not hypotheses invented
for this manuscript. The argument also applies to a multiplicity-space
constituent after splitting an endomorphism algebra.

\begin{prelemma}\label{p:small}
An irreducible symplectic weight-one constituent of dimension $2$, $4$, or
$6$, acted on by a connected semisimple image, has full image $\Sp_2$,
$\Sp_4$, or $\Sp_6$, respectively. An irreducible orthogonal constituent of
dimension $6$ has full image $\SO_6$.
For a symplectic constituent of dimension $12$, the only possible
complex semisimple Lie-algebra images are $\mathfrak{sp}_{12}$ and
$\mathfrak{sl}_2\oplus\mathfrak{sl}_4$, the latter on
$\C^2\otimes\bigwedge^2\C^4$.
\end{prelemma}
\begin{proof}
Write an irreducible representation of a semisimple algebra as the tensor
product of irreducible representations of its simple ideals. If the tensor
product is self-dual, each factor is self-dual. The self-dual classical
minuscule representations of dimension at most $12$ have the following
possibilities, with coincidences between low-rank root systems understood:
\[
\begin{array}{c|c|c}
\text{dimension}&\text{symplectic}&\text{orthogonal}\\\hline
2& A_1\text{ standard}&\text{none}\\
4& C_2\text{ standard}&\text{none}\\
6& C_3\text{ standard}&A_3,\ \bigwedge^2\C^4\\
8& C_4\text{ standard}&B_3\text{ spin},\ D_4\text{ vector or half-spin}\\
10&C_5\text{ standard}&D_5\text{ vector}\\
12&C_6\text{ standard}&D_6\text{ vector}
\end{array}
\]
This list follows by inspecting $\bigwedge^j\C^{\ell+1}$ in Type A,
spin representations in Types B and D, and the standard representation in
Type C. In Type A self-duality requires $2j=\ell+1$; its dimension is
$\binom{2j}{j}$ and its sign is $(-1)^j$.
There is no nontrivial self-dual classical minuscule representation of odd
dimension. Thus dimension $6$ cannot factor as $2\cdot3$, and dimension $4$
can only factor as $2\cdot2$, whose invariant form is symmetric, not
alternating. In dimension $12$ the only possible nontrivial factorization is
$2\cdot6$. An alternating form then requires the six-dimensional factor to
be orthogonal, giving the displayed $A_1\times A_3$ possibility.
The listed simple images are the full classical groups on these vector
spaces. Connectedness converts Lie-algebra equality into group equality.
\end{proof}

\begin{prelemma}\label{p:goursat}
Let $G\subseteq\prod_{i=1}^e\Sp(W_i)$ be connected, with $\dim W_i=2r$
independent of $i$ and with surjective projections. If the $W_i$ are
pairwise nonisomorphic $G$-modules, then $G$ is the full product.
\end{prelemma}
\begin{proof}
The semisimple Lie-algebra form of Goursat's lemma expresses a subdirect
product of isomorphic simple algebras as a product of diagonally identified
blocks. If a block contains $i\ne j$, their standard representations become
isomorphic. Every automorphism of $\mathfrak{sp}_{2r}$ is inner (including
$r=1$), so a twisted diagonal has the same consequence. This contradicts
the hypothesis. All blocks therefore consist of one factor.
\end{proof}

\subsection*{Reusable constructions of actual examples}
A ``very general'' example below is specified by a fixed rational
polarized endomorphism datum and a nonempty period domain; the generic
endomorphism algebra is then justified. It is not an assertion of an
explicit arithmetic point or an equation for every family.

\begin{preproposition}\label{p:realconstruction}
Let $F$ be totally real of degree $e$, let $r\geq1$, and choose
$Z_\sigma\in\mathfrak H_r$ for every real embedding $\sigma:F\hookrightarrow\R$.
In $\C^{er}$ take the lattice which is the image of
$\mathcal O_F^r\oplus\mathfrak d_F^{-1\,r}$ under
\[
 (x,y)\longmapsto(\sigma(x)+Z_\sigma\sigma(y))_\sigma .
\]
The quotient is a polarized abelian variety of dimension $er$ with an
$F$-action. At a very general point its Hodge group is the restriction of
scalars of $\Sp_{2r}$ and its endomorphism algebra is $F$; in particular it
is simple.
\end{preproposition}
\begin{proof}
The trace alternating form
$\Tr_{F/\Q}(x\cdot y'-y\cdot x')$ is integral on the given lattice. The
symmetry and positive imaginary part of each $Z_\sigma$ give the Riemann
bilinear relations. Multiplication by $\mathcal O_F$ preserves the lattice
and the complex structure. The family is the Hilbert--Siegel PEL family
attached to this datum. Its connected monodromy is Zariski dense in
$\Res_{F/\Q}\Sp_{2r}$: after imposing a level, it is an arithmetic subgroup
of this group, which has no compact real factors, so Borel density applies.
Connected monodromy is contained in the generic Hodge group. The reverse
inclusion follows from the polarization and $F$-action. Equality follows.
The centralizer on the standard restriction-of-scalars module is $F$.
Extra Hodge tensors occur in a countable union of proper Hodge loci;
outside it the claimed group and endomorphism algebra hold. A nontrivial
isogeny decomposition would give a nontrivial rational idempotent in $F$,
which is impossible. The PEL and Hodge-locus facts used here are reviewed in
\cite{Moonen}.
\end{proof}

\begin{preproposition}\label{p:quaternionconstruction}
Let $D_0=(-1,3)_\Q$, with generators $i^2=-1$, $j^2=3$, $ij=-ji$.
Let $F=\Q$ or $\Q(2\cos(2\pi/7))$, and set $D_F=D_0\otimes_\Q F$.
For $n=3/[F:\Q]$, on $D_F^n$ put
\[
 \psi(x,y)=\Tr_{F/\Q}\Trd_{D_F/F}\left(\sum_k\bar x_k i y_k\right).
\]
Choose an order lattice and multiply $\psi$ by a positive integer to make
it integral. The compatible complex structures commuting with left $D_F$
form a nonempty PEL period domain $\mathfrak H_n^{[F:\Q]}$. Very general
members are simple sixfolds of Type II with endomorphism algebra $D_F$.
\end{preproposition}
\begin{proof}
The quaternion algebra $D_0$ splits at infinity and is a division algebra:
its Hilbert symbol at $3$ is $(-1,3)_3=-1$, since $-1$ is a nonsquare modulo
$3$. An odd-degree field extension cannot split a quaternion division
algebra, so $D_F$ remains division. It splits at every real place of $F$.
Quaternion conjugation shows that $\psi$ is alternating and nondegenerate.
The adjoint on the left $D_F$-action is the positive involution
$a\mapsto i^{-1}\bar a i$. A base compatible complex structure is right
multiplication by $-i$. To check positivity, at a real place identify
$i$ with $J=\left(\begin{smallmatrix}0&-1\\1&0\end{smallmatrix}\right)$
in $M_2(\R)$. Then the contribution to $\psi(x,x(-i))$ is
$\mathrm{tr}(x^T x)>0$ for $x\ne0$.
The Riemann relations therefore give abelian varieties. The centralizer of
the datum is the symplectic group on the Morita multiplicity space, and the
nonempty domain is the stated product of Siegel domains.
Arithmetic monodromy, Borel density, and the commutant argument used in
Proposition~\ref{p:realconstruction} give the generic Hodge group and
endomorphism algebra. Division of $D_F$ implies simplicity.
\end{proof}

\begin{preproposition}\label{p:products}
Fix a partition $6=d_1+\cdots+d_k$. In the product of Siegel moduli spaces
$\prod_i\mathcal A_{d_i}(N)$ choose a very general tuple $(B_1,\ldots,B_k)$,
with $N\geq3$, and put $A=\prod_i B_i$. Each $B_i$ is simple,
$\End^0(B_i)=\Q$, the factors are pairwise nonisogenous, and
\[
 \Hg(A)=\prod_i\Sp_{2d_i},\qquad
 \End^0(A)=\Q^k,\qquad
 P_A(t)=\prod_i(1+t+\cdots+t^{d_i}).
\]
All Hodge classes on all powers of $A$ are divisor-generated.
\end{preproposition}
\begin{proof}
The product universal family has independent arithmetic monodromy in each
symplectic factor. Borel density gives the full product as connected
monodromy; the polarization decompositions give the reverse containment
for the generic Hodge group. At a very general point these groups agree.
The commutant is $\Q^k$, so each factor is simple and there are no
cross-factor homomorphisms. Lemma~\ref{p:fft} proves the all-powers statement.
On one copy, the $\Sp_{2d_i}$ invariants of the exterior algebra are
$1,\theta_i,\ldots,\theta_i^{d_i}$; independence of the factors proves the
polynomial formula. On powers the additional degree-two generators are
Poincare divisor classes between repeated copies of the same $B_i$.
\end{proof}

\subsection*{Promoting algebraic Weil classes to all powers}
We retain the following established geometric input:
\begin{pretheorem}[Markman]\label{p:markman}
For a polarized abelian sixfold of Weil type over an imaginary quadratic
field, the Weil classes are algebraic if its Hermitian discriminant is
$-1$ in $\Q^\times/N_{K/\Q}K^\times$.
\end{pretheorem}
This is \cite[Theorem 1.5.1]{Markman}, also stated as the split sixfold case
of \cite[Theorem 1.2]{MarkmanSurvey}. It is used as an external theorem,
not reproduced as a new proof here. Its discriminant hypothesis is essential.
For the standard hyperbolic Hermitian form of rank $6$, that discriminant
is $(-1)^3=-1$.

\begin{pretheorem}[New external inputs]\label{p:openai-inputs}
We additionally use the following two statements, with the preprint
status specified above.
\begin{enumerate}
\item If $X$ is an abelian eightfold with an imaginary quadratic
$k$-action, compatible polarization, and hyperbolic rational
$k$-Hermitian form of signature $(4,4)$, then its entire rational
Weil plane $W_k(X)\subset H^8(X,\Q)$ is algebraic. This includes
every compatible polarization type and all special periods.
\item Every rational Hodge class on every complex CM abelian variety
is algebraic, in every dimension and codimension.
\end{enumerate}
\end{pretheorem}
These are respectively \cite[Theorem 1.1]{OpenAIEightfold} and
\cite[Theorem 1.1]{OpenAICM}. They are distinct inputs: the first
applies to split Weil-type eightfolds that need not be CM. Neither
statement is reproved in this manuscript.

\begin{prelemma}[Rational hyperbolicity]\label{p:hyperbolic-test}
Let $k/\Q$ be imaginary quadratic. A nondegenerate $k$-Hermitian
form of rank $2r$ and signature $(r,r)$ is hyperbolic if and only if
its determinant class is $[(-1)^r]$ in
$\Q^\times/N_{k/\Q}(k^\times)$.
\end{prelemma}
\begin{proof}
A hyperbolic plane has determinant $-1$, proving necessity.
For sufficiency, the underlying rational quadratic form $x\mapsto h(x,x)$
is nondegenerate of dimension $4r$ and signature $(2r,2r)$.
When $r\geq2$, Meyer's theorem for indefinite rational quadratic
forms of dimension at least five supplies a nonzero isotropic vector;
compare the same reduction in \cite[Lemma 3.2]{OpenAIWeilPowers}.
Choose a second vector pairing to $1$ with it and subtract half its
self-pairing times the first vector. The resulting two isotropic
vectors span a hyperbolic Hermitian plane. Its orthogonal complement
has signature $(r-1,r-1)$ and determinant class $[(-1)^{r-1}]$.
Iterating reduces to rank two. Write that last form as
$\langle a,b\rangle$ with $a,b\in\Q^\times$.
The determinant condition gives $-ab=N_{k/\Q}(t)$ for some
$t\in k^\times$, and $(t/a,1)$ is isotropic, since
$aN_{k/\Q}(t/a)+b=0$. It splits the final hyperbolic plane.
\end{proof}

\begin{pretheorem}[All sixfold Weil planes from the eightfold input]
\label{p:all-sixfold-weil}
Assume the split-eightfold input in Theorem~\ref{p:openai-inputs}.
Let $A$ be any abelian sixfold of Weil type over an imaginary quadratic
field $k$, with a compatible polarization and signature $(3,3)$.
Then every class in $W_k(A)\subset H^6(A,\Q)$ is algebraic,
without a discriminant restriction.
\end{pretheorem}
\begin{proof}
Choose a $k$-basis of $H_1(A,\Q)$ and let $d<0$ be the determinant
of the associated Hermitian form $h_A$.
Take $k$-CM elliptic curves $E_+$ and $E_-$ with opposite $k$-types.
Their rank-one Hermitian coefficients have opposite signs.
Positive rational rescaling of their polarizations therefore gives
$P=E_+\times E_-$ the Hermitian form $\operatorname{diag}(1,d)$.
The product $X=A\times P$ has signature $(4,4)$ and determinant
$d^2=N_{k/\Q}(d)$. Lemma~\ref{p:hyperbolic-test} makes its
Hermitian form hyperbolic. Clear denominators in the product
polarization by a positive integer multiple; hyperbolicity is unchanged.
The external eightfold theorem now makes all of $W_k(X)$ algebraic.

It remains to give an algebraic map recovering $W_k(A)$.
For the two embeddings of $k$ write
\[
 W_k(A)_\C=L_+\oplus L_-,\qquad
 W_k(P)_\C=M_+\oplus M_-.
\]
Here $L_\pm=\bigwedge^6 H^1(A,\C)_\pm$ and
$M_\pm=\bigwedge^2 H^1(P,\C)_\pm$ are one-dimensional.
The top exterior power of each embedding summand of
$H^1(X)=H^1(A)\oplus H^1(P)$ gives, in the Kunneth decomposition,
\[
 W_k(X)_\C=(L_+\otimes M_+)\oplus(L_-\otimes M_-).
\]
The surface $P$ has signature $(1,1)$, so $W_k(P)$ consists of
rational $(1,1)$-classes. Choose $0\ne\beta\in W_k(P)$ and a
rational divisor $D$ with $\mathrm{cl}(D)=\beta$, using Lefschetz
$(1,1)$. Both complex components $\beta_+,\beta_-$ are nonzero:
they are conjugate, and $\beta$ is rational and nonzero.
The cup-product pairing $M_+\otimes M_-\to H^4(P,\C)$ is
nonzero, since it is the exterior product of the two complementary
two-dimensional embedding spaces of $H^1(P,\C)$.
The same-sign pairings vanish.

Define the rational algebraic correspondence
\[
 C_\beta(\gamma)=(\operatorname{pr}_A)_*
       (\gamma\cup\operatorname{pr}_P^*\beta).
\]
On each summand $L_\pm\otimes M_\pm$ it is contraction with
$\beta_\mp$, a nonzero scalar map onto $L_\pm$.
Consequently $C_\beta:W_k(X)\to W_k(A)$ is an isomorphism over
$\Q$, as can be checked after complexification.
For $w\in W_k(A)$ choose its rational preimage $\gamma$ and
$Z\in\CH^4(X)_\Q$ with $\mathrm{cl}(Z)=\gamma$.
Then
\[
 (\operatorname{pr}_A)_*(Z\cdot\operatorname{pr}_P^*D)
       \in\CH^3(A)_\Q
\]
has class $w$. The codimension is $4+1-\dim P=3$.
Thus the argument uses an actual algebraic contraction, not merely
an abstract isomorphism of Hodge structures.
\end{proof}

\begin{prelemma}[Products of CM factors]\label{p:CMproducts-new}
Assume the CM input in Theorem~\ref{p:openai-inputs}. Every finite
product of CM abelian varieties satisfies the rational Hodge conjecture
on all its powers, without a Hodge-group independence assumption.
\end{prelemma}
\begin{proof}
Products and powers of CM abelian varieties are CM: the products of
their commutative semisimple endomorphism algebras have dimension
twice that of the corresponding product variety. Apply the external
CM theorem to each such power, and use Lemma~\ref{p:isogeny} if needed.
\end{proof}

\begin{prelemma}\label{p:polarization}
Let $A$ have dimension $6$, imaginary quadratic multiplication by $K$, and
$\Hg(A)_\C=\SL_6$ on $W\oplus W^*$. If
$W_K(A)=\bigwedge_K^6 H^1(A,\Q)$ is algebraic, then the Hodge conjecture
holds for every power $A^n$.
\end{prelemma}
\begin{proof}
For $\SL_6$ the tensor invariant theorem gives evaluation pairings and the
two determinant tensors. After alternation, evaluation pairings give
rational divisor classes, while determinant invariants in degree $6$ on
$n$ copies have multiplicity spaces $\Sym^6\C^n$ for $W$ and for $W^*$.
This last assertion also follows from the exterior Cauchy formula: the
partition $(1^6)$ on $W$ has transpose $(6)$ on the copy space.
For $a=(a_1,\ldots,a_n)\in K^n$ consider the rational homomorphism
$f_a:A^n\to A$, $\sum_i a_i x_i$. Clear denominators to obtain an actual
homomorphism; on degree $6$ this changes its pullback by a nonzero rational
scalar. The pullbacks of $W_K(A)$ contain the pure powers
$(\sigma(a_i))_{i=1}^n{}^{\otimes6}$ in each determinant multiplicity space.
The image of $K^n$ under its two embeddings is Zariski dense in
$\C^n\times\C^n$, because $K\otimes_\Q\C=\C\times\C$ and a rational
vector space is Zariski dense in its complexification. Pure powers span
$\Sym^6\C^n$ in characteristic zero. Thus these algebraic pullbacks span
all determinant generators. Products with divisor classes span every
invariant. The rational descent argument in Lemma~\ref{p:fft} completes the
proof. This argument does not infer algebraicity from invariant theory;
the algebraicity of $W_K(A)$ is its geometric input.
\end{proof}


\subsection*{Six-dimensional reductions and geometric inputs}
\begin{prelemma}\label{p:HLsix}
For a smooth projective sixfold, the rational Hodge conjecture follows
from its cases in codimensions two and three.
\end{prelemma}
\begin{proof}
Codimension one is Lefschetz $(1,1)$. For an ample rational divisor $L$,
Hard Lefschetz gives $L^2:H^4\simeq H^8$ and $L^4:H^2\simeq H^{10}$,
with their respective Hodge bidegree shifts. The inverse images of rational
Hodge classes are rational Hodge classes. If those inverse images are
algebraic, multiplication by the indicated powers of $L$ makes the
original classes algebraic. Degrees zero and twelve are generated by the
fundamental and point classes. No algebraicity of a Lefschetz inverse is
needed in this argument.
\end{proof}
\begin{pretheorem}[Low-dimensional inputs]\label{p:lowdim}
We use the following established results with rational coefficients.
\begin{enumerate}
\item A simple abelian variety of prime dimension has full Lefschetz
Hodge group and divisor-generated Hodge rings on every power
\cite[Theorem (2.7)]{MZ} (Tankeev--Ribet).
\item If $\Hom(E,X)=0$ for an elliptic curve $E$, failure of
$\Hg(X\times E)=\Hg(X)\times\Hg(E)$ requires that $E$ have CM by an
imaginary quadratic field embedding in a simple field component of
the center of $\End^0(X)$
\cite[Proposition (3.8)]{MZ}.
\item The Hodge conjecture holds for products of abelian varieties of
dimension at most two \cite[Proposition 2.18 in the preprint]{RamonMari}.
\item All Weil classes on abelian fourfolds of Weil type are algebraic,
for every imaginary quadratic field and every discriminant
\cite[proof of Corollary 1.6.1]{Markman}.
\end{enumerate}
\end{pretheorem}
These are external inputs, not new theorems of this manuscript.
In addition, the Hodge conjecture holds for every abelian fourfold
\cite[Corollary 1.6.1]{Markman}.
The second assertion will only be applied after excluding elliptic
isogeny factors; the third includes repeated and CM factors.

\begin{prelemma}[Odd-rank polarization adjustment]
\label{p:oddscale}
Suppose an imaginary-quadratic Weil sixfold is $K$-linearly isogenous to
$X\times Y$, where the $K$-rank $r$ of $H_1(X,\Q)$ is odd.
Then its Weil classes are algebraic.
\end{prelemma}
\begin{proof}
Pull back a $K$-compatible polarization of the given Weil sixfold
along the $K$-linear isogeny. Its restrictions to $X$ and $Y$ are ample
and $K$-compatible. Take the product of these restricted polarizations. Its total $K$-Hermitian form has
signature $(3,3)$ and determinant $d<0$ in a $K$-basis. Scale the
polarization on $X$ by the positive rational number $c=-d$, keeping the
other factor fixed. The new determinant is
\[
 d'=c^r d=-d^{r+1}
     =-N_{K/\Q}\bigl(d^{(r+1)/2}\bigr).
\]
Thus its discriminant is $-1$. Clear denominators in the polarization;
a common integral scaling changes the determinant by a sixth power,
which is a norm. The Weil space depends on $(A,K)$ and not on this choice
of polarization. Theorem~\ref{p:markman} and isogeny invariance finish
the proof.
\end{proof}

\begin{prelemma}[Real-quadratic polarization adjustment]\label{p:quartic-scale}
Let a sixfold $A$ be of Weil type for an imaginary quadratic field $k$.
Suppose its endomorphisms contain $K=kF$, where $F=\Q(\sqrt d)$ is
real quadratic, $d>0$, and a polarization has Rosati involution equal
to complex conjugation on $K$. Then $A$ admits a $k$-compatible
polarization of discriminant $-1$. Consequently its $k$-Weil classes
are algebraic.
\end{prelemma}
\begin{proof}
The $K$-rank of $H_1(A,\Q)$ is three. Fix a purely imaginary
$\delta\in k^\times$. Express the polarization as
$\psi=\Tr_{K/\Q}(\delta H)$ for a nondegenerate $K/F$-Hermitian
form $H$ on $K^3$; the equivalent $k$-Hermitian form is
$h=\Tr_{K/k}H$. Such $H$ is obtained by the nondegenerate trace pairing
and the Rosati identity. Put $b=\det_K H\in F^\times$.
In the $F$-basis $(1,\sqrt d)$ the trace Gram matrix is
$\operatorname{diag}(2,2d)$, of determinant $4d$. Taking the trace of
a rank-three Hermitian form therefore gives
\[
 D:=\det_k h=(4d)^3N_{F/\Q}(b).
\]
For completeness, diagonalize $H$ over $K$ with diagonal entries in
$F$ and apply the rank-one trace identity to each entry. A change of
$K$-basis changes both sides by the same norm factor, so this proves
the formula in the induced $k$-basis.

The Weil signature of $h$ is $(3,3)$, hence $D<0$ and the two real
conjugates of $b$ have opposite signs. The element
$a=\sqrt d\,b$ or its negative is therefore totally positive in $F$.
Replace the polarization by $\psi_a(x,y)=\psi(ax,y)$. The element $a$
is Rosati-self-adjoint, and its two positive real eigenvalues show that
$\psi_a$ remains positive. It is rational; an integral multiple is a
polarization. Its Hermitian form is $aH$, and
$N_{F/\Q}(a)=-dN_{F/\Q}(b)$. Thus its determinant over $k$ is
\[
 \begin{split}
 D'&=(4d)^3N(b)N(a)^3\\
   &=-64d^6N(b)^4
    =-\bigl(8d^3N(b)^2\bigr)^2.
 \end{split}
\]
This represents $-1$ modulo $N_{k/\Q}(k^\times)$.
Clearing denominators multiplies it by a sixth power, also a norm.
The Weil space is independent of the polarization, so
Theorem~\ref{p:markman} proves its algebraicity.
\end{proof}

\begin{prelemma}[A split trace polarization for CM Weil sixfolds]
\label{p:CMsplit}
Let $A$ be a simple CM sixfold with $K=\End^0(A)$, and let
$k\subset K$ be imaginary quadratic with signature $(3,3)$ on $A$.
Then $A$ admits a compatible split $k$-Weil polarization. Consequently
$W_k(A)$ consists of algebraic classes.
\end{prelemma}
\begin{proof}
Put $F=K^+$, so $[F:\Q]=6$, $K=kF$, and identify $H_1(A,\Q)$
with $K$. A polarization has Rosati involution equal to conjugation on
$K$. Fix $0\ne\delta\in k$ with $\bar\delta=-\delta$.
The polarization pairing can therefore be written
\[
 \psi(x,y)=\Tr_{K/\Q}(\delta b x\bar y),\qquad b\in F^\times.
\]
Indeed the trace pairing identifies a $K$-compatible alternating form
with an anti-invariant element of $K$, and every such element is
$\delta b$. The associated $k$-Hermitian form is
$h(x,y)=\Tr_{K/k}(b x\bar y)$. Over a fixed complex embedding of $k$
it diagonalizes with entries $\sigma(b)$ for the six embeddings
$\sigma:F\hookrightarrow\R$. Its signature $(3,3)$ means that exactly
three of these entries have each sign.

Order the six embeddings so that their prescribed signs alternate,
starting with a negative sign. By weak approximation choose
$\theta\in F$ with
$\sigma_1(\theta)<\cdots<\sigma_6(\theta)$ in this chosen order.
One elementary justification is that an embedding matrix of a rational
basis identifies $F$ with an invertible real-linear image of $\Q^6$,
which is dense in $\R^6$; the strict ordering specifies a nonempty open
set. These six values are distinct, so $\theta$ generates $F$.
Let $f$ be its monic minimal polynomial and put $b'=1/f'(\theta)$.
The derivative at the $j$th ordered root has sign $(-1)^{6-j}$.
Thus $a=b'/b$ is totally positive. The form
$\psi'(x,y)=\psi(ax,y)$ remains a rational polarization: $a$ is
Rosati-self-adjoint and positive on each complex eigenspace. An
integral positive multiple is an actual polarization on $A$.

Lagrange interpolation, followed by comparison of the coefficient of
$T^5$, gives the exact identities
\[
 \Tr_{F/\Q}\!\left(\frac{\theta^m}{f'(\theta)}\right)=0
 \quad(0\le m\le4),\qquad
 \Tr_{F/\Q}\!\left(\frac{\theta^5}{f'(\theta)}\right)=1.
\]
For $h'(x,y)=\Tr_{K/k}(b'x\bar y)$ it follows that
$L=k\langle1,\theta,\theta^2\rangle$ is totally isotropic: each pairing
of its indicated generators is one of the first five traces.
The generators are independent over $k$, since $[K:k]=6$.
The trace form is nondegenerate, so $L$ is a maximal isotropic
$k$-subspace. This proves splitting directly, without a local--global
argument for norms. Clearing denominators preserves $L$ and splitting.
The space $W_k(A)$ does not depend on the polarization, and
Theorem~\ref{p:markman} proves its algebraicity.
\end{proof}

\begin{prelemma}[A prime-dimensional factor and an elliptic curve]
\label{p:prime-ell}
Let $B$ be simple of dimension $g\in\{3,5\}$ and let $E$ be elliptic.
All Hodge classes on $B\times E$ are divisor-generated except possibly
for a single rational Weil plane in codimension $(g+1)/2$.
The exceptional case occurs precisely when $E$ has CM by a field
$k\subset\End^0(B)$ whose signature on $B$ is
$((g+1)/2,(g-1)/2)$, up to interchange. In that case
\[
 B^\bullet(B\times E)=D^\bullet(B\times E)+W_k(B\times E),
\]
where the $k$-action on $E$ is chosen to balance the signature.
Every class in the displayed space is algebraic.
\end{prelemma}
\begin{proof}
By the prime-dimensional theorem, the possible endomorphism algebras are
$\Q$, a totally real field of degree $g$, an imaginary quadratic field,
or a CM field of degree $2g$; the Hodge group is the full Lefschetz group.
Indeed the dimension formulas exclude Types II and III and force $d=1$
in Type IV. If the joint group is a product, its invariants are tensor
products of the divisor-generated factor invariants.
By Theorem~\ref{p:lowdim}(2), only CM elliptic $E$ with field
$k\subset\End^0(B)$ remains. There is at most one such $k$, since
$4$ does not divide $2g$.

Write $G=\Hg(B)$, and let $\chi:G\to U_k$ be the $k$-determinant on
$H^1(B,\Q)$. In the CM case $G=U_F$ and $\chi=N_{F/k}$.
Orient the $k$-embedding on $E$ so its Hodge-circle character is $z$.
If $k$ acts on $B$ with signature $(r,g-r)$, the determinant character
of its Hodge circle is $z^q$, with $q=2r-g$ (a simultaneous reversal of
signs is harmless). Simplicity excludes a pure signature: in that case
all of $\End_k H_1(B,\Q)$ would commute with its complex structure,
forcing an elliptic isogeny decomposition. Thus $0<|q|\leq g-2$.
The joint Hodge group is
\[
 H=\{(v,u)\in G\times U_k:\chi(v)=u^q\}.
\]
To justify equality, the displayed relation holds on the Hodge circle,
its kernel is connected, and it projects surjectively onto $G$.
Connectedness follows over $\C$ from the primitive character vector:
$\det$ on $\GL_g$, or $(1,\ldots,1)$ on the split CM torus.
The actual Hodge group also projects onto $G$, hence has dimension at
least $\dim G=\dim H$ and must equal $H$.

For quadratic $\End^0(B)$, $\ker\chi$ is $\SL_g$ after complexification.
Its invariants on $\bigwedge(W\oplus W^*)$, $\dim W=g$, are pairings
and the two volumes. Pairings have degree two and trivial $\chi$-weight;
the volumes have degree $g$ and weights $\chi^{\pm1}$.
For CM $F$, use embedding coordinates $x_1,\ldots,x_g$ selected over
$k$. A monomial invariant under $\ker N_{F/k}$ has exponent differences
$m(1,\ldots,1)$. On one copy each difference is in $\{-1,0,1\}$,
so either $m=0$ and the monomial consists of conjugate pairs, or $m=\pm1$
and it is a volume of degree $g$.

On $E$, only $H^1(E)$ has nonzero weights, namely $u^{\pm1}$.
A volume from $B$ can therefore be neutralized exactly when $|q|=1$.
The two resulting complex lines give the rational $k$-Weil plane of
$B\times E$. No further factor can multiply a volume nontrivially,
since it already uses every selected $B$ coordinate and one $E$
coordinate. Rational descent gives the stated decomposition.
For $g=3$ the plane is algebraic by the fourfold theorem; for $g=5$
apply Lemma~\ref{p:oddscale} to the rank-one elliptic summand.
\end{proof}

\begin{prelemma}[A threefold and repeated elliptic factors]
\label{p:three-ell}
Let $B$ be a simple threefold, and $P$ a product of elliptic curves.
Then $B\times P$ satisfies the Hodge conjecture. Its Hodge ring is
generated by divisors and pullbacks of the fourfold Weil plane on
$B\times E$, if there is an elliptic isogeny class $E$ in $P$ with
$k=\End^0(E)\subset\End^0(B)$ imaginary quadratic.
\end{prelemma}
\begin{proof}
Collect elliptic isogeny classes. All classes not matching such a $k$
split off in the Hodge group by successive applications of
Theorem~\ref{p:lowdim}(2); for a new elliptic curve its CM field would
have to embed in a simple component of the center of the existing
endomorphism algebra. Distinct elliptic classes have distinct CM fields,
and the threefold contains at most one imaginary quadratic field.
Non-CM elliptic classes split off by the same proposition.
If there is no matching field, factorwise divisor generation proves the
claim. Otherwise it suffices to treat $B\times E^m$.

The joint group is the group from Lemma~\ref{p:prime-ell} with $g=3$
and $|q|=1$, acting identically on the $m$ copies of $E$.
The kernel-of-determinant invariants on the single $B$ factor consist of
pairings and the two degree-three volumes, by that lemma's proof.
For a volume, central neutrality requires elliptic weight $\mp1$.
An elliptic exterior monomial of that weight is a product of one
unpaired elliptic coordinate and weight-zero pairings. Therefore all
invariants are generated by divisors and a volume times one elliptic
coordinate. Pullbacks of $W_k(B\times E)$ under
$(b,e_1,\ldots,e_m)\mapsto(b,\sum a_i e_i)$, $a_i\in k$, span these
last generators: after scalar extension the two embeddings of $k^m$
are independent and Zariski dense. Clear denominators in the maps.
The fourfold theorem makes all these pullbacks algebraic, and rational
descent finishes the proof. Only one copy of $B$ was used.
\end{proof}

\begin{prelemma}[A surface separates from a threefold and elliptic curves]
\label{p:surface-three}
If $S$ is a simple abelian surface, $B$ a simple threefold, and $P$ a
product of elliptic curves, then
$\Hg(S\times B\times P)=\Hg(S)\times\Hg(B\times P)$.
\end{prelemma}
\begin{proof}
If $S$ is not CM, apply \cite[Proposition (2.4) and Lemma (3.4)]{MZ}:
its Hodge Lie algebra is simple over $\Q$, of noncompact type, with
unique irreducible length-one representation. There is no homomorphism
from $S$ to $B\times P$, since their simple factors have different dimensions.

If $S$ is CM, its Hodge group is the simple rank-two torus $U_F$ for a
primitive quartic CM field $F$. Its splitting-field Galois group is
$C_4$ or $D_4$ (the latter of order eight), by \cite[\S4]{MZ}.
A failure of the asserted product equality would, by Goursat, make
$U_F$ an isogeny quotient of the center of $\Hg(B\times P)$.
Adding elliptic factors only adds rank-one quadratic Galois constituents,
so it suffices to examine the center for $B$.
For real multiplication this center is finite, and for quadratic
imaginary multiplication its connected part has rank one.
For CM $B$ the torus $U_L$ has rank three and, by
\cite[Proposition (2.4)]{MZ}, is $\Q$-simple unless the sextic field $L$
contains an imaginary quadratic $k$. A simple rank-three torus has no
rank-two quotient. In the remaining case the normal closure of $L$ has
Galois group $C_2\times C_3$ or $C_2\times S_3$, since $L=kL^+$.
Neither group has a quotient isomorphic to $C_4$ or $D_4$.
Thus $U_F$ cannot be an isogeny constituent of $U_L$ either.
Here we use the equivalence, up to isogeny, between rational tori and
rational finite-image Galois modules; these modules are semisimple and
have the same splitting-field kernel for isogenous tori.
This excludes every possible common quotient.
\end{proof}


\begin{prelemma}[Small unitary constituents]\label{p:smallunitary}
Let $A$ be a simple sixfold with $\End^0(A)=K$ a CM field of degree
$2e$, where $e=2$ or $3$, and put $q=6/e$. Then
\[
 \Hg(A)_\C^{\dd}=\SL_q^e.
\]
Its determinant image $T\subset U_K$ has rank either $e$ or $e-1$.
In the latter case there is an imaginary quadratic subfield $k\subset K$
with $T=\ker(N_{K/k}:U_K\to U_k)$, and $A$ is of $k$-Weil type.
\end{prelemma}
\begin{proof}
Put $G=\Hg(A)$. Over $\C$ write
$H^1(A)=\bigoplus_{j=1}^e(W_j\oplus W_j^*)$ with $\dim W_j=q$.
The commutant is $K\otimes\C=\C^{2e}$, so these $2e$ constituents
are irreducible and pairwise inequivalent. The connected center acts by
scalars on each, and irreducibility therefore also holds for $G^{\dd}$.
In dimensions three and two the minuscule list forces the projected
Lie algebra to be $\mathfrak{sl}_3$ and $\mathfrak{sl}_2$, respectively.
A subdirect product consists of diagonally identified blocks, allowing
the duality automorphism for $\mathfrak{sl}_3$.

For $e=2$, suppose the two projected $\mathfrak{sl}_3$'s are identified.
Choose the orientation in each conjugate pair so that $W_1,W_2$ both
carry its standard module. The derived Hodge-circle projection is
nontrivial: otherwise the Hodge circle and all its rational conjugates
would be central, contrary to the definition of $G$. On both modules
this projection has the same eigenvalues. Adding a scalar central
weight must yield exactly the two weight-one Hodge-circle weights
$+1,-1$. The required scalar is unique, hence the central weights on
$W_1,W_2$ coincide. This equality is Galois-stable: a Galois automorphism
acts on the common $A_2$ diagram either trivially, preserving both
standard constituents (possibly interchanging them), or by duality,
replacing both by their duals. Writing their central characters as
$\chi_1,\chi_2$, the relation is $\chi_1\chi_2^{-1}=1$; each
Galois conjugate is the same relation or its inverse. Its kernel is
therefore a rational subtorus containing the central projection of
the Hodge circle. By the minimality defining the Hodge group it
contains every central Hodge-circle conjugate. Consequently the connected central torus generated by the conjugates
of the Hodge circle also acts equally on $W_1,W_2$. These modules
would be isomorphic for $G$, contradicting the commutant. The two
$\SL_3$ factors are therefore independent.

For $e=3$, the diagonal-block partition of the three $A_1$ factors is
Galois-invariant and the real embeddings are transitive. Since three
is prime, either all three factors are independent or they form one
diagonal block. Write the signatures as $(r_j,2-r_j)$.
If all $r_j=1$, the determinant Hodge circle is trivial; its conjugates
then force $G$ to be semisimple. Each $W_j$ is self-dual as an $A_1$
module, so $W_j\cong W_j^*$ for $G$, again contradicting the commutant.
If all $r_j$ are pure ($0$ or $2$), the Hodge circle is scalar in each
block and $G$ is a torus, contradicting its irreducible two-dimensional
constituents. Thus a pure place and a balanced place both occur.
At the former the derived Hodge-circle projection is zero and at the
latter it is nonzero; this excludes one diagonal block. The derived
group is therefore $\SL_2^3$.

The determinant cocharacter has coordinates $2r_j-q$.
For $e=2$, both coordinates are odd and nonzero, so $\rk T\geq1$.
For $e=3$, after division by two the coordinates are in $\{0,1,-1\}$,
with both zero and nonzero entries. A Galois-stable line in the signed
permutation representation cannot contain such a vector: transitivity
would force its zero-coordinate set to be both empty and nonempty.
Thus $\rk T\geq2$. Since $\rk U_K=e$, only the asserted ranks remain.
If $T$ has codimension one, \cite[Lemma (3.7)]{MZ} identifies it with
$\ker N_{K/k}$ for an imaginary quadratic subfield $k$.
The composition of the Hodge circle with the $k$-determinant is then
trivial; its exponent is the difference of the two $k$-multiplicities.
These multiplicities sum to six, so they are $(3,3)$.
\end{proof}

\part{The eleven simple rows}
\section{S1: Type I(1) (\provedstatus)}\label{sec:S1}

\status{Full row: divisor generation on all powers, under exactly
$\End^0(A)=\Q$. The argument below strengthens the old very-general claim.}
\begin{theorem}\label{t:S1}
If $A$ is a simple abelian sixfold and $\End^0(A)=\Q$, then
$\Hg(A)=\Sp_{12}$, and $\B^\bullet(A^n)=\D^\bullet(A^n)$ for every $n\geq1$.
On $A$ itself $P_A(t)=1+t+\cdots+t^6$.
\end{theorem}
\begin{proof}
The commutant of $\Hg(A)$ on $V_\C$ is $\C$, so $V_\C$ is absolutely
irreducible. Its connected center acts by scalars, and scalar matrices in
$\Sp_{12}$ form a finite group. Thus $\Hg(A)$ is semisimple.
Lemma~\ref{p:small} leaves either the full symplectic algebra or the image
$\mathfrak{sl}_2\oplus\mathfrak{sl}_4$ on
$\C^2\otimes\bigwedge^2\C^4$.
In the latter case the two nonisomorphic complex simple ideals are each
Galois-stable, hence are defined over $\Q$ (as normal factors, up to a finite
central isogeny). A nontrivial cocharacter on an irreducible representation
of a semisimple factor has at least two distinct weights: otherwise its
image is scalar, which is impossible for a nontrivial one-parameter
subgroup of a finite center. If the Hodge circle were nontrivial on both
factors, the weights on the tensor product would include at least three
distinct sums. Weight one permits only two. Consequently the circle is
trivial on one rational normal factor. Its smallest rational algebraic
containing group would omit that factor, contradicting the defining
minimality of $\Hg(A)$. The tensor possibility is excluded.
Connectedness now gives $\Hg(A)=\Sp_{12}$. Lemma~\ref{p:fft} proves divisor
generation on every power. On one standard symplectic module the only
exterior invariants are the successive powers of the polarization, giving
$P_A$.
\end{proof}
\begin{example}
Use Proposition~\ref{p:realconstruction} with $F=\Q$, $r=6$.
Explicitly $A_Z=\C^6/(\Z^6+Z\Z^6)$ for $Z\in\mathfrak H_6$ very general.
The construction proves simplicity and $\End^0(A_Z)=\Q$. The seven
one-dimensional Hodge spaces are generated by $1,\theta,\ldots,\theta^6$;
on powers add the Poincare divisors between the copies.
\end{example}
\section{S2: Type I(2) (\provedstatus)}\label{sec:S2}

\status{Full row: divisor generation on all powers.}
\begin{theorem}
Let $A$ be a simple sixfold with $\End^0(A)=F$ totally real of degree $2$.
Then
\[
 \Hg(A)=\Res_{F/\Q}\Sp_{6},\qquad
 P_A(t)=(1+t+\cdots+t^{3})^{2},
\]
and $\B^\bullet(A^n)=\D^\bullet(A^n)$ for every $n$.
\end{theorem}
\begin{proof}
The polarization is the trace of a nondegenerate $F$-alternating form.
Thus the Hodge group embeds in the indicated restriction of scalars.
Over $\C$, $V$ splits into $2$ summands $V_\sigma$ of dimension $6$.
Its commutant is $F\otimes\C=\C^{2}$; reductivity therefore implies that
each $V_\sigma$ is irreducible and the summands are pairwise nonisomorphic.
The center of the image on each summand is scalar and hence finite inside
its symplectic group. Lemma~\ref{p:small} gives surjectivity onto every
symplectic factor. Lemma~\ref{p:goursat} excludes diagonal identifications,
proving equality with the full product. Rational descent determines the
stated restriction-of-scalars group. Lemma~\ref{p:fft} proves the all-powers
conclusion. The exterior invariants on a single $V_\sigma$ are
$1,\theta_\sigma,\ldots,\theta_\sigma^{3}$, proving the polynomial.
\end{proof}
\begin{example}
Set $F=\Q(\sqrt5)$ and $r=3$ in
Proposition~\ref{p:realconstruction}. This specifies the lattice, period
matrices, and polarization and proves that a very general member belongs
to this simple row. Choose a rational basis $a_1,\ldots,a_{2}$ of $F$.
The rational $(1,1)$ classes defined by the forms
$\psi(a_jx,y)$ are divisor classes and generate $\B^\bullet(A)$.
After splitting $F$, the ring is generated by $\theta_\sigma$ with
relations $\theta_\sigma^{4}=0$; its dimension in codimension $p$ is the
coefficient of $t^p$ in the polynomial above. Endomorphism-translated
Poincare divisors complete the generators on powers.
\end{example}
\section{S3: Type I(3) (\provedstatus)}\label{sec:S3}

\status{Full row: divisor generation on all powers.}
\begin{theorem}
Let $A$ be a simple sixfold with $\End^0(A)=F$ totally real of degree $3$.
Then
\[
 \Hg(A)=\Res_{F/\Q}\Sp_{4},\qquad
 P_A(t)=(1+t+\cdots+t^{2})^{3},
\]
and $\B^\bullet(A^n)=\D^\bullet(A^n)$ for every $n$.
\end{theorem}
\begin{proof}
The polarization is the trace of a nondegenerate $F$-alternating form.
Thus the Hodge group embeds in the indicated restriction of scalars.
Over $\C$, $V$ splits into $3$ summands $V_\sigma$ of dimension $4$.
Its commutant is $F\otimes\C=\C^{3}$; reductivity therefore implies that
each $V_\sigma$ is irreducible and the summands are pairwise nonisomorphic.
The center of the image on each summand is scalar and hence finite inside
its symplectic group. Lemma~\ref{p:small} gives surjectivity onto every
symplectic factor. Lemma~\ref{p:goursat} excludes diagonal identifications,
proving equality with the full product. Rational descent determines the
stated restriction-of-scalars group. Lemma~\ref{p:fft} proves the all-powers
conclusion. The exterior invariants on a single $V_\sigma$ are
$1,\theta_\sigma,\ldots,\theta_\sigma^{2}$, proving the polynomial.
\end{proof}
\begin{example}
Set $F=\Q(2\cos(2\pi/7))$ and $r=2$ in
Proposition~\ref{p:realconstruction}. This specifies the lattice, period
matrices, and polarization and proves that a very general member belongs
to this simple row. Choose a rational basis $a_1,\ldots,a_{3}$ of $F$.
The rational $(1,1)$ classes defined by the forms
$\psi(a_jx,y)$ are divisor classes and generate $\B^\bullet(A)$.
After splitting $F$, the ring is generated by $\theta_\sigma$ with
relations $\theta_\sigma^{3}=0$; its dimension in codimension $p$ is the
coefficient of $t^p$ in the polynomial above. Endomorphism-translated
Poincare divisors complete the generators on powers.
\end{example}
\section{S4: Type I(6) (\provedstatus)}\label{sec:S4}

\status{Full row: divisor generation on all powers.}
\begin{theorem}
Let $A$ be a simple sixfold with $\End^0(A)=F$ totally real of degree $6$.
Then
\[
 \Hg(A)=\Res_{F/\Q}\Sp_{2},\qquad
 P_A(t)=(1+t+\cdots+t^{1})^{6},
\]
and $\B^\bullet(A^n)=\D^\bullet(A^n)$ for every $n$.
\end{theorem}
\begin{proof}
The polarization is the trace of a nondegenerate $F$-alternating form.
Thus the Hodge group embeds in the indicated restriction of scalars.
Over $\C$, $V$ splits into $6$ summands $V_\sigma$ of dimension $2$.
Its commutant is $F\otimes\C=\C^{6}$; reductivity therefore implies that
each $V_\sigma$ is irreducible and the summands are pairwise nonisomorphic.
The center of the image on each summand is scalar and hence finite inside
its symplectic group. Lemma~\ref{p:small} gives surjectivity onto every
symplectic factor. Lemma~\ref{p:goursat} excludes diagonal identifications,
proving equality with the full product. Rational descent determines the
stated restriction-of-scalars group. Lemma~\ref{p:fft} proves the all-powers
conclusion. The exterior invariants on a single $V_\sigma$ are
$1,\theta_\sigma,\ldots,\theta_\sigma^{1}$, proving the polynomial.
\end{proof}
\begin{example}
Set $F=\Q(\zeta_{13}+\zeta_{13}^{-1})$ and $r=1$ in
Proposition~\ref{p:realconstruction}. This specifies the lattice, period
matrices, and polarization and proves that a very general member belongs
to this simple row. Choose a rational basis $a_1,\ldots,a_{6}$ of $F$.
The rational $(1,1)$ classes defined by the forms
$\psi(a_jx,y)$ are divisor classes and generate $\B^\bullet(A)$.
After splitting $F$, the ring is generated by $\theta_\sigma$ with
relations $\theta_\sigma^{2}=0$; its dimension in codimension $p$ is the
coefficient of $t^p$ in the polynomial above. Endomorphism-translated
Poincare divisors complete the generators on powers.
\end{example}
\section{S5: Type II(1) (\provedstatus)}\label{sec:S5}

\status{Full row: divisor generation on all powers. The old $\SO_6$
model and the claimed Type II volume obstruction are replaced.}
\begin{theorem}
Let $A$ be a simple sixfold whose endomorphism algebra is an indefinite
quaternion division algebra over $\Q$. Then $\Hg(A)_\C=\Sp_6$ on
$V_\C=U\otimes W$, $\dim U=2$, $\dim W=6$, where the group acts trivially
on $U$. Every Hodge class on every power is divisor-generated, and
\[
 P_A(t)=1+3t+6t^2+10t^3+6t^4+3t^5+t^6.
\]
\end{theorem}
\begin{proof}
For Type II the positive involution on $D$ induces a symmetric form on
$U$. The alternating polarization on $U\otimes W$ therefore induces an
alternating form on $W$, so the connected Lefschetz group has complex form
$\Sp(W)$; see \cite[\S3.2]{Vial}.
The commutant $D\otimes\C=M_2(\C)$ implies $\End_{\Hg_\C}(W)=\C$,
so $W$ is irreducible. The center of the image on $W$ is finite by
Schur's lemma and the symplectic form. Lemma~\ref{p:small} gives full image
$\Sp_6$. Lemma~\ref{p:fft} applies to the $2n$ copies of $W$ on $A^n$.
For the dimension formula use the exterior Cauchy identity
\[
 \bigwedge^{2p}(U\otimes W)
   =\bigoplus_{\lambda\vdash2p} S_\lambda U\otimes S_{\lambda'}W.
\]
A summand has a symplectic invariant precisely when
$\lambda=(2a,2b)$, $3\geq a\geq b\geq0$. Its invariant multiplicity is
$\dim S_{(2a,2b)}U=2a-2b+1$, and $p=a+b$.
Summing these numbers gives the displayed coefficients.
\end{proof}
\begin{example}
Take $D=(-1,3)_\Q$, $n=3$, and the explicit trace polarization of
Proposition~\ref{p:quaternionconstruction}. Very general compatible
complex structures give simple Type II(1) sixfolds. The three-dimensional
$\NS(A)_\Q$ corresponds to the self-adjoint subspace of the orthogonal
Rosati involution. Its divisor classes generate all Hodge spaces with
respective dimensions $(1,3,6,10,6,3,1)$.
\end{example}
\section{S6: Type II(3) (\provedstatus)}\label{sec:S6}

\status{Full row: divisor generation on all powers. The split group is
$\SL_2^3$, not a torus; a nontrivial torus cannot be a derived group.}
\begin{theorem}
Let $A$ be a simple sixfold of Type II with real center $F$ of degree $3$.
The complexified Hodge group is $\SL_2^3$, acting on
$\bigoplus_\sigma U_\sigma\otimes W_\sigma$ with
$\dim U_\sigma=\dim W_\sigma=2$. All Hodge classes on all powers are
divisor-generated, and $P_A(t)=(1+3t+t^2)^3$.
\end{theorem}
\begin{proof}
Here $\dim_\Q D=12=\dim_\Q H_1(A)$, so the left $D$-rank is one.
The Type II involution gives the symplectic, rather than orthogonal,
form on each $W_\sigma$. The complex commutant is
$D\otimes\C=\prod_{\sigma=1}^3M_2(\C)$.
Therefore each $W_\sigma$ is irreducible and different multiplicity-space
modules are nonisomorphic. Each projected group is $\Sp_2=\SL_2$, by
Lemma~\ref{p:small}. A diagonal identification would create nonzero
cross-summand intertwiners, enlarging that commutant. Lemma~\ref{p:goursat}
therefore gives the full product. Lemma~\ref{p:fft} proves the conclusion.
The calculation in S5 with $\dim W=2$ instead of $6$ gives the single-factor
polynomial $1+3t+t^2$; take its third power.
\end{proof}
\begin{example}
Set $F=\Q(2\cos(2\pi/7))$ and $D=(-1,3)_\Q\otimes F$. This is an
indefinite quaternion division algebra by the odd-degree argument.
Proposition~\ref{p:quaternionconstruction} with $n=1$ gives a nonempty
three-dimensional family and simple very general members.
Here $\rho(A)=9$ and the Hodge-space dimensions are
$(1,9,30,45,30,9,1)$. The nine rational divisor generators come from a
basis of the Rosati-fixed subspace of $D$; Poincare divisors and $D$
endomorphisms provide generators on powers.
\end{example}
\section{S7: Type III(1) (\provedstatus)}\label{sec:S7}

\status{The whole row is proved on all powers using the new external
split-eightfold input and Theorem~\ref{p:all-sixfold-weil}.
Quaternionic endomorphisms generate the seven-dimensional exceptional
space from a Weil plane. The older split-subfield criterion and its
obstructed example are retained to distinguish the two constructions.}
\begin{theorem}\label{t:S7structure}
For a simple Type III(1) sixfold,
$V_\C=U\otimes W$, $\dim U=2$, $\dim W=6$, and
$\Hg(A)_\C=\SO(W)$. Moreover
\[
 P_A(t)=1+t+t^2+8t^3+t^4+t^5+t^6,
 \qquad\dim_\Q\bigl(\B^3(A)/\D^3(A)\bigr)=7.
\]
On powers, the invariant algebra is generated by degree-two pairings and
six-linear determinant invariants on the multiplicity space.
\end{theorem}
\begin{proof}
For the definite quaternion algebra the Rosati involution is canonical
quaternion conjugation. Its split alternating form on $U$ leaves a
symmetric form on $W$ in the alternating polarization of $U\otimes W$.
Thus the connected Lefschetz group is $\SO_6$, not $\Sp_6$.
The commutant $M_2(\C)$ makes $W$ an irreducible constituent.
The center of its orthogonal image is finite.
Lemma~\ref{p:small} forces its semisimple image to be $\SO_6$.
For $\mathrm O(W)$ the invariant theorem gives pairings; for $\SO(W)$ it
adds the volume tensor $\bigwedge^6 W$.
In the exterior Cauchy formula, orthogonal invariants have
$\lambda=(a,a)$ on $U$, $0\leq a\leq6$, with one-dimensional multiplicity
$S_{(a,a)}U$. These give one divisor-power invariant in each even degree.
The volume contribution first occurs in degree $6$: its partition on $W$
is $(1^6)$ and its multiplicity on $U$ is $\Sym^6 U$, of dimension $7$.
In degree $6$ these classes are independent from the orthogonal invariants.
There are no additional contributions below degree $6$; the hard Lefschetz
isomorphisms between degrees $2p$ and $12-2p$ determine the higher degrees.
This proves the polynomial. The tensor invariant theorem followed by
alternation gives the asserted generators on arbitrary powers.
\end{proof}
\begin{theorem}[Complete Type III row]\label{t:S7complete}
Assuming the split-eightfold input, every simple Type III(1) sixfold
satisfies the rational Hodge conjecture on every power.
\end{theorem}
\begin{proof}
Choose a nonzero pure quaternion in the definite algebra $D$.
Its square is negative rational, so it generates an imaginary
quadratic field $k\subset D$. The positive Rosati involution on
$D$ is canonical conjugation, hence this action is compatible with
the polarization. The action on $H^{1,0}(A)$ extends to
$D\otimes_\Q\C=M_2(\C)$; this six-dimensional module is the
sum of three standard modules. The two $k$-eigenvalues therefore
have multiplicities $(3,3)$.
Theorem~\ref{p:all-sixfold-weil} makes $W_k(A)$ algebraic.
In the notation of Theorem~\ref{t:S7structure}, its two complex
lines have multiplicity vectors $u_+^6,u_-^6$ in $\Sym^6 U$.
For every $n$, pull these classes back under rational homomorphisms
$A^n\to A$, $(x_i)\mapsto\sum_i d_i x_i$, with $d_i\in D$.
After clearing denominators these are algebraic maps; the induced
nonzero scalar changes do not affect rational cycle spans.
The tuples $D^n$ are Zariski dense in $(M_2(\C))^n$.
For a fixed $u_+\ne0$, the vectors $(d_i u_+)_i$ are therefore
Zariski dense in $U^{\oplus n}$, and their sixth powers span
$\Sym^6(U^{\oplus n})$. These are all volume generators on
$A^n$. The other generators in Theorem~\ref{t:S7structure} are
degree-two pairings, which are divisor classes. Products and
rational descent prove the assertion.
\end{proof}

\begin{theorem}[The earlier split-subfield case]\label{t:S7split}
In the setting of Theorem~\ref{t:S7structure}, suppose that a subfield
$K\subset D$ makes the polarized variety a sixfold of split Weil type.
Then all Hodge classes on all powers of $A$ are algebraic.
\end{theorem}
\begin{proof}
Theorem~\ref{p:markman} makes $W_K(A)$ algebraic. After splitting $D$,
$K$ selects two eigenlines $\ell_+,\ell_-$ in $U$. Its complex Weil
classes are the determinant tensors with multiplicity vectors
$\ell_+^6$ and $\ell_-^6$.
On $A^n$ use rational homomorphisms
$f_d:A^n\to A$, $\sum_i d_i x_i$, with $d_i\in D$, clearing denominators.
The tuples $D^n$ are Zariski dense in $(M_2(\C))^n$ after complexification.
For a fixed nonzero $u\in\ell_+$, the vectors $(d_i u)_i$ therefore range
Zariski densely in $U^{\oplus n}$. Their sixth pure powers span
$\Sym^6(U^{\oplus n})$.
Consequently the pullbacks of the algebraic Weil plane span every volume
invariant on $A^n$, not just the two on $A$. Pairings are divisor classes.
Theorem~\ref{t:S7structure} and rational descent now prove the assertion.
\end{proof}
\begin{corollary}[A single exceptional cycle in Type III]
\label{p:S7onecycle}
For a simple Type III(1) sixfold $A$, the Hodge conjecture holds on
all powers if and only if there exists $Z\in\CH^3(A)_\Q$ with
$\mathrm{cl}(Z)\notin\Q\theta^3$, where $\theta$ is a polarization.
\end{corollary}
\begin{proof}
Theorem~\ref{t:S7structure} gives
$\B^3(A)=\Q\theta^3\oplus\mathcal E(A)$ and
$\mathcal E(A)_\C=\Sym^6 U\otimes\det W$.
The exceptional summand is primitive: degree eight has only the
one-dimensional divisor contribution, and $D^\times$-equivariance
of multiplication by $\theta$ (with its reduced-norm character twist)
excludes a map from the seven-dimensional irreducible summand to it. Hence its projection is
algebraic on Hodge classes, explicitly
\[
 \mathrm{cl}(Z)\longmapsto\mathrm{cl}(Z)
 -\frac{\int_A\mathrm{cl}(Z)\theta^3}{\int_A\theta^6}\,\theta^3.
\]
The action of $D^\times$ on $\mathcal E(A)$ is absolutely irreducible:
after complexification its multiplicity action is the irreducible
$\Sym^6$ representation of $\GL_2$. Rational quaternion elements
are Zariski dense in $M_2(\C)$, and their pullbacks are algebraic
correspondences after clearing denominators. Thus the orbit of any
nonzero algebraic exceptional projection spans all of $\mathcal E(A)$.
Its complex span contains the pure volume vectors $u^6$.
The $D$-linear sum maps of Theorem~\ref{t:S7split} then span
$\Sym^6(U^{\oplus n})$ on every power. All pairings are divisors,
proving sufficiency; necessity follows from $\dim\mathcal E(A)=7$.
This criterion does not construct the required $Z$.
\end{proof}

\begin{example}\label{ex:hamilton}
Let $D=(-1,-1)_\Q$ with $i^2=j^2=-1$, and $V=D^3$. Put
\[
 \psi(x,y)=\Trd\left(\sum_{k=1}^3\bar x_k y_k i\right).
\]
On an order lattice, clear denominators. Right multiplication by $-i$ is
a compatible complex structure, since
$\psi(x,x(-i))=\Trd(\sum_k\bar x_kx_k)>0$ for $x\ne0$.
The left $D$-action has Rosati adjoint $a\mapsto\bar a$.
The compatible PEL domain is $\SO^*(6)/\mathrm U(3)$, of complex dimension
$3$. Its group is noncompact and $\Q$-simple; arithmetic monodromy is dense
in its connected $\SO_6$ complex form. The generic commutant is $D$;
very general members are therefore simple Type III(1) sixfolds.
Take $K=\Q(i)$ acting on the left. The associated Hermitian form can be
written as $h(x,y)=\psi(ix,y)+i\psi(x,y)$. On the $K$-basis $(1,j)$ of
each copy of $D$, it has matrix $\mathrm{diag}(2,-2)$.
Thus on $D^3$ it is three hyperbolic planes; the span of the three vectors
$1+j$ in separate summands is a three-dimensional totally isotropic
$K$-subspace. Its discriminant is $-64$, the class of $-1$ modulo norms.
This computation uses the fixed rational polarization, so it remains valid
throughout the compatible family. The signature is $(3,3)$: left $j$
interchanges the two $K$ eigenspaces and commutes with the complex
structure. Theorem~\ref{t:S7split} applies. The eight middle Hodge classes
consist of $\theta^3$ and seven volume classes; the latter are spanned by
$D$-translates of the algebraic Weil plane.
\end{example}
\begin{remark}
The following criterion determines exactly when this split-subfield
method applies. The subsequent example shows that it does not apply
to every Type III sixfold.
\end{remark}
\begin{proposition}[An arithmetic split-subfield criterion]
\label{p:quaternion-split-test}
Let $A$ be a simple Type III(1) sixfold with definite quaternion
endomorphism algebra $D$. Diagonalize its rational polarization on
$H_1(A,\Q)=D^3$ as
\[
 \psi(x,y)=\Trd\!\left(\sum_{s=1}^3\bar x_s y_s q_s\right),
 \qquad \bar q_s=-q_s,\quad q_s\ne0,
 \qquad N=\prod_{s=1}^3\operatorname{Nrd}(q_s)>0.
\]
There exists an imaginary quadratic subfield $k\subset D$ for which
this polarization has $k$-Weil discriminant $-1$ if and only if
$\Q(\sqrt{-N})$ embeds in $D$. When it does, the Hodge conjecture
holds on every power of $A$.
\end{proposition}
\begin{proof}
The displayed diagonalization is the usual orthogonalization of a
nondegenerate quaternionic skew-Hermitian form. Fix any imaginary
quadratic $k=\Q(i)\subset D$, and choose $j$ with
\[
 i^2=-a,\qquad j^2=-b,\qquad ij=-ji,\qquad a,b\in\Q_{>0}.
\]
The canonical Rosati involution makes the polarization compatible with
$k$. Left multiplication by $j$ interchanges its two eigenspaces and
commutes with the complex structure, so its signature is $(3,3)$.
For one diagonal entry $q=xi+yj+zij$, the associated $k$-Hermitian
form $h(v,w)=\psi(iv,w)+i\psi(v,w)$, on the basis $(1,j)$, is
\[
 \begin{pmatrix}
 2ax&2abz-2byi\\
 2abz+2byi&-2abx
 \end{pmatrix}.
\]
Its determinant is $-4ab\operatorname{Nrd}(q)$, since
$\operatorname{Nrd}(q)=ax^2+by^2+abz^2$.
Taking the three blocks, and using that $a$ and every rational square
are norms from $k$, gives
\[
 [\det_k h]=[-bN]\quad\hbox{in }\Q^\times/N_{k/\Q}(k^\times).
\]
Consequently the split condition is $bN\in N_{k/\Q}(k^\times)$.
The quaternion algebra $D=(k,-b)$ then also has the presentation
$(k,-N)$: the ratio $N/b$ is a norm, so rescale its second generator.
This produces an element of square $-N$, proving necessity.

Conversely, if $u\in D$ has square $-N$, choose a nonzero pure
quaternion $i$ anticommuting with $u$. Definiteness gives $i^2=-a<0$,
and $D=(\Q(i),-N)$. In the preceding determinant calculation take
$b=N$; then $bN=N^2$ is a norm. The discriminant is $-1$.
Theorem~\ref{t:S7split} proves the all-powers assertion.
\end{proof}
\begin{example}[A quaternionic family outside every split subfield]
\label{ex:quaternion-nonsplit}
Take $D=(-2,-5)_\Q$, with generators $i,j$, and
$(q_1,q_2,q_3)=(i,j,ij)$. Their reduced norms are $2,5,10$, hence
$N=100$ and $\Q(\sqrt{-N})=\Q(\sqrt{-1})$.
The algebra $D$ is division at $5$: the local Hilbert symbol is
$(-2,-5)_5=(\frac{-2}{5})=-1$. But $\Q(\sqrt{-1})$ splits at $5$,
so it cannot embed in $D$. Indeed tensoring such an embedding with
$\Q_5$ would embed $\Q_5\times\Q_5$ in a division algebra.
Thus Proposition~\ref{p:quaternion-split-test} excludes every split
imaginary quadratic subfield for this polarization.

This is a nonempty polarized datum. On the $s$th real quaternionic
coordinate take right multiplication by
$-q_s/\sqrt{\operatorname{Nrd}(q_s)}$ as complex structure. Its square
is $-1$ and
\[
 \psi(x,Jx)=2\sum_s\sqrt{\operatorname{Nrd}(q_s)}\,
                  \operatorname{Nrd}(x_s)>0\quad(x\ne0).
\]
The Riemann relations and an order lattice give abelian varieties.
The same arithmetic-monodromy and commutant argument as in
Example~\ref{ex:hamilton} gives simple very general members with
$\End^0=D$ in a three-dimensional PEL family.
For such a member every rational polarization is a positive rational
multiple of this one, because the Rosati-fixed subspace of $D$ is
$\Q$. A multiplier $c$ replaces $N$ by $c^6N$, which defines the
same quadratic field. Hence changing polarization cannot remove this
obstruction. This is an obstruction to the split-subfield method,
not a nonalgebraicity assertion. Nevertheless,
Theorem~\ref{t:S7complete} now proves the Hodge conjecture on all
powers of every simple member, using the eightfold input. For example,
choose the subfield $k=\Q(\sqrt{-2})\subset D$ generated by $i$ and stabilize its
sixfold Weil structure as in Theorem~\ref{p:all-sixfold-weil};
no split imaginary quadratic subfield inside $D$ is required.
\end{example}
\section{S8: Type IV(1,1) (\provedstatus)}\label{sec:S8}

\status{The whole row is proved on all powers. Unbalanced signatures
are divisor-generated. For balanced signature, arbitrary rational
Hermitian discriminants are covered by Theorem~\ref{p:all-sixfold-weil},
using the new external split-eightfold input.}
\begin{theorem}\label{t:S8}
Suppose $A$ is a simple sixfold with $\End^0(A)=K$ imaginary quadratic.
For the $K$-action on $H^{1,0}$ write its signature as $(r,6-r)$, with
$1\leq r\leq5$. If $r\ne3$, then $\Hg(A)$ is the full unitary group
and every Hodge class on every power is divisor-generated.
If $r=3$, then $\Hg(A)$ is the special unitary group,
\[
 P_A(t)=1+t+t^2+3t^3+t^4+t^5+t^6,
\]
and the Hodge conjecture on all powers is equivalent to algebraicity of
$W_K(A)$. Under the split-eightfold input it holds for every
discriminant, on all powers.
\end{theorem}
\begin{proof}
Over $\C$ write $V=W\oplus W^*$ with $\dim W=6$. The two summands are
irreducible and nonisomorphic for the full Hodge group, since its commutant
is $K\otimes\C=\C^2$. Its center acts by scalars on each, so the derived
representation on $W$ is irreducible.
The classical minuscule list in dimension $6$ leaves: the standard
representation of $\SL_6$, the standard symplectic representation of
$\Sp_6$, the orthogonal $\bigwedge^2\C^4$ of $\SL_4$, or
$\C^2\otimes\C^3$ for $\SL_2\times\SL_3$.
The last case has two different absolute root types and hence two rational
normal factors. A nontrivial Hodge cocharacter on both would give at least
three weights on $W$; one inactive rational factor contradicts minimality
of the Hodge group, as in S1. It is excluded.
In the symplectic or orthogonal case a nontrivial cocharacter with only
two weights has opposite weights of equal multiplicity $3$. Adding a
central weight does not change those multiplicities, so the signature must
be $(3,3)$. For that signature the determinant of $K$-linear $H^1$ has zero
Hodge-circle weight. As $K$ is imaginary quadratic, all its rational
conjugates also have weight zero; hence $\Hg(A)\subset\SU_K(h)$.
Its connected center is then finite: a scalar on $W$ with determinant one
has sixth power one. Thus $W$ and $W^*$ would be isomorphic for $\Hg(A)$
in either of the self-dual alternatives, enlarging the commutant. This
excludes both. The derived image is $\SL_6$.
If $r=3$, its full group is $\SU_K(h)$. If $r\ne3$, the determinant weight
$2r-6$ is nonzero; it generates the one-dimensional unitary determinant
torus. The full group is therefore $\mathrm U_K(h)$.
General-linear invariant theory proves divisor generation in the latter
case. In the former, contractions give polarization powers and the two
determinants give the rational Weil plane, proving the polynomial.
Necessity of Weil algebraicity is immediate on $A$, and sufficiency on all
powers is Lemma~\ref{p:polarization}. Theorem~\ref{p:all-sixfold-weil}
supplies its algebraicity for every discriminant under the new input.
For discriminant $-1$, Markman's theorem already suffices.
The excluded signatures $(0,6)$ and $(6,0)$ would make the Hodge circle
$K$-scalar on $W$; its rational containing group is then a torus with a
larger commutant, contradicting simplicity with endomorphism algebra $K$.
\end{proof}
\begin{example}
Take $K=\Q(i)$, $V=K^6$, and
$h=\mathrm{diag}(1,1,1,-1,-1,-1)$ with alternating form
$\Tr_{K/\Q}(i h)$ on an order lattice. Compatible positive complex
structures form the nonempty unitary domain of signature $(3,3)$, of
complex dimension $9$. The generic monodromy is the full special unitary
group, whose commutant is $K$; very general members are simple with
$\End^0=K$. The Hermitian form is hyperbolic and has discriminant $-1$.
The middle Hodge space has dimension $3$, generated by $\theta^3$ and an
algebraic rational basis of $W_K(A)$. Their pullbacks under $K$-linear sum
maps, together with divisors, generate the Hodge classes on every power.
For a comparison of boundaries, replacing the last diagonal entry by
$-3$ gives determinant $-3$. The positive number $3$ is not a norm from
$\Q(i)$ (its $3$-adic valuation is odd), so the split theorem does not
apply directly to that family. Its product with two oppositely oriented
$\Q(i)$-CM elliptic curves with Hermitian coefficients $1,-3$
has signature $(4,4)$ and determinant $9$, hence is split.
Theorem~\ref{p:all-sixfold-weil} contracts its algebraic Weil plane
to the original sixfold. Lemma~\ref{p:polarization} then proves
algebraicity on every power, including this formerly unresolved family.
\end{example}
\section{S9: Type IV(2,1) (\provedstatus)}\label{sec:S9}

\status{The entire row is proved, including the exceptional Weil case.
The real-quadratic polarization adjustment makes its Weil plane split;
the determinant argument then also proves algebraicity on all powers.}
\begin{theorem}
Let $A$ be simple with CM endomorphism field $K$ of degree $4$.
Its derived Hodge group has complex form $\SL_{3}^{2}$.
There are exactly the following alternatives for its determinant image.
\begin{enumerate}
\item The image is all of $U_K$. Then $\Hg(A)_\C=\GL_{3}^{2}$,
all powers are divisor-generated, and $P_A(t)=(1+t+t^2+t^3)^2$.
\item The image is $\ker N_{K/k}$ for an imaginary quadratic subfield
$k\subset K$. Then $A$ is of $k$-Weil type,
\[
 B^\bullet(A)=D^\bullet(A)+W_k(A),\qquad
 P_A(t)=(1+t+t^2+t^3)^2+2t^3.
\]
The sum adds a two-dimensional space only in codimension three.
Every class is algebraic, and the Hodge conjecture holds on all powers.
\end{enumerate}
\end{theorem}
\begin{proof}
Lemma~\ref{p:smallunitary} gives the derived group and determinant torus.
The full Hodge group is the preimage of that torus under the determinant
map on the full unitary Lefschetz group: it contains the entire derived
kernel, and the two connected groups have the same dimension.
For the full image, the general-linear invariant theorem gives pairings.
For the norm kernel, orient the $K$ embeddings over one embedding of $k$.
The split determinant relation is $z_1\cdots z_{2}=1$.
The invariant computation in Proposition~\ref{p:coupled-det} gives pairings
and the two simultaneous volumes, with the indicated dimensions. That
computation does not require algebraicity. For the geometric assertion,
write $F=K^+$ and $K=kF$. A polarization has Rosati conjugation on $K$,
so Lemma~\ref{p:quartic-scale} supplies a polarization with $k$-discriminant
$-1$ and makes the Weil plane algebraic. Proposition~\ref{p:coupled-det}
then proves algebraicity on all powers.
\end{proof}
\begin{example}
Let $K=\Q(\zeta_5)$, $F=\Q(\sqrt5)$, $V=K^3$, and
$h=\mathrm{diag}(1,1,-1)$. Choose a nonzero totally imaginary
$\delta\in K$ and the trace alternating form $\Tr_{K/\Q}(\delta h)$,
with compatible positive complex structures and an order lattice.
The signatures at the two real embeddings of $F$ are $(2,1)$.
This gives a nonempty PEL domain of complex dimension $4$.
Arithmetic monodromy gives generic derived group
$\Res_{F/\Q}\SU(h)$, whose complex form is $\SL_3^2$.
The determinant Hodge cocharacter is nonzero. On the two-dimensional
anti-conjugate character space, a generator of the cyclic Galois group
$\mathrm{Gal}(K/\Q)$ acts by a signed rotation of order $4$; this
representation is irreducible over $\Q$. The orbit of any nonzero
determinant vector therefore spans that entire space. The generic
central torus is full, so the generic Hodge group is the full unitary group.
Its centralizer is $K$; very general members are simple with $\End^0=K$.
Their Hodge dimensions are $(1,2,3,4,3,2,1)$, generated by the two
rational Rosati-fixed divisor classes, with the usual Poincare generators
on powers. This verifies the group hypothesis rather than assuming it.
\end{example}

\begin{proposition}\label{p:coupled-det}
Suppose a sixfold has a CM-field action by $K$ of degree
$2e$, $q=6/e$, $e\in\{2,3,6\}$, and
\[
 \Hg(A)_\C^{\dd}=\prod_{j=1}^e\SL(W_j),\quad \dim W_j=q.
\]
Choose one embedding in each conjugate pair. Suppose the determinant
image of $\Hg(A)_\C$ is the connected torus
$\{(z_1,\ldots,z_e):\prod z_j=1\}$, and these selected embeddings are
exactly those restricting to a fixed embedding of an imaginary quadratic
subfield $k\subset K$. If $W_k(A)$ consists of algebraic classes,
the Hodge conjecture holds for $A^n$ for every $n$. In particular
this holds when the induced $k$-Weil structure has discriminant $-1$.
On $A$,
\[
 P_A(t)=(1+t+\cdots+t^q)^e+2t^3.
\]
\end{proposition}
\begin{proof}
The exterior invariant algebra for the derived group is generated by
pairings between $W_j$ and $W_j^*$ and determinant tensors on $q$ copies
of either space. This is the first fundamental theorem for $\SL_q$,
applied to the tensor algebra and then its exterior quotient. For $q=1$
the derived group is trivial and the determinant generators are the
individual degree-one coordinates, so the assertion is immediate. A determinant
and a dual determinant in the same factor have product equal, up to the
alternating convention, to a determinant of pairings. After eliminating
such pairs, a central invariant determinant monomial must have the same
signed determinant exponent in every factor: the kernel of the characters
on the displayed torus is generated by $(1,\ldots,1)$. Consequently the
invariant algebra is generated by pairings and products of one determinant
from each factor, all positive or all negative. Those products have
cohomological degree $eq=6$.

On one copy the two products constitute precisely the complexification
of the rational $k$-Weil plane $\bigwedge_k^6 H^1(A,\Q)$, embedded in
$H^6(A,\Q)$. They are algebraic by hypothesis; discriminant $-1$
is a sufficient way to verify this hypothesis by
\cite[Theorem 1.5.1]{Markman}. To treat powers, pull back this whole plane
by the algebraic sum maps $A^n\to A$ with coefficients in $K$. Clearing
denominators realizes all coefficients as endomorphisms. At the selected
$e$ embeddings the coefficients vary Zariski densely and independently
in $(\C^n)^e$, since $K\otimes_\Q\C$ is the product of its embedding
coordinates. Thus these pullbacks span
$\bigotimes_{j=1}^e\Sym^q(\C^n)$, the space of all positive determinant
products; pure $q$th powers span each symmetric power. The conjugate plane
gives all negative products. The pairings are divisor classes, by
Lefschetz $(1,1)$. Surjectivity descends from $\C$ to $\Q$ by scalar
extension of the rational cycle-class spans, proving the assertion.
On one copy pairings contribute the displayed product polynomial. Only
the two simultaneous volumes survive in addition, both in codimension
three; this proves the dimension formula.
\end{proof}
\begin{example}[A split coupled quartic-CM family]
Take $K=\Q(i,\sqrt2)$, $F=\Q(\sqrt2)$, $k=\Q(\sqrt{-2})$,
$V=K^3$, and $h_K=\operatorname{diag}(1,1,-1)$. Use the alternating form
$\psi=\Tr_{K/\Q}(i h_K)$ on an order lattice. The signatures relative
to $\Q(i)$ are $(2,1),(2,1)$; the compatible PEL domain has dimension
four. Its generic derived group is $\Res_{F/\Q}\SU(h_K)$, with complex
form $\SL_3^2$, by arithmetic monodromy. The determinant cocharacter
relative to these two embeddings is proportional to $(1,1)$. The
biquadratic Galois group interchanges its coordinates and negates them
simultaneously, so its orbit spans exactly this line. This determines the
generic central determinant torus.

Relative instead to the two embeddings restricting to one embedding of
$k$, one $\Q(i)$-embedding is conjugated. Thus the signatures are
$(2,1),(1,2)$, the central relation is $z_1z_2=1$, and the total $k$
signature is $(3,3)$. Writing $\delta'=i\sqrt2$, the $k$-Hermitian form
associated with the same $\psi$ is
\[
 h_k=\Tr_{K/k}\bigl((1/\sqrt2)h_K\bigr).
\]
On the $k$-basis $(1,\sqrt2)$ its first two blocks are
$\left(\begin{smallmatrix}0&2\\2&0\end{smallmatrix}\right)$, and the
third is the negative of this block. The span of the first vector in each
block is a three-dimensional isotropic $k$-subspace. This proves splitness
explicitly. The generic centralizer on $V$ is $K$: the two derived
standard modules are inequivalent and each is distinguished from its dual
by the nontrivial center. Hence a very general member is simple with
$\End^0(A)=K$. Proposition~\ref{p:coupled-det} applies: every power
satisfies the Hodge conjecture, while
$P_A(t)=1+2t+3t^2+6t^3+3t^4+2t^5+t^6$.
In particular its two exceptional middle classes are genuinely present
and algebraic; this example is not divisor-generated.
\end{example}

\begin{example}[A nonsplit polarization that can be corrected]
Let $k=\Q(i)$, $F=\Q(\sqrt3)$, $K=kF$, and
$H=\operatorname{diag}(1,-1,1+\sqrt3)$ on $K^3$.
Use $\psi=\Tr_{K/\Q}(iH)$ and a compatible very general complex structure.
The signatures are $(2,1)$ and $(1,2)$. The product derived monodromy
and the determinant-circle relation give the group in the second
alternative of S9; its commutant is $K$ and the sixfold is simple.
Here $b=-1-\sqrt3$, $N_{F/\Q}(b)=-2$, and the initial $k$-determinant is
\[
 D=12^3(-2)=-3456=-6\cdot24^2.
\]
It is not the class of $-1$: the factor $3$ occurs to odd order and
is inert in $\Q(i)$. But $a=3+\sqrt3$ is totally positive and has
norm $6$. Replacing $H$ by $aH$ gives
\[
 D'=D\cdot6^3=-746496=-864^2.
\]
This is a polarization of discriminant $-1$ on the same complex abelian
variety. Its Weil plane is unchanged and is therefore algebraic.
The Hodge dimensions are $(1,2,3,6,3,2,1)$.
This example tests the nontrivial step of Lemma~\ref{p:quartic-scale}:
scaling by a totally positive real-quadratic endomorphism can change
the norm class even though a rational scalar on a sixfold cannot.
\end{example}

\section{S10: Type IV(3,1) (\provedstatus)}\label{sec:S10}

\status{The whole row is proved on all powers using the new external
split-eightfold input for its exceptional Weil plane. The full-unitary
case is divisor-generated. The explicit split and nonsplit cyclic
constructions below remain independent geometric examples of the row.}
\begin{theorem}
Let $A$ be simple with CM endomorphism field $K$ of degree $6$.
Its derived Hodge group has complex form $\SL_{2}^{3}$.
There are exactly the following alternatives for its determinant image.
\begin{enumerate}
\item The image is all of $U_K$. Then $\Hg(A)_\C=\GL_{2}^{3}$,
all powers are divisor-generated, and $P_A(t)=(1+t+t^2)^3$.
\item The image is $\ker N_{K/k}$ for an imaginary quadratic subfield
$k\subset K$. Then $A$ is of $k$-Weil type,
\[
 B^\bullet(A)=D^\bullet(A)+W_k(A),\qquad
 P_A(t)=(1+t+t^2)^3+2t^3.
\]
The sum adds a two-dimensional space only in codimension three.
Under the split-eightfold input, every class on every power is
algebraic, for arbitrary induced $k$-Hermitian discriminant.
\end{enumerate}
\end{theorem}
\begin{proof}
Lemma~\ref{p:smallunitary} gives the derived group and determinant torus.
The full Hodge group is the preimage of that torus under the determinant
map on the full unitary Lefschetz group: it contains the entire derived
kernel, and the two connected groups have the same dimension.
For the full image, the general-linear invariant theorem gives pairings.
For the norm kernel, orient the $K$ embeddings over one embedding of $k$.
The split determinant relation is $z_1\cdots z_{3}=1$.
The invariant computation in Proposition~\ref{p:coupled-det} gives pairings
and the two simultaneous volumes, with the indicated dimensions. That
computation does not require algebraicity. Theorem~\ref{p:all-sixfold-weil}
now makes the entire $k$-Weil plane algebraic for arbitrary discriminant.
Proposition~\ref{p:coupled-det} then proves algebraicity on all powers.
For discriminant $-1$, the earlier Theorem~\ref{p:markman} suffices.
\end{proof}
\begin{example}
Set $K=\Q(\zeta_7)$, $F=\Q(\alpha)$ with
$\alpha=2\cos(2\pi/7)$, and $V=K^2$. Choose
$h=\mathrm{diag}(1,\alpha-1)$ and an integral multiple of
$\Tr_{K/\Q}(\delta h)$ for a nonzero totally imaginary $\delta$.
The conjugates of $\alpha$ are approximately
$1.24698,-0.44504,-1.80194$. Thus the signatures are $(2,0),(1,1),(1,1)$,
and the nonempty compatible period domain has complex dimension $2$.
The rational derived unitary group is simple and has noncompact real
factors. Arithmetic monodromy is dense in it, including its compact real
factor, so its generic complex form is $\SL_2^3$.
The determinant cocharacter has just one nonzero component in the
three-dimensional anti-conjugate character space. The cyclic Galois group
of $K$ permutes its three coordinate directions with signs; that vector's
orbit spans all three directions. The central determinant torus is full,
so the generic Hodge group is $\GL_2^3$ after complexification.
The commutant is $K$, proving simplicity and $\End^0=K$ for very general
members. Their Hodge dimensions are $(1,3,6,7,6,3,1)$; a basis of $F$
gives three rational divisor generators. Although the label is Type IV,
these members are not CM points: their Hodge group has nontrivial derived
part. Calling the entire row toric would be incorrect.
\end{example}

\begin{example}[A split coupled sextic-CM family]
Let $K=\Q(\zeta_7)=kF$, where $k=\Q(\sqrt{-7})$ and
$F=\Q(\alpha)$, $\alpha=2\cos(2\pi/7)$,
$\alpha^3+\alpha^2-2\alpha-1=0$. Set
\[
 a=\alpha^2+\alpha-1,\qquad b=\alpha-1,\qquad
 V=K^2,\qquad h_K=\operatorname{diag}(a,b),
\]
and use $\Tr_{K/\Q}(\sqrt{-7}\,h_K)$ on an order lattice.
At the three real embeddings the signs of $a$ are $(+,-,+)$ and those
of $b$ are $(+,-,-)$, so the signatures are $(2,0),(0,2),(1,1)$.
The compatible domain has dimension one and generic derived group
$\Res_{F/\Q}\SU(h_K)$, of complex form $\SL_2^3$.
The determinant cocharacter is proportional to $(1,-1,0)$.
The cyclic Galois action on $F$ generates the plane of coordinate sum
zero; conjugation negates it. Thus the determinant image is exactly
$z_1z_2z_3=1$. The three standard derived modules are distinct and their
central characters are nonzero, so none is identified with its dual;
the generic centralizer is $K$. Very general members are simple with
$\End^0(A)=K$.

The induced $k$-Hermitian form is $h_k=\Tr_{K/k}(h_K)$ and has signature
$(3,3)$. In the basis $(1,\alpha,\alpha^2)$ the trace-form determinant
of $F$ is $49$. The multiplication matrices are
\[ M_a=\begin{pmatrix}-1&1&0\\1&1&1\\1&0&1\end{pmatrix},\qquad
M_b=\begin{pmatrix}-1&0&1\\1&-1&2\\0&1&-2\end{pmatrix}. \]
Their determinants are $-1$ and $1$, respectively. Therefore
\[
 \det(h_k)=49^2\operatorname{Norm}_{F/\Q}(ab)=-49^2.
\]
In $\Q^*/\operatorname{Norm}_{k/\Q}(k^*)$ this is $-1$, since
$49^2$ is the norm of the rational number $49$.
This is exactly the discriminant hypothesis of
\cite[Theorem 1.5.1]{Markman}; no additional identification of Hermitian
forms is needed. Proposition~\ref{p:coupled-det} gives the Hodge
conjecture on all powers and
\[ P_A(t)=1+3t+6t^2+9t^3+6t^4+3t^5+t^6. \]
The two additional degree-six classes here are the algebraic $k$-Weil
plane. The full-unitary example above has middle dimension seven, so
endomorphism degree alone demonstrably does not determine Hodge dimensions.
\end{example}

\begin{lemma}[All powers of the Fermat-eighteen Jacobian]\label{p:fermat18}
Let $F_{18}$ be the Fermat curve of degree eighteen. The Hodge
conjecture holds for every power of $J(F_{18})$.
\end{lemma}
\begin{proof}
The $(\mu_{18})^2$ characters on $H^1(F_{18})$ have multiplicity
one and labels $(a,b)$ with $a,b,a+b\ne0$ modulo eighteen.
Their rational Galois orbits have size two or six. The corresponding
abelian factors have CM by subfields of
$K=\Q(\zeta_{18})=\Q(\zeta_9)$ and dimension one or three.
The cyclic sextic field $K$ has a unique proper CM subfield
$k=\Q(\sqrt{-3})$. Of its eight CM types, two are induced from
$k$, while the other six form one orbit under $\operatorname{Gal}(K/\Q)$.
Thus there is only one primitive threefold isogeny class $X$;
an imprimitive threefold factor is isogenous to $E^3$, where
$E$ has CM by $k$. Consequently $J(F_{18})\sim X^a\times E^b$.
Corollary 1.11 of \cite{MilneTateStandard} proves the Hodge
conjecture for every $X^r\times E^s$, and hence the assertion.
The exact character enumeration gives $a=36,b=28$, although
these multiplicities are not needed for the argument.
\end{proof}

\begin{theorem}[A nonsplit cyclic-eighteen family]\label{p:S10cyclic18}
For $t\in\C\setminus\{0,1\}$ let $C_t$ be the smooth projective
model of
\[
 y^{18}=x(x-1)^4(x-t)^4,
\]
and let $B_t$ be its primitive $K=\Q(\zeta_{18})$ Jacobian factor.
Then $\dim B_t=6$, its $k=\Q(\sqrt{-3})$ action has signature
$(3,3)$, and its Weil plane $W_k(B_t)$ is algebraic for every $t$.
The Hodge conjecture holds for every power of every $B_t$.
For very general $t$, $B_t$ is simple with $\End^0(B_t)=K$, and
\[
 \Hg(B_t)=\det_K^{-1}(\ker(N_{K/k}:U_K\to U_k)),\qquad
 P_{B_t}(z)=(1+z+z^2)^3+2z^3.
\]
The inherited polarization has $k$-Hermitian determinant class
$-2$, not $-1$. For these very general members no change of
polarization makes this Weil structure split.
\end{theorem}
\begin{proof}
Put $\zeta=\zeta_{18}$ and orient the three $K$ embeddings over
one embedding of $k$ by $j=1,7,13$. The branch vector is
$(1,4,4,9)$; its entries are all nonzero modulo eighteen, so
every primitive eigenspace has dimension two. The formula
in \cite[Lemma 2.7]{MoonenCyclic} gives their $H^{1,0}$
multiplicities $(2,1,0)$. Thus the three signatures are
$(2,0),(1,1),(0,2)$ and their $k$-sum is $(3,3)$.
The signed determinant cocharacter is $(2,0,-2)$. Its cyclic
Galois orbit spans the coordinate-sum-zero plane; hence its
torus is exactly $\ker N_{K/k}$ for every $t$. In particular
\[
 T=\bigwedge\nolimits_K^2H^1(B_t,\Q)\subset H^2(B_t,\Q)
\]
has dimension six and no Tate subquotient. The displayed
inclusion again means the sum of the same-embedding determinant
lines, as in Lemma~\ref{p:cyclic-transfer}.

\emph{A fixed quotient with algebraic tensor classes.}
Write $B=B_t$, $C=C_t$, and let $\sigma$ be the covering automorphism.
The dihedral quotient of $C\times C$ by
$(\sigma,\sigma^{-1})$ and interchange of the factors has field
\[
 a=x_1+x_2,\quad b=x_1x_2,\quad w=y_1y_2,\qquad
 w^{18}=b(b-a+1)^4(b-ta+t^2)^4.
\]
The right side has valuation one along $b=0$, so this field has
degree eighteen over $\C(a,b)$ and is the full invariant field.
After homogenizing, its branch divisor consists of the four
lines $l_1,l_2,l_3,l_4$ in Lemma~\ref{p:cyclic-transfer}, now with
exponents $(1,4,4,9)$. The same linear change of plane coordinates
and a rescaling of $w$ identify a resolution $R_t$ birationally with
\[
 R:\quad W^{18}=XY^4Z^9(X+Y+Z)^4.
\]
The Fermat surface $A^{18}+B^{18}+C^{18}+D^{18}=0$ dominates
this normalization by
\[
 [A:B:C:D]\longmapsto[A^{18}:B^{18}:C^{18}:AB^4C^9D^4].
\]
As in Lemma~\ref{p:cyclic-transfer}, that Fermat surface is
dominated by $F_{18}\times F_{18}$. Both maps are generically
finite. Lemma~\ref{p:fermat18}, Abel--Jacobi pullback, the point
blowup formula, and generically finite push--pull therefore prove
the Hodge conjecture on every $R_t^n$. This uses the all-powers
statement of that lemma, not only algebraicity on a single surface.

The sum Abel--Jacobi map $C^2\to B$ pulls $T$ back by
antisymmetrization. This is fixed by the dihedral group, since
interchange acts by minus the tensor flip on $H^1(C)^{\otimes2}$.
Transfer gives an algebraic injection $\alpha:T\to H^2(R_t)$.
The algebraic determinant projector and orthogonality argument
in Lemma~\ref{p:cyclic-transfer} apply with three conjugate pairs
instead of two. For an associated $K$-compatible polarization,
the adjoint is now
\[
 \alpha^\dagger=p_T L_B^{-4}\,{}^t\!\alpha:
 H^2(R_t)\longrightarrow T,
 \qquad L_B^{-4}:H^{10}(B)\longrightarrow H^2(B).
\]
This inverse Lefschetz operator is algebraic. Both copies of $T$
are primitive because they have no Tate quotient. Positivity
and Cayley--Hamilton again give an algebraic left inverse
$\beta=(\alpha^\dagger\alpha|_T)^{-1}\alpha^\dagger$.
It follows from the Hodge conjecture for $R_t^3$ that every Hodge
class in $T^{\otimes3}\subset H^6(B^3)$ is algebraic.
The cup-product map from $T^{\otimes3}$ to $H^6(B)$ contains
$W_k(B)$ in its image: over $\C$ each $k$-volume is the product
of the three $K$-determinant lines above that embedding of $k$.
Semisimplicity of polarizable rational Hodge structures lifts
each Weil class to a Hodge class in $T^{\otimes3}$. Applying
the preceding algebraic correspondences and then the small
diagonal pullback proves algebraicity of $W_k(B)$.

\emph{The group and the special parameters.}
The balanced eigenspace has connected monodromy $\SU(1,1)$ by
\cite[Proposition 4.7(i)]{MoonenCyclic}, so the generic fiber is
not CM. More generally, at any non-CM parameter the derived
projections to the three two-dimensional blocks must all be
$\SL_2$, by Galois transitivity. A Galois-stable diagonal
identification would identify all three factors, which is
impossible since the adjoint Hodge circle is trivial on the pure
blocks and nontrivial on the balanced block. Thus the derived
group is $\SL_2^3$ and the determinant computation gives the
stated full preimage. Its commutant is $K$: the three standard
blocks are distinct and their nontrivial determinant characters
separate each block from its dual. This proves simplicity,
the endomorphism assertion, and the displayed polynomial.
Proposition~\ref{p:coupled-det} now gives algebraicity on all
powers from the Weil plane just constructed.

At a CM parameter, if $B$ is simple use Theorem~\ref{p:CMcomplete}.
Otherwise the isotypic projectors commute with $K$, so each
isotypic summand has dimension divisible by three. A single
isotypic summand is a power of a simple variety of dimension at
most three, whose powers are already covered by the low-dimensional
inputs. If there are two summands, each has dimension three and
is either a simple threefold or a cube of an elliptic curve.
Its $H^1$ has $K$-rank one, so its endomorphisms commuting with
$K$ lie in $K$. A simple threefold summand therefore has CM
field $K$ and belongs to the unique primitive class $X$ above;
an elliptic summand has field $k$. The only nonisotypic possibility
is consequently $X\times E^3$, whose powers are covered by
\cite[Corollary 1.11]{MilneTateStandard}. This proves the assertion
at every allowed parameter without using a cycle-deformation principle.

\emph{The nonsplit determinant.}
Put $F=K^+$ and $\delta=1+2\zeta^6=\sqrt{-3}$.
The geometric twisted-intersection matrix
\cite[Proposition 7.4]{YuZheng} gives, up to a scalar in $F$
and a $K$-basis change, the following skew-Hermitian form on
the primitive rank-two space, with $x=\zeta$, $y=z=\zeta^4$:
\[
 G=\begin{pmatrix}
 \dfrac{1-xy}{(1-x)(1-y)}&\dfrac1{y-1}\\[6pt]
 \dfrac{y}{y-1}&\dfrac{1-yz}{(1-y)(1-z)}
 \end{pmatrix},\qquad
 \det G=\frac{1-xyz}{(1-x)(1-y)(1-z)}.
\]
The matrix is the intersection pairing, not a form chosen merely
to have the desired signature. One can use the balanced embedding
$j=7$ in the cited construction and Galois-conjugate back to this
expression. Set $D=(1-\zeta)(1-\zeta^4)^2$. Then
\[
 D\in F,\qquad N_{F/\Q}(D)=3,\qquad
 \det(G/\delta)=-\frac2{3D},\qquad
 N_{F/\Q}\det(G/\delta)=-\frac8{81}.
\]
For $\alpha=\zeta+\zeta^{-1}$ one has
$\alpha^3-3\alpha-1=0$ and $D=2\alpha^2-3\alpha-2$,
which also verifies the norm above. The trace Gram matrix in the basis
$(1,\alpha,\alpha^2)$ is
\[
 \begin{pmatrix}3&0&6\\0&6&3\\6&3&18\end{pmatrix},
 \qquad \det=81.
\]
The determinant of the $k$-Hermitian trace form is therefore
$81^2(-8/81)=-648=-2\cdot18^2$. A conformal scalar changes the
rank-two $K$-determinant by a square in $F$, and a basis change
by a $K/F$ norm; both become $k/\Q$ norms after taking the trace
determinant. Thus the polarization class really is $-2$.
The prime $2$ is inert in $k$, so $2$ is not a $k/\Q$ norm;
this class is not $-1$. When $\End^0(B)=K$, every change of
polarization multiplies the determinant by $N_{F/\Q}(a)^2$
for some totally positive $a\in F$. It cannot remove the obstruction.
\end{proof}

The script \texttt{verify\_cyclic18.py} checks the six eigenspace
multiplicities, both rational maps, the Fermat character factors,
and the intersection matrix and norms in $\Q[\zeta]/(\zeta^6-\zeta^3+1)$.
This explicit cyclic construction proves a specific nonsplit family.
The complete row for all sextic fields and rational Hermitian forms
instead follows from Theorem~\ref{p:all-sixfold-weil} and the
all-powers argument above, using the new external input.

\section{S11: Type IV(6,1) (\provedstatus)}\label{sec:S11}

\status{Proved for every simple CM sixfold, on all powers, without
a Galois hypothesis. Its Hodge rank is five or six. In rank five the
sole extra generator is an algebraic split Weil plane. The finite
character test and explicit Galois computations below provide additional
certificates, not restrictions on the theorem.}
\begin{proposition}\label{t:CM}
Let $A$ be a CM sixfold with field $K$ of degree $12$. If its Hodge group
is the full norm-one torus
$T_K=\ker(\Res_{K/\Q}\mathbf G_m\to\Res_{K^+/\Q}\mathbf G_m)$,
then all Hodge classes on all powers are divisor-generated and
$P_A(t)=(1+t)^6$.
For a specified CM type, $\Hg(A)$ is the smallest rational subtorus
containing its Hodge circle; its cocharacter space is the rational span of
the Galois translates of the signed CM-type vector in a normal closure.
\end{proposition}
\begin{proof}
Over $\C$ the representation has six independent conjugate character
pairs $\chi_i,-\chi_i$, each of multiplicity one. On powers a monomial
is invariant under the full norm-one torus exactly when, for every $i$,
the numbers of its positive and negative occurrences agree. It is then
an exterior product of opposite-character degree-two pairings.
By the $(1,1)$ theorem and rational descent these are divisor-generated
algebraic classes. On $A$ one either chooses both members of a pair or
neither, giving $\binom6p$ invariants in codimension $p$.
The cocharacter assertion follows directly from the definition of the
smallest rational torus: its cocharacter space must contain the Galois
orbit and conversely the orbit span defines a rational subtorus after
saturating the lattice.
\end{proof}

\begin{proposition}[An exact finite test for every specified CM type]
\label{p:CMtest}
Let $A$ be simple with $\End^0(A)=K$ of degree twelve, and fix its CM
type. Form a rational matrix $M$ whose rows are all Galois translates
of its signed CM-type vector in six conjugate-pair coordinates. For
$\delta\in\{-1,0,1\}^6$ write $w(\delta)$ for its support size.
Let $n_j$ count the vectors with $M\delta=0$ and $w(\delta)=j$.
Then
\[
 \dim B^2(A)=15+n_4,\qquad
 \dim B^3(A)=20+2n_4+n_6.
\]
The divisor-generated dimensions are $15$ and $20$ respectively.
Thus $n_4=n_6=0$ proves the Hodge conjecture for $A$ even without a
prior full-rank assertion about $M$.
\end{proposition}
\begin{proof}
The twelve complex eigenlines have characters $x_i,x_i^{-1}$.
For an exterior monomial, let $\delta_i$ be the number of selected
$x_i$ coordinates minus the number of selected $x_i^{-1}$ coordinates.
The monomial is fixed by $\Hg(A)$ exactly when $M\delta=0$, by
Proposition~\ref{t:CM}. A coordinate with $\delta_i=0$ contributes
either zero or two to exterior degree; a nonzero one contributes one.
Consequently
\[
 \dim B^p(A)=\sum_{\substack{\delta\in\{-1,0,1\}^6\\M\delta=0}}
 \binom{6-w(\delta)}{p-w(\delta)/2},
\]
with a binomial interpreted as zero unless its lower index is an integer
between zero and its upper index.
There are no weight-two nonzero relations: such a relation would identify
two nonconjugate eigencharacters or make a conjugate character equal to
another one, enlarging the commutant beyond $K\otimes\C$.
The six opposite-coordinate pairings therefore form a basis of
$B^1(A)_\C$. Their squarefree products span $D^p(A)_\C$ with dimension
$\binom6p$. A weight-four relation contributes once in degree four and
twice in degree six; a weight-six relation contributes once in degree six.
This gives the formulas. Scalar extension preserves invariant dimensions
and the image of the divisor multiplication maps, proving them over
$\Q$. The last assertion follows from Lemma~\ref{p:HLsix}.
\end{proof}
\begin{remark}
This is an exhaustive computation for a specified field and CM type,
using at most $3^6=729$ exponent vectors. A character-neutral exceptional
monomial is not by itself an algebraic cycle. Nonzero $n_4$ or $n_6$
still requires a geometric algebraicity argument for its rational
Galois-stable span. Theorem~\ref{p:CMcomplete} below in fact shows
that for a simple CM sixfold $n_4=0$ and $n_6\in\{0,2\}$;
the test remains a useful independent check on the supplied CM data.
\end{remark}

\begin{corollary}[A checkable degenerate CM stratum]\label{p:CMquartic}
In Proposition~\ref{p:CMtest}, suppose $n_4=0$ and every weight-six
relation, up to sign, selects precisely the six embeddings above one
embedding of an imaginary quadratic subfield $k\subset K$.
Then all Hodge classes on $A$ are algebraic; no subfield condition
on $K^+$ is needed.
\end{corollary}
\begin{proof}
The vanishing character says that the selected $k$-determinant is a
Hodge class, equivalently that $A$ has $k$-signature $(3,3)$.
Lemma~\ref{p:CMsplit} makes its Weil plane algebraic. Apply this
to every listed relation. Proposition~\ref{p:CMtest} shows that these
planes and divisors exhaust codimensions two and three. Their rational
Galois-stable spans descend as in Lemma~\ref{p:fft}, and
Lemma~\ref{p:HLsix} finishes the proof.
\end{proof}

\begin{corollary}[All CM sixfolds of Hodge rank at least five]
\label{p:CMrankfive}
If $A$ is a simple CM sixfold and $\dim\Hg(A)\ge5$, the Hodge
conjecture holds for every power of $A$.
\end{corollary}
\begin{proof}
Write $K=\End^0(A)$. If the dimension is six, $\Hg(A)=U_K$ and all
invariants on all powers are pairings, hence divisor-generated.
If the dimension is five, \cite[Lemma (3.7)]{MZ} gives an imaginary
quadratic $k\subset K$ with
$\Hg(A)=\ker(N_{K/k}:U_K\to U_k)$.
The determinant is trivial on the Hodge circle, so the $k$-signature
is $(3,3)$. Lemma~\ref{p:CMsplit} proves algebraicity of $W_k(A)$;
Proposition~\ref{p:coupled-det} with $e=6,q=1$ proves the assertion on all
powers. On $A$ itself one has $n_4=0$, $n_6=2$, and
$P_A(t)=(1+t)^6+2t^3$.
\end{proof}
\begin{lemma}[A rank bound without a Galois hypothesis]\label{p:CMsixrank}
Every simple CM sixfold has Hodge rank at least five.
\end{lemma}
\begin{proof}
Put $K=\End^0(A)$, and let $G$ be the Galois group of a normal closure
of $K$. It acts by signed permutations on $\Q^6$, transitively on
the twelve signed coordinate vectors; complex conjugation acts as $-1$.
Let $f\in\{1,-1\}^6$ be the signed CM type and
$W=\langle Gf\rangle_\Q$. Then $r=\dim W=\dim\Hg(A)$.
The twelve weights of $H^1(A)_\C$ are distinct: equality of two would
enlarge its Hodge commutant beyond $K\otimes\C$.
Choose a basis of $W$ from $Gf$. Each weight takes only the values
$1,-1$ on these $r$ vectors, so there are at most $2^r$ different
weights. Hence $12\le2^r$ and $r\ge4$.

Suppose $r=4$. Use the standard $G$-invariant positive definite dot
product on $\Q^6$ and set $R=W^\perp$, of dimension two. Write
$p_i=\operatorname{pr}_R(e_i)$. None is zero, since transitivity
would otherwise make them all zero. The distinct vectors in
\[
 \mathcal P=\{p_1,-p_1,\ldots,p_6,-p_6\}
\]
form one orbit of the finite rational orthogonal group $\Gamma$
induced by $G$ on $R$, and span $R$. Each occurs with the same
multiplicity in this twelve-element multiset. Thus its cardinality
$m$ is an even divisor of twelve, and $m\ne2$; consequently
$m\in\{4,6,12\}$.

The rotation subgroup of a finite subgroup of $O(R_\R)$ is cyclic.
Since $-1\in\Gamma$, its order $n$ is even. A rational finite-order
two-dimensional rotation has integral trace $2\cos(2\pi/n)$;
thus $n\in\{2,4,6\}$. The group has order $n$ or $2n$.
In particular, if $m$ is six or twelve, $\Gamma$ contains a rotation
$\rho$ of order three, with $1+\rho+\rho^2=0$.

If $m=4$, orient the six coordinates so that three projections equal
$u$ and three equal $v$, where $u,v$ are independent. The relation
\[
 \sum_{i=1}^6\varepsilon_i p_i=0,
 \qquad \varepsilon_i\in\{1,-1\},
 \tag{*}
\]
would require a sum of three signs to vanish. No such signed vector
lies in $W$, contradicting $f\in W$.

If $m=6$, after orienting and ordering coordinates the projections
are $u,u,\rho u,\rho u,\rho^2u,\rho^2u$.
The unique relation among the three distinct oriented vectors is
$u+\rho u+\rho^2u=0$. Thus (*) says
\[
 \varepsilon_1+\varepsilon_2
 =\varepsilon_3+\varepsilon_4
 =\varepsilon_5+\varepsilon_6.
\]
Either all three sums are zero, or all six signs are equal.
The first possibility lies in the three-dimensional subspace $V$
whose sums on the three pairs vanish; the second lies in the line
$L=\Q(1,1,1,1,1,1)$. Both are $G$-stable. Indeed $G$ permutes the
three pairs, and its signs on the three oriented projection vectors
must all agree, by uniqueness of their relation.
Consequently the orbit of $f$ lies in $V$ or in $L$, contradicting
$\dim\langle Gf\rangle=4$.

Finally let $m=12$. The projection vectors are distinct even up to
sign. The four orbits of $\langle\rho\rangle$ pair under negation,
so orient coordinates to obtain
\[
 u,\rho u,\rho^2u,\quad v,\rho v,\rho^2v.
\]
A signed sum of one triple is either zero, when all its signs agree,
or is $\pm2$ times one of that triple's vectors. Nonzero sums from
the two triples cannot cancel, since their vectors are distinct up
to sign. Therefore every signed vector in $W$ is constant on each
triple. All of $Gf$ lies in this two-dimensional subspace, again a
contradiction. This excludes $r=4$ and proves the lemma.
\end{proof}
This recovers the degree-twelve rank conclusion in the Dodson--Zoller
classification, summarized in \cite[\S2.2.2]{Borselli}. The rank in
that CM-type convention is $\dim\MT(A)=r+1$. The proof above does
not require the enumeration of the possible Galois groups.

\begin{theorem}[All simple CM sixfolds]\label{p:CMcomplete}
For every simple CM abelian sixfold $A$ and every $n\ge1$, the
rational Hodge conjecture holds for $A^n$.
On $A$ itself the only possible Hodge-dimension polynomials are
\[
 P_A(t)=(1+t)^6\quad\hbox{or}\quad (1+t)^6+2t^3.
\]
In the first case all classes on every power are divisor-generated.
In the second case divisors together with pullbacks of the algebraic
Weil plane, under $K$-linear sum maps, generate them.
\end{theorem}
\begin{proof}
Lemma~\ref{p:CMsixrank} gives Hodge rank five or six. Apply
Corollary~\ref{p:CMrankfive}, including its invariant-generation
argument on all powers. Its rank-five Weil plane is algebraic by
the split trace polarization of Lemma~\ref{p:CMsplit} and the
sixfold theorem of Markman, not merely by the rank computation.
\end{proof}

\begin{remark}[Independent finite cross-check]
The rank-four exclusion also has a small exact check. Normalize
$f=(1,\ldots,1)$ and complete it to four independent signed orbit
vectors. Distinct Hodge weights make the six columns distinct among
the eight vectors $(1,a,b,c)$, $a,b,c\in\{1,-1\}$. The two omitted
vertices have Hamming distance one, two, or three. Up to cube
symmetry these are all possibilities. For their row spaces, exact
signed-permutation stabilizer computations give, respectively,
\[
\begin{array}{c|c|c|c}
\text{omitted distance}&|\operatorname{Aut}_{\pm}(W)|&
 \text{coordinate orbit sizes}&
 \dim\langle\operatorname{Aut}_{\pm}(W)f\rangle\\\hline
1&96&6&3\\
2&32&4+2&4\\
3&96&6&1
\end{array}
\]
The middle model is not transitive; the others cannot be generated
by the orbit of $f$. The script \texttt{verify\_cm\_sixfold\_rank.py}
checks all $6!2^6$ signed permutations in each model using rational
linear algebra. This verifies the rank lemma independently of
listing Galois groups, but is not used in place of its proof.
\end{remark}

\begin{corollary}[All abelian Galois CM fields of degree twelve]
\label{p:CMabelian}
Let $A$ be a simple CM sixfold. If $K=\End^0(A)$ is Galois and
abelian over $\Q$, the Hodge conjecture holds for every power of $A$.
\end{corollary}
\begin{proof}
We show $\dim\Hg(A)\ge5$ and apply Corollary~\ref{p:CMrankfive}.
For an abelian Galois CM field, a CM type is primitive precisely when
its stabilizer under translation by $\operatorname{Gal}(K/\Q)$ is
trivial: a nontrivial stabilizer would induce the type from its fixed
CM subfield. The stabilizer cannot contain complex conjugation.
The Hodge dimension is the number of nonzero odd Fourier components
of the type; equivalently it is the dimension of the orbit span of
its signed vector. This is also the character description in
\cite[\S1]{Yanai}, with the scalar component omitted when passing
from $\MT$ to $\Hg$.

The only abelian groups of order twelve are $C_{12}$ and $C_2\times C_6$.
For $C_{12}=\langle\rho\rangle$, conjugation is $\rho^6$.
Write the signed vector as $f_0,\ldots,f_5\in\{\pm1\}$ with
$f_{j+6}=-f_j$, and put $F(T)=\sum_{j=0}^5 f_jT^j$.
The rational odd-character factors are
$T^6+1=\Phi_4(T)\Phi_{12}(T)$, of dimensions two and four.
The factor $\Phi_4=T^2+1$ cannot divide $F$, since this would require
$f_0-f_2+f_4=f_1-f_3+f_5=0$, impossible for sums of three odd integers.
If $\Phi_{12}=T^4-T^2+1$ divided $F$, then
\[
 F(T)=(aT+b)(T^4-T^2+1),\qquad a,b\in\{\pm1\}.
\]
The full signed vector would be periodic with period four, contradicting
primitivity. Thus the Hodge dimension is six.

For $C_2\times C_6$, choose the first factor to be complex conjugation
and a generator $\rho$ of the complementary $C_6$. The six signs
$f_j$ are arbitrary, and the rational Fourier factors of their cyclic
orbit are $\Phi_1,\Phi_2,\Phi_3,\Phi_6$, of dimensions $1,1,2,2$.
If $\Phi_3$ divides $F$, the three sums $f_j+f_{j+3}$ are equal.
If their common value is $\pm2$, the type is constant; if it is zero,
translation by $c\rho^3$ stabilizes it. Both violate primitivity.
If $\Phi_6$ divides $F$, the differences satisfy
\[
 (f_0-f_3,f_1-f_4,f_2-f_5)=(t,-t,t),\quad t\in\{-2,0,2\}.
\]
For $t=0$ the vector has period three; for $t=\pm2$ it has period two.
Again primitivity is violated. Finally $F(1)=F(-1)=0$ cannot both
hold, because it would make each sum of three even-index or odd-index
signs zero. Both two-dimensional components and at least one
one-dimensional component therefore occur. This gives Hodge dimension
at least five, as required.
\end{proof}

\begin{corollary}[All Galois CM fields of degree twelve]
\label{p:CMgalois}
Let $A$ be a simple CM sixfold whose endomorphism field $K$ is Galois
over $\Q$. The Hodge conjecture holds for every power of $A$.
If its Galois group is nonabelian, $\Hg(A)=U_K$ and all these
Hodge classes are divisor-generated.
\end{corollary}
\begin{proof}
The abelian case is Corollary~\ref{p:CMabelian}. Among the groups of
order twelve, the nonabelian groups having a central involution are
$C_2\times S_3$ and
\[
 G=\langle a,b:a^3=b^4=1,\ bab^{-1}=a^{-1}\rangle.
\]
The other nonabelian group, $A_4$, has trivial center and cannot have
a central complex conjugation. In the second displayed group the
conjugation is $c=b^2$; in the first it generates the $C_2$ factor.
A primitive CM type has trivial right-translation stabilizer: a
nontrivial stabilizer would make it induced from its fixed CM subfield.
We prove that every such type has Hodge rank six.

For $C_2\times S_3$, the signed type is a function
$f:S_3\to\{1,-1\}$. Its translates span a submodule of the rational
regular representation
\[
 \Q[S_3]\simeq\Q\oplus\Q\oplus M_2(\Q).
\]
Changing $f$ to $-f$, assume its negative set $S$ has size at most
three. Size zero is imprimitive. If $|S|=1$, the trivial and sign
components of $f$ are nonzero, and its matrix component is
$-2\rho(s)$, hence invertible. Thus its translates span all six
dimensions. If $S=\{s,t\}$ and $s^{-1}t$ has order two, right
translation by $s^{-1}t$ fixes $S$, contradicting primitivity.
Otherwise $s^{-1}t$ has order three, so $s,t$ have the same parity.
Both one-dimensional components are nonzero and
$-2\rho(s)(1+\rho(s^{-1}t))$ is invertible: the eigenvalues of
$\rho(s^{-1}t)$ are the two primitive cube roots of unity.
Again the rank is six.

If $|S|=3$ and all its elements have the same parity, its type is
imprimitive. Otherwise $S$ has two elements of one parity and one
of the other. Exactly two of the three right transpositions pair
two elements inside $S$. The third transposition $t$ therefore
exchanges $S$ with its complement. Right translation by $ct$ fixes
the full CM type, also contradicting primitivity. This exhausts
the $C_2\times S_3$ case.

For the second group the odd group algebra is
\[
 \Q[G]/(c+1)\simeq\Q(\sqrt{-1})\oplus(-1,-3)_\Q.
\]
One verifies this directly by sending $(a,b)$ to $(1,\sqrt{-1})$
in the first summand, and to $((-1+j)/2,i)$ in the second, where
$i^2=-1$, $j^2=-3$, $ij=-ji$. These are surjections onto distinct
simple algebras and their dimensions sum to six. The quaternion
summand is division, since it is definite over $\R$.
Write the signed group-algebra element using the transversal
$1,a,a^2,b,ab,a^2b$ as
\[
 x=\sum_{r=0}^2u_ra^r+\sum_{r=0}^2v_ra^rb,
 \qquad u_r,v_r\in\{1,-1\}.
\]
Its first component is $\sum u_r+\sqrt{-1}\sum v_r$, which is
nonzero because both sums are odd. Its quaternion component vanishes
only if each triple $(u_0,u_1,u_2)$ and $(v_0,v_1,v_2)$ is constant:
use the direct sum $\Q(a)\oplus\Q(a)i$ and the sole relation
$1+a+a^2=0$. Such a type is invariant under right translation by
$a$ and is imprimitive. For a primitive type both division-algebra
components are nonzero, hence invertible. Its translates span the
whole six-dimensional odd algebra. This proves the asserted Hodge
rank and, by Proposition~\ref{t:CM}, divisor generation on all powers.
\end{proof}

\begin{example}[Exact matrix certificate for the trace construction]
Let $F=\Q(\theta)$ with
$f(T)=T^6+T^5-5T^4-4T^3+6T^2+3T-1$; one can take
$\theta=2\cos(2\pi/13)$. This is the real sextic subfield of
$\Q(\zeta_{13})$. For $K=\Q(i)F$ and $b'=1/f'(\theta)$ the trace
Hermitian form, in the $\Q(i)$-basis $1,\theta,\ldots,\theta^5$, is
\[
 \left(\Tr_{F/\Q}\frac{\theta^{u+v}}{f'(\theta)}\right)_{u,v=0}^5
 =\begin{pmatrix}
 0&0&0&0&0&1\\
 0&0&0&0&1&-1\\
 0&0&0&1&-1&6\\
 0&0&1&-1&6&-7\\
 0&1&-1&6&-7&27\\
 1&-1&6&-7&27&-35
 \end{pmatrix}.
\]
Its determinant is $-1$ and its upper-left $3\times3$ block is zero.
Choose the CM type by the positivity signs of $i/f'(\theta)$; after
clearing denominators this trace form polarizes the associated CM
abelian variety. The example verifies the split-form construction.
For an arbitrary prescribed balanced CM type, the ordering step in
Lemma~\ref{p:CMsplit} is essential; this particular $\theta$ is not
asserted to fit every type.
\end{example}

\begin{example}\label{ex:CM13}
Let $C$ be the smooth projective curve $y^2=x^{13}-1$ and $A=J(C)$.
There are $14$ branch points of the hyperelliptic map including infinity,
so $g(C)=6$. The automorphism $(x,y)\mapsto(\zeta_{13}x,y)$ acts on the
basis $x^{a-1}dx/y$, $1\leq a\leq6$, with eigenvalues $\zeta_{13}^a$.
It gives an embedding $K=\Q(\zeta_{13})\hookrightarrow\End^0(A)$ and
CM type $\Phi=\{1,2,3,4,5,6\}$ modulo $13$.
In the basis of six conjugate pairs, six Galois translates of the signed
CM vector have matrix
\[
 M=\begin{pmatrix}
1&1&1&1&1&1\\
1&1&1&-1&-1&-1\\
1&1&-1&-1&1&1\\
1&-1&-1&1&-1&-1\\
1&-1&1&-1&-1&1\\
1&-1&1&-1&1&-1
\end{pmatrix},\qquad\det(M)=-160.
\]
The row with multiplier $a$ has entry $+1$ at column $b$ if
$ab\bmod13\in\Phi$, and $-1$ otherwise; an inverse-multiplier convention
gives the same orbit span. Exact elimination gives the displayed nonzero
determinant, so the Hodge group is the full six-dimensional torus.
Its rational centralizer is $K$ (over $\C$ the twelve distinct character
spaces have diagonal centralizer). Therefore $\End^0(A)=K$ and $A$ is
simple. Proposition~\ref{t:CM} applies, with dimensions
$(1,6,15,20,15,6,1)$ and six rational divisor generators.
\end{example}
\begin{remark}
For a proper CM torus the vanishing of the total character must be tested
on that smaller torus, not on $T_K$. Additional monomials may survive.
Their exact degrees, rational dimensions, and algebraic representatives
must be supplied before a degenerate type is counted as solved.
Absolute Hodge classes are not thereby algebraic. Andre's reduction to
split Weil classes also does not keep all auxiliary varieties in dimension
six; see \cite[Theorem 1.4]{MarkmanSurvey}.
\end{remark}

\begin{example}[A degenerate CM sixfold with algebraic exceptional classes]
Let $F=\Q(\zeta_{13}+\zeta_{13}^{-1})$ and $K=F(i)$.
Then $[K:\Q]=12$ and $\mathrm{Gal}(K/\Q)=C_6\times C_2$,
where complex conjugation is the second factor. Order the real
embeddings of $F$ cyclically under a generator of $C_6$. Choose the
CM type whose signs on the six embeddings of $i$ are
\[
 f=(1,1,-1,1,-1,-1).
\]
The complex torus $A=\C^6/\Phi(\mathcal O_K)$ is polarizable:
choose $b\in F$ with prescribed signs by weak approximation so that
$ib$ defines a positive CM Riemann form. Clear denominators in
\[\psi(x,y)=\Tr_{K/\Q}(ib\,x\overline y).\]
The signed CM matrix consists of the six cyclic shifts of $f$.
Its rows sum to zero. A maximal minor is
\[
 \det\begin{pmatrix}
 1&1&-1&1&-1\\
 1&-1&1&-1&-1\\
 -1&1&-1&-1&1\\
 1&-1&-1&1&1\\
 -1&-1&1&1&-1
 \end{pmatrix}=-16.
\]
Thus its rational kernel is exactly $\Q(1,1,1,1,1,1)$ and its rank
is five. All twelve restricted eigencharacters are distinct: their
pairwise differences cannot be multiples of this all-ones vector.
The commutant is therefore $K$, which proves $\End^0(A)=K$ and simplicity.
Its Hodge group is $\ker(N_{K/\Q(i)}:U_K\to U_{\Q(i)})$.
The finite test gives $n_4=0$, $n_6=2$, and
\[
 P_A(t)=(1+t)^6+2t^3,
 \qquad (1,6,15,22,15,6,1).
\]
The real field $F$ contains $\Q(\sqrt{13})$, so
Lemma~\ref{p:quartic-scale} applied to $\Q(i,\sqrt{13})$ makes the
two exceptional $\Q(i)$-Weil classes algebraic.
Divisors and these classes exhaust $A$; Proposition~\ref{p:coupled-det}
with $e=6$, $q=1$ additionally proves algebraicity on every power.
This is a degenerate CM example, not merely a full-torus example.
\end{example}

\part{The ten nonsimple rows}
\subsection*{Product strata and the source of mixed classes}
The following sections retain each original row rather than identifying
an arbitrary product with its generic example. If a product $A\sim\prod B_i$
has the actual joint Hodge group $\prod\Sp(H^1(B_i))$, then
Lemma~\ref{p:fft} proves divisor generation on all its powers.
For repeated factors, first collect isogeny classes; the group acts
diagonally on their copies, and the invariant calculation includes
Poincare divisors between those copies.
If one already knows algebraicity on every power of the factors and
has verified a product Hodge-group equality, the Kunneth invariant
algebra is the tensor product of the factor invariant algebras, hence
is algebraic. Neither product-group equality nor algebraicity on all
powers follows merely from the Hodge conjecture for the individual factors.
The symplectic examples use jointly very general tuples as in
Proposition~\ref{p:products}. The CM and unitary examples verify their
different joint groups separately.
\section{N1: The \texorpdfstring{$5+1$}{5+1} case (\provedstatus)}\label{sec:N1}

\status{Proved for every sixfold in the $5+1$ row. The all-powers
strengthening is not asserted for the entire row.}
\begin{theorem}
Let $A\sim B\times E$ with $B$ a simple fivefold and $E$ elliptic.
Every rational Hodge class on $A$ is algebraic. Its Hodge ring is generated
by divisors, except that a rational two-dimensional Weil space in
codimension three must be added if $E$ has CM by a field acting on $B$
with signature $(3,2)$, up to interchange.
\end{theorem}
\begin{proof}
Apply Lemma~\ref{p:prime-ell} with $g=5$ and isogeny invariance.
In the exceptional case the $k$-action on the elliptic factor is oriented
so the combined signature is $(3,3)$. Its rank-one summand permits
Lemma~\ref{p:oddscale}, hence Markman's discriminant-$-1$ theorem.
This step proves the algebraicity of the mixed classes, rather than
inferring it from the individual factors.
\end{proof}
\begin{example}
Choose $(B_1,\ldots,B_{2})$ jointly very generally in
$\prod_i\mathcal A_{d_i}(N)$ with $d_i=(5,1)$,
as constructed in Proposition~\ref{p:products}. Then each factor is
simple with $\End^0(B_i)=\Q$, and the factors are pairwise nonisogenous.
Their product is an actual sixfold of this partition, with
$\End^0(A)=\Q^{2}$ and the proved joint group above.
The Hodge-space dimensions in codimensions $0$ through $6$ are
$(1,2,2,2,2,2,1)$.
An explicit monomial basis over $\Q$ on $A$ is
$\theta_1^{a_1}\cdots\theta_{2}^{a_{2}}$ with
$0\leq a_i\leq d_i$; every such monomial is an intersection of divisors.
On powers add Poincare divisors between the copies of each individual
$B_i$. There are no cross-factor divisors or higher mixed invariants in
this verified jointly generic example.
\end{example}
\begin{example}[A genuinely exceptional $5+1$ product]
Let $k=\Q(i)$ and take a very general polarized fivefold $B$ with
$k$-Hermitian datum $(k^5,\operatorname{diag}(1,1,1,-1,-1))$.
The unitary PEL domain has dimension six; its dense derived monodromy
and nonzero determinant circle give $\Hg(B)=U(3,2)$ and
$\End^0(B)=k$. Thus $B$ is simple. Let $E:y^2=x^3-x$ with the opposite
$k$-action. Then $A=B\times E$ has signature $(3,3)$ and determinant
$-1$. Lemma~\ref{p:prime-ell} gives
\[
 P_A(t)=(1+t+\cdots+t^5)(1+t)+2t^3,
 \qquad (\dim B^p(A))_{p=0}^6=(1,2,2,4,2,2,1).
\]
Products of the two polarization classes span the divisor part;
the remaining middle plane is $W_k(A)$, algebraic by the split Weil
theorem. This example explicitly exhibits the mixed classes absent
from the jointly symplectic example.
\end{example}
\section{N2: The \texorpdfstring{$4+2$}{4+2} case (\provedstatus)}\label{sec:N2}

\status{All cases with a CM fourfold are now covered, using the new
external CM theorem for the CM--CM products and the earlier independent
strata for a non-CM surface. The other previously proved strata and
weighted covering families remain covered. For the residual non-CM
fourfold with primitive quartic CM field and a matching CM surface,
Theorem~\ref{p:N2auxiliary} constructs a nonzero mixed cycle from
the auxiliary-spin and CM inputs. This completes the sixfold row;
the residual models also satisfy the all-powers conclusion.}

\begin{corollary}[All CM--CM products in N2]\label{p:N2CM-new}
Assume the CM input in Theorem~\ref{p:openai-inputs}.
If $B$ is a simple CM fourfold and $S$ a simple CM surface, then
the Hodge conjecture holds on every power of $B\times S$.
In particular every coupled case in Proposition~\ref{p:N2CMstructure}
is algebraic; no independence of the two Hodge groups is required.
Together with Proposition~\ref{p:N2strata}, this settles every
sixfold in N2 whose fourfold factor is CM.
\end{corollary}
\begin{proof}
The product $B\times S$ is CM, even when its Hodge torus is a proper
subgroup of $\Hg(B)\times\Hg(S)$.
Lemma~\ref{p:CMproducts-new} applies to each of its powers.
If instead $S$ is not CM, Proposition~\ref{p:N2strata} already
proves the Hodge conjecture for the sixfold. These two alternatives
exhaust surfaces and prove the final assertion.
\end{proof}
\begin{example}[A CM fourfold times a CM surface]\label{ex:N2CM-new}
Let $L=\Q(\zeta_{15})$ and take the CM type with embedding
exponents $\Phi=\{1,2,4,7\}$ modulo $15$.
It selects one embedding from each conjugate pair.
Its multiplicative stabilizer in $(\Z/15\Z)^\times$ is trivial:
any stabilizing element belongs to $\Phi$, and direct multiplication
excludes $2,4,7$. Since $L/\Q$ is Galois, the type is primitive.
The standard CM construction therefore gives a simple fourfold
$B=\C^4/\Phi(\mathcal O_L)$.
For its polarization, choose $\delta\in L$ with
$\bar\delta=-\delta$ and $\operatorname{Im}\phi(\delta)>0$
for every $\phi\in\Phi$, by weak approximation.
An integral multiple of $\Tr_{L/\Q}(\delta x\bar y)$ on
$\mathcal O_L$ is a Riemann form. Primitivity of the CM type
gives simplicity.
Take $S=J(y^2=x^5-1)$, the simple CM surface with field
$\Q(\zeta_5)$. Corollary~\ref{p:N2CM-new} proves the Hodge
conjecture for every $(B\times S)^n$ by the external CM theorem.
This argument does not require computing the joint torus or assuming
that its exceptional classes are divisor-generated.
\end{example}

\begin{proposition}
Let $A\sim B_1\times\cdots\times B_{2}$ with dimensions
$(4,2)$. Assume the actual joint Hodge group is the
independent product of the full symplectic groups on these factors.
Then all Hodge classes on all powers of $A$ are divisor-generated,
and $P_A(t)=(1+t+\cdots+t^{4})(1+t+t^2)$.
\end{proposition}
\begin{proof}
Use Lemma~\ref{p:isogeny} to replace the isogenous variety by the product.
Under the stated group equality,
\[
 \left(\bigwedge^\bullet\bigoplus_i H^1(B_i)^{\oplus n}\right)^{\Hg(A)}
 =\bigotimes_i\left(\bigwedge^\bullet H^1(B_i)^{\oplus n}\right)^{\Sp(H^1(B_i))}.
\]
Each factor is generated by degree-two contractions by Lemma~\ref{p:fft}.
These are divisor classes on the product, so every invariant is algebraic.
On one copy the generators are the polarization powers through the
dimension of each factor, giving the displayed polynomial.
\end{proof}
\begin{proposition}[Two further complete strata]\label{p:N2strata}
Let $A=B\times S$, with $B$ a simple fourfold and $S$ a simple surface.
The Hodge conjecture holds for $A$ if either $S$ is not CM, or the center
of $\End^0(B)$ is totally real or imaginary quadratic.
\end{proposition}
\begin{proof}
If $S$ is not CM, \cite[Proposition (2.4), Lemma (3.4)]{MZ} and
$\Hom(S,B)=0$ give $\Hg(A)=\Hg(B)\times\Hg(S)$.
If $S$ is CM, its Hodge group is a $\Q$-simple rank-two torus.
The connected center of $\Hg(B)$ embeds in the norm-one torus of the
CM center of $\End^0(B)$; it is trivial for a totally real center and
has rank at most one for an imaginary quadratic center.
Goursat's lemma therefore excludes a common nontrivial quotient with
$\Hg(S)$, proving the same product equality. The exterior Kunneth
invariants are consequently tensor products of factor Hodge classes.
The fourfold theorem in \cite[Corollary 1.6.1]{Markman} and the
prime-dimensional surface theorem make these factor classes algebraic.
Their external products prove the assertion.
\end{proof}
\begin{proposition}[A splitting-field separation test]\label{p:N2galois}
Let $A=B\times S$, where $B$ is a simple fourfold and $S$ a simple
CM surface with field $F$. Suppose the center $K$ of $\End^0(B)$
is CM. Let $L_F,L_K\subset\overline\Q$ be their normal closures.
If $L_F\not\subset L_K$, then
$\Hg(A)=\Hg(B)\times\Hg(S)$ and the Hodge conjecture holds for $A$.
In particular this holds when $K/\Q$ is multiquadratic, including
every biquadratic quartic center.
\end{proposition}
\begin{proof}
The torus $\Hg(S)=U_F$ is $\Q$-simple of rank two, with splitting
field $L_F$ and Galois group $C_4$ or $D_4$.
The connected center $Z$ of $\Hg(B)$ embeds into $U_K$ and is
therefore split by $L_K$. By Goursat, any failure of independence
would make $U_F$ an isogeny quotient of $Z$. Its rational character
module would then be split by $L_K$, forcing $L_F\subset L_K$.
This contradiction gives independence; the fourfold and surface
theorems imply algebraicity of the tensor-product Hodge ring.
If $K/\Q$ is multiquadratic, every quotient of its Galois group has
exponent at most two, and thus cannot be $C_4$ or $D_4$.
\end{proof}
\begin{proposition}[All residual non-CM-fourfold groups in N2]
\label{p:N2nonCMstructure}
Let $A=B\times S$ be an N2 case with $B$ not CM and not already
covered by Propositions~\ref{p:N2strata}--\ref{p:N2galois}.
Then $\End^0(B)=K$ is a primitive quartic CM field,
$\End^0(S)=F$ is a primitive quartic CM field, and the two
$K$-signatures on $B$ are $(2,0),(1,1)$ after orientation.
In particular $\Hg(B)$ is the full $K$-unitary group.
Orient the two complex coordinates on $U_F$ at the CM type of $S$.
There are precisely two alternatives:
\begin{enumerate}
\item There is no rational torus morphism $\nu:U_F\to U_K$
whose complex coordinate expression is $(u,v)\mapsto(uv,u/v)$,
up to conjugating the second $K$ embedding. Then the joint group
is the independent product and the Hodge conjecture holds.
\item Such a morphism exists. It is an isogeny, and the joint group is
\[
 \{(g,s)\in\Hg(B)\times U_F:\det_K g=\nu(s)\}.
\]
Its Hodge ring and polynomial are exactly those in
Lemma~\ref{p:quartic-mixed}. Algebraicity is equivalent to the
existence of one codimension-two algebraic cycle with nonzero
projection to the four-dimensional space $Q$ in that lemma.
\end{enumerate}
The torus-morphism condition is a finite Galois-equivariance test
on the displayed integral matrix; it is not an algebraicity assumption.
\end{proposition}
\begin{proof}
The previous strata leave a larger CM center on the non-CM
fourfold and a simple CM surface. The dimension formula forces
the fourfold's endomorphism algebra to be a quartic field $K$;
a degree-eight field would make $B$ CM. Biquadratic $K$ was
already separated by Proposition~\ref{p:N2galois}, so $K$ is
primitive. Two pure signatures would make the Hodge circle,
and hence its rational group, central and would make $B$ CM.
Two balanced signatures would put the group in $\SL_2^2$;
the standard--dual intertwiners would contradict the endomorphism
commutant $K$. Only the stated mixed signatures remain. The
derived factors are $\SL_2^2$, not a diagonal group, because
only one block has nontrivial adjoint Hodge circle. The determinant
cocharacter and its conjugates span $U_K$, giving the full group.

Both $U_K$ and $U_F$ are rationally simple two-dimensional tori.
If the joint group is not independent, Goursat gives an isogeny
between these determinant tori. On rational cocharacter spaces
it is a Galois-equivariant isomorphism $f$ sending $(1,1)$ to
$(2,0)$. The Galois actions are irreducible signed permutation
actions of $C_4$ or $D_4$ on a plane. Their invariant positive
quadratic forms are unique up to a scalar, so $f$ is a similarity
for the standard Euclidean forms, with multiplier $2$.
Orthogonality to $(1,1)$ consequently gives
\[
 f(1,-1)=(0,\pm2),\qquad
 f=\begin{pmatrix}1&1\\1&-1\end{pmatrix}
 \quad\text{up to the stated conjugation}.
\]
This integral Galois-equivariant matrix gives the claimed rational
torus morphism. Conversely, if it exists, the displayed group is
connected, has dimension $\dim\Hg(B)$, and contains the Hodge
circle. Surjectivity of the actual group onto $\Hg(B)$ proves
equality. If it does not exist, Goursat forces independence;
the factor Hodge classes are already algebraic.

In the coupled case Lemma~\ref{p:quartic-mixed} leaves exactly
$Q$. Its projection is algebraic: use the determinant projector
on $H^2(B)$ constructed in Lemma~\ref{p:cyclic-transfer},
and the degree-two projector on $S$. For a general
$a\in F$, pullback by $a$ has four distinct eigenvalues on the
four mixed character lines of $H^2(S)$, namely the products
obtained by choosing one embedding from each conjugate pair.
The Galois group acts transitively on these four products, so
their polynomial over $\Q$ is irreducible of degree four.
The same four eigenvalues occur once each on $Q$.
Thus the rational algebra generated by this algebraic operator
makes $Q$ a one-dimensional vector space over a quartic field.
Every nonzero algebraic projection generates all of $Q$ by
algebraic correspondences, proving the criterion.
\end{proof}

\begin{proposition}[Rank-three CM fourfolds are independent of surfaces]
\label{p:N2CMrankthree}
Let $B$ be a simple CM fourfold with $L=\End^0(B)$.
Its Hodge rank is three or four. In rank three, $\Hg(B)$ is a
$\Q$-simple torus. Consequently, for every simple CM surface $S$,
\[
 \Hg(B\times S)=\Hg(B)\times\Hg(S),
\]
and the Hodge conjecture holds for $B\times S$. In this case
\[
 P_{B\times S}(t)=(1+t)^6+2t^2(1+t)^2
 =1+6t+17t^2+24t^3+17t^4+6t^5+t^6.
\]
\end{proposition}
\begin{proof}
The eight weights of $H^1(B)_\C$ are distinct because its Hodge
commutant is $L\otimes\C$. Evaluating them on a basis of signed
CM-type orbit vectors gives $8\leq2^r$, where $r=\dim\Hg(B)$,
so $3\leq r\leq4$.
Suppose $r=3$. Lemma (3.7) of \cite{MZ} gives an imaginary
quadratic $k\subset L$ with $\Hg(B)=\ker(N_{L/k}:U_L\to U_k)$.
Orient the four embedding coordinates over one embedding of $k$.
The signed CM vector has two positive and two negative entries,
and the Galois action is by permutations with a common sign.
Its permutation image $P\subset S_4$ is transitive. On the
hyperplane $\sum x_i=0$ the action is the standard permutation
representation, possibly twisted by a one-dimensional character.

The transitive subgroups of $S_4$ are $C_4,V_4,D_4,A_4,S_4$.
For the first three, an orbit of a vector with two positive and
two negative entries spans at most a plane: in the square model
the two positive positions are either adjacent or opposite.
The $V_4$ case is a subgroup of this same model. Thus rank three
forces $P=A_4$ or $S_4$. Both are doubly transitive, so their
standard three-dimensional representation is irreducible over
$\Q$, as is its character twist. This proves rational simplicity.
The rank-two torus $\Hg(S)$ is also rationally simple and cannot
be a quotient of this rank-three torus. Goursat gives independence.
The Hodge ring is the tensor product of the factor Hodge rings;
the fourfold theorem of \cite{Markman} and the surface theorem
give its algebraicity.
For the norm-kernel fourfold, the only non-pairing weights on one
copy are the two $k$-volumes. Thus $P_B(t)=(1+t)^4+2t^2$;
multiplication by $P_S(t)=(1+t)^2$ gives the stated polynomial.
\end{proof}

\begin{proposition}[All CM-fourfold couplings in N2]
\label{p:N2CMstructure}
Let $B$ be a simple CM fourfold, $L=\End^0(B)$, and let $S$ be
a simple CM surface with field $F$. If their joint Hodge group
is not the independent product, then $\Hg(B)=U_L$ and there is
a primitive quartic CM subfield $K\subset L$ such that the
$K$-signatures on $B$ are $(2,0),(1,1)$ after orientation.
There is a rational torus isogeny $\nu:U_F\to U_K$ with complex
coordinates $(u,v)\mapsto(uv,u/v)$, and the actual joint group is
\[
 H=\{(b,s)\in U_L\times U_F:N_{L/K}(b)=\nu(s)\}.
\]
Conversely, these field, signature, and equivariant-isogeny data
with $\Hg(B)=U_L$ imply this group equality. Put
\[
 T=\bigwedge\nolimits_K^2H^1(B,\Q),\qquad
 Q=\bigl(T\otimes H^2(S,\Q)\bigr)\cap H^{2,2}(B\times S),
\]
using the rational-summand convention of Lemma~\ref{p:cyclic-transfer}.
Then $\dim_\Q Q=4$, divisors and $Q$ generate the Hodge ring, and
\[
 P_{B\times S}(t)=(1+t)^6+4t^2(1+t)^2
 =1+6t+19t^2+28t^3+19t^4+6t^5+t^6.
\]
The Hodge conjecture is equivalent to the existence of a
codimension-two algebraic cycle whose projection to $Q$ is nonzero.
\end{proposition}
\begin{proof}
Proposition~\ref{p:N2CMrankthree} leaves only $\Hg(B)=U_L$.
Since $U_F$ is rationally simple of rank two, a non-independent
joint torus projects isogenously onto $U_L$. On rational
cocharacter spaces it is the graph of a surjection
\[
 f:\Q^4\longrightarrow\Q^2,\qquad f(h_B)=h_S=(1,1),
 \quad h_B\in\{1,-1\}^4.
\]
This map is Galois-equivariant. Write $a_i=f(e_i)$.
Transitivity on the eight signed $L$ embeddings shows that all
$a_i$ are nonzero and that the multiset $\{\pm a_i\}_{i=1}^4$
is one orbit, each vector occurring with equal multiplicity.
The Galois image on the target is the standard signed-plane
group $\Gamma=C_4$ or $D_4$. Its nonzero orbits have four or
eight elements, so there are exactly two possibilities.

If all eight vectors are distinct, equivariance identifies the
signed permutation action on the source with the regular action
of $D_4$ on its eight elements; an element acting trivially on
the target must fix every signed source coordinate. The source
is therefore the negative-central part of the regular module.
Writing
\[
 R=\begin{pmatrix}0&-1\\1&0\end{pmatrix},\qquad
 D=\begin{pmatrix}1&0\\0&-1\end{pmatrix},
\]
identifies this four-dimensional module with $M_2(\Q)$ under
left multiplication. A signed CM vector corresponds to
\[
 M=\epsilon_1 I+\epsilon_2 R+\epsilon_3 D+\epsilon_4 RD
 =\begin{pmatrix}
 \epsilon_1+\epsilon_3&-\epsilon_2+\epsilon_4\\
 \epsilon_2+\epsilon_4&\epsilon_1-\epsilon_3
 \end{pmatrix},\qquad \epsilon_i\in\{1,-1\}.
\]
Its determinant is
$\epsilon_1^2+\epsilon_2^2-\epsilon_3^2-\epsilon_4^2=0$,
so it has rank one. Its left orbit consequently spans only two
dimensions, contradicting $\langle Gh_B\rangle=\Q^4$.

There are thus four distinct projection vectors, each repeated
twice. Orient and reorder the coordinates so that
\[
 (a_1,a_2,a_3,a_4)=(a,a,b,b),\qquad b\perp a,\quad |b|=|a|.
\]
The orbit is a square. The image of a signed vector is $n_1a+n_2b$,
where $n_i\in\{-2,0,2\}$ are its pair sums. If both sums are
nonzero, that signed vector is constant on each pair. The
pair-constant plane is Galois-stable, contradicting full Hodge
rank. Both sums cannot be zero because $f(h_B)=(1,1)$.
Hence, after changing orientations,
\[
 h_B=(1,1,1,-1),\qquad
 a=(1/2,1/2),\quad b=(1/2,-1/2).
\]
The four fibers of the map from signed embeddings to these
projection vectors form a Galois-stable block system. It defines
a degree-four subfield $K\subset L$. Conjugation interchanges
opposite blocks, so $K$ is CM; the associated rank-two signed
module is irreducible, so $K$ is primitive. The displayed CM
vector gives its signatures $(2,0),(1,1)$.

The norm $N_{L/K}$ is the pair-sum map on cocharacters. If
$\nu_*=\left(\begin{smallmatrix}1&1\\1&-1\end{smallmatrix}\right)$,
then $\nu_*f=(N_{L/K})_*$. Surjectivity of $f$ proves that
$\nu_*$ is Galois-equivariant. It is integral in these standard
cocharacter lattices and therefore defines the asserted torus
isogeny. Conversely, the displayed group $H$ is connected
because $\ker(N_{L/K})$ is connected: its complex coordinate map
is $(x_1,x_2,y_1,y_2)\mapsto(x_1x_2,y_1y_2)$.
It contains the joint Hodge circle and has dimension four.
Surjectivity of the actual group onto $U_L$ gives equality.

For the invariants, split each two-dimensional block in
Lemma~\ref{p:quartic-mixed} into its two independent CM characters.
The determinant-kernel invariants in one block are its two
opposite-character pairings, their product, and its two volume
lines. The same neutrality equations allow only one active
volume block. The pairings give $(1+t)^6$, and the four mixed
degree-four lines contribute $4t^2(1+t)^2$. They form $Q$ over
$\Q$ and generate the quotient by the divisor algebra.
The determinant projector on $B$ and the degree-two projector
on $S$ are algebraic. On $Q$, a general $a\in F$ has four
distinct eigenvalues in one Galois orbit, the four mixed products
of embeddings of $F$. Its minimal polynomial is irreducible
of degree four. Thus $Q$ is rank one over the quartic field
generated by this algebraic operator. Every nonzero algebraic
projection generates all of $Q$, proving the final criterion.
\end{proof}

The script \texttt{verify\_n2\_cm\_structure.py} checks all thirty
subgroups of $S_4$, all sixteen signed vectors in the regular
dihedral model, and the pair-sum alternatives. These finite
checks supplement the preceding proofs; they assert no cycle
existence. Together with Proposition~\ref{p:N2nonCMstructure},
the result leaves only the single mixed-cycle construction in
each residual N2 case, not an unclassified common torus quotient.

\begin{remark}[A finite-kernel obstruction to an automatic reduction]
In either coupled N2 model, projection of the joint group onto
$\Hg(B)$ has kernel $\mu_2$: its nontrivial element acts trivially
on $H^1(B)$ and as $-1$ on $H^1(S)$. Thus $H^1(S)$ is not an
object of the tensor category generated by $H^1(B)$ and Tate
structures. Even an all-powers theorem for $B$ does not make $S$
an available factor by tensor operations. The even-degree
mixed Hodge tensors still require an algebraic correspondence;
the projection criterion above records exactly that missing step.
\end{remark}

\begin{lemma}[The quartic mixed invariant calculation]\label{p:quartic-mixed}
Let $B$ be a fourfold and $P$ an abelian surface. Suppose that over
$\C$ their actual joint Hodge representation is
\[
 H^1(B)_\C=V_1\oplus V_1^*\oplus V_2\oplus V_2^*,\quad
 H^1(P)_\C=\C_u\oplus\C_{u^{-1}}\oplus\C_v\oplus\C_{v^{-1}},
 \qquad \dim V_j=2,
\]
with group
\[
 H=\{(g_1,g_2,u,v):\det g_1=uv,\ \det g_2=u/v\}
 \subset\GL_2^2\times\mathbf G_m^2.
\]
Suppose also that the sum $T_\C$ of the four determinant lines
$\bigwedge^2V_j,\bigwedge^2V_j^*$ descends to
$T\subset H^2(B,\Q)$. There is a four-dimensional rational space
\[
 Q=\bigl(T\otimes H^2(P,\Q)\bigr)\cap H^{2,2}(B\times P)
 \subset\B^2(B\times P)
\]
such that the Hodge ring is generated by divisors and $Q$. More precisely,
\[
 P_{B\times P}(t)=(1+t+t^2)^2(1+t)^2+4t^2(1+t+t^2)
 =1+4t+12t^2+14t^3+12t^4+4t^5+t^6.
\]
Thus the Hodge conjecture on this sixfold is equivalent to the
algebraicity of $Q$, not merely of a middle-degree Weil plane.
\end{lemma}
\begin{proof}
The derived group is $\SL_2^2$. On
$\bigwedge^\bullet(V_j\oplus V_j^*)$ its invariants are the
pairing powers $1,\theta_j,\theta_j^2$ of weight zero and the two
degree-two volumes of weights $d_j,d_j^{-1}$. A pairing times a
volume in the same block is zero. Write $m_j\in\{-1,0,1\}$ for
the chosen determinant exponent. The required $P$ weights are
$-(m_1+m_2)$ and $-(m_1-m_2)$, each of which can only be $-1,0,1$
on a single copy of $P$. Therefore both $m_j$ cannot be nonzero.
When both vanish the invariants are exactly the pairing monomials,
with polynomial $(1+t+t^2)^2(1+t)^2$. Otherwise there are four
possibilities: one $B$ volume and one coordinate of each of the
$u$ and $v$ pairs. They give the four lines of $Q_\C$ in degree four.
Each can only be multiplied by $1,\theta_{3-j},\theta_{3-j}^2$
from the inactive $B$ block, giving $4t^2(1+t+t^2)$.
Rational descent of the surjective multiplication map
$\D^\bullet\otimes Q\to\B^\bullet/\D^\bullet$
proves the claimed generation over $\Q$.
\end{proof}

\begin{proposition}[Tensor contraction for every residual mixed model]
\label{p:mixed-allpowers}
Let $A=B\times P$ be a residual coupled case of N2 or N3 in
Propositions~\ref{p:N2nonCMstructure}, \ref{p:N2CMstructure}, or
\ref{p:N3residual}; thus $B$ is simple and $P$ is a CM surface or a
product of the two specified CM elliptic curves. The following are
equivalent:
\begin{enumerate}
\item some codimension-two algebraic cycle on $A$ has nonzero
projection to the four-dimensional mixed space $Q$;
\item the Hodge conjecture holds for $A$;
\item the Hodge conjecture holds for $A^n$ for every $n\ge1$.
\end{enumerate}
\end{proposition}
\begin{proof}
The single-cycle criteria in the cited propositions give
$(1)\Longleftrightarrow(2)$, and $(3)\Longrightarrow(2)$ is immediate.
It remains to prove $(1)\Longrightarrow(3)$. Those criteria first make
all of $Q$ algebraic using field endomorphisms; fix this input.

We work in tensor powers of $H^1(A)_\C$; alternation will recover
cohomology of arbitrary powers. Kunneth projectors, factor permutations,
endomorphisms, cup products, and polarization contractions on abelian
varieties are algebraic correspondences. In particular an algebraic
class in $Q\subset H^2(B)\otimes H^2(P)$ can be separated into four
$H^1$ slots by addition-map pullback and degree-one projectors.
Inverse polarization contractions are algebraic by the Lefschetz
standard conjecture for abelian varieties. We may therefore use these
tensor operations on algebraic classes.

First take invariants under $D=\ker(\det_K:\Hg(B)\to U_K)$.
For non-CM $B$, $D_\C=\SL_2^2$. Its tensor invariants are generated
by evaluation pairings and the two block volumes and their duals.
For CM $B$, $D_\C=\mathbf G_m^2$: in each rank-two $K$ block its
characters are $z,z^{-1}$, with the two dual characters on the dual
block. Pairwise opposite-character tensors are exactly the individual
evaluation pairings, the block volume, and its dual. In this CM case the endomorphisms
in $L$ separate the two lines: applying their eigenspace projectors
to individual slots of the antisymmetric volume gives both ordered
opposite-character tensors. With these algebraic projectors the
listed tensors again generate all tensor invariants. The evaluation pairings in either
case are rational divisor classes after taking the full rational span.
An opposite pair of volumes in the same block can be replaced by the
determinant of evaluation pairings. Thus it suffices to handle volume
monomials with signed exponents $r,s\in\Z$ for the two blocks.

The remaining torus has coordinates $u,v$ on $H^1(P)$, and the block
volumes have characters $uv,u/v$. Neutrality requires residual $P$
character $u^{-(r+s)}v^{-(r-s)}$. A positive or negative block-one
volume together with one coordinate of each of the $u,v$ pairs is a
line of $Q_\C$; the analogous block-two tensors give the other two
lines. Take $|r|+|s|$ appropriate signed copies of these algebraic
$Q$ tensors. Their $B$ part is the desired volume monomial.
Their $P$ part has the correct total character, but may have extra
opposite $u$ or $v$ slots. Contract every such opposite pair with the
polarization. Each character line is one-dimensional and its pairing
with the opposite line is nonzero, so this yields a nonzero multiple
of the required residual tensor. Any additional opposite slots in a
given target tensor are supplied by divisor pairings.

This constructs every torus-neutral tensor from $Q$ and algebraic
pairings, even when it is not an ordinary exterior product of classes
pulled back from $A$. Character projections and tensor permutations
may be performed after scalar extension: endomorphism eigenspace
projectors are complex linear combinations of rational algebraic
correspondences. Consequently the complex span of rational algebraic
classes contains every invariant tensor. Alternating the slots into
any prescribed copies of $A$, and then descending the surjective
rational cycle-class span, proves the assertion on every $A^n$.
No implication from $\mathrm{HC}(A)$ to all powers is used outside
these explicitly verified group models.
\end{proof}

\begin{lemma}[A Fermat correspondence for weighted cyclic covers]
\label{p:cyclic-transfer}
Let $m\in\{5,8,10,12\}$ and let $\boldsymbol a=(a_1,a_2,a_3,a_4)$
satisfy $1\le a_i<m$, $\sum a_i\equiv0\pmod m$, and
$\gcd(m,a_1,a_2,a_3,a_4)=1$. Assume the four primitive multiplicities
\[
 d_j=-1+\sum_{i=1}^4\left\langle-\frac{ja_i}{m}\right\rangle,
 \qquad j\in(\Z/m\Z)^\times,
\]
have multiset $\{2,1,1,0\}$. For $t\in\C\setminus\{0,1\}$ let
$C_{m,\boldsymbol a,t}$ be the smooth projective model of
\[
 y^m=x^{a_1}(x-1)^{a_2}(x-t)^{a_3},
\]
and let $B_{m,\boldsymbol a,t}$ be the primitive $K=\Q(\zeta_m)$
isogeny factor of its Jacobian. This factor has dimension four;
the entire Jacobian need not have dimension four. Put
\[
 T=\bigwedge\nolimits_K^2 H^1(B_{m,\boldsymbol a,t},\Q)
 \subset H^2(B_{m,\boldsymbol a,t},\Q),
\]
where the inclusion denotes the rational summand comprising the four
same-embedding determinant lines. For every abelian surface $P$, all
rational Hodge classes in $T\otimes H^2(P,\Q)$ are algebraic on
$B_{m,\boldsymbol a,t}\times P$, for every allowed $t$.
For $m=5,8,12$ and $\boldsymbol a=(1,1,1,m-3)$ retain the shorthand
$C_{m,t},B_{m,t}$ used below.
\end{lemma}
\begin{proof}
Write $B=B_{m,\boldsymbol a,t}$ and $C=C_{m,\boldsymbol a,t}$.
Connectedness follows from the gcd assumption. Since every primitive
character is nontrivial at each of the four branch points, its
$H^1$ eigenspace has dimension two. The displayed multiplicity formula
is \cite[Lemma 2.7]{MoonenCyclic}. Orient the conjugate pairs so that
the pure block precedes the balanced block. The determinant cocharacter
$(2,0)$ and its Galois conjugates span $U_K$. Consequently the
determinant image of $\Hg(B)$ is $U_K$ at every parameter, even
when extra endomorphisms appear. In particular $T$ has no Tate
subquotient and is primitive for any polarization of $B$.

We also record a group dichotomy used below. Each derived projection
to a two-dimensional block is either trivial or $\SL_2$, and
Galois transitivity makes the two projections simultaneous. In the
nontrivial case a diagonal identification is impossible: the
adjoint Hodge circle is trivial on the pure-signature block and
nontrivial on the balanced block. Thus either $B$ is CM, or
$\Hg(B)_\C=\GL_2^2$ and $\End^0(B)=K$. This argument only
uses the signatures and works for any quartic CM field with these
signatures. No simplicity or endomorphism-equality assumption is
needed for the correspondence constructed here.

\emph{The quotient surface.} Let $\sigma(x,y)=(x,\zeta_m y)$.
On $C\times C$, the automorphisms
$r=(\sigma,\sigma^{-1})$ and $s(p,q)=(q,p)$ generate a dihedral
group of order $2m$. Its invariant function field has coordinates
\[
 a=x_1+x_2,\qquad b=x_1x_2,\qquad w=y_1y_2,
 \qquad w^m=b^{a_1}(b-a+1)^{a_2}(b-ta+t^2)^{a_3}.
\]
Indeed these functions are invariant. The Kummer extension they
define has degree $m$: the valuations along the three displayed
lines, together with the line at infinity, are $a_1,a_2,a_3,a_4$
modulo $m$, whose common gcd with $m$ is one. The function field
$\C(C\times C)$ has degree $2m^2$ over $\C(a,b)$, so the stated
field is the full dihedral invariant field.
Let $R_t$ be a smooth projective resolution of the quotient.
Homogenizing with $z$ gives a cyclic cover of $\mathbf P^2$
branched on the four lines
\[
 l_1=b,\quad l_2=b-a+z,\quad l_3=b-ta+t^2z,\quad l_4=z,
 \qquad w^m=l_1^{a_1}l_2^{a_2}l_3^{a_3}l_4^{a_4}.
\]
Set $X=(1-t)l_1$, $Y=t l_2$, $Z=t(t-1)l_4$.
Then $l_3=X+Y+Z$. After rescaling $w$ by a nonzero complex
constant, $R_t$ is birational to the normalization of
\[
 R_{m,\boldsymbol a}:\quad W^m=X^{a_1}Y^{a_2}Z^{a_4}(X+Y+Z)^{a_3}.
\]
This means the normalization of the cyclic cover of $\mathbf P^2$;
if $\sum a_i=dm$, the coordinate $W$ has weight $d$ in
$\mathbf P(1,1,1,d)$. The identification holds for every allowed $t$,
not merely at a special fibre.

\emph{Domination and known algebraic cycles.} Choose $\eta^m=-1$. The Fermat surface
$F_m^2:A^m+B^m+C^m+D^m=0$ dominates $R_{m,\boldsymbol a}$ by
\[
 [A:B:C:D]\longmapsto
 [A^m:B^m:C^m:\eta^{a_3} A^{a_1}B^{a_2}C^{a_4}D^{a_3}].
\]
The Fermat curve $F_m:a^m+b^m+c^m=0$ in turn gives a dominant
rational map $F_m\times F_m\dashrightarrow F_m^2$ by
\[
 ([a:b:c],[d:e:f])\longmapsto[af:bf:\eta cd:\eta ce].
\]
Both maps are generically finite: the first is finite over a general
point of the plane, and the two independent ratios $a/b,d/e$
show that the second has two-dimensional image. Resolving the
composition gives a generically finite morphism
$\widetilde{F_m\times F_m}\to R_t$, with source obtained by
point blowups.

The characters of $(\mu_m)^2$ in $H^1(F_m,\C)$ occur once and
are $(a,b)$ with $a,b,a+b\ne0$ in $\Z/m\Z$. Their rational
Galois orbits have at most $\varphi(m)=4$ elements. The rational
group-algebra idempotents therefore decompose $J(F_m)$, up to
isogeny, into abelian varieties of dimension at most two.
By Theorem~\ref{p:lowdim}(3), the Hodge conjecture holds for
$J(F_m)^2\times P$. The Abel--Jacobi maps give a surjection on
cohomology to $F_m\times F_m\times P$: they are isomorphisms in
degree one, and a polarization restricts to a nonzero class in
degree two. This surjection splits in rational polarizable Hodge
structures, so every Hodge class on the latter product lifts to a
Hodge class, and is algebraic. Point blowups of $F_m\times F_m$
add, after product with $P$, only Tate twists of $H^\bullet(P)$.
The blowup formula and generically finite push--pull now prove
the Hodge conjecture for $R_t\times P$.

\emph{An algebraic retraction, not just a Hodge isomorphism.}
Choose an Abel--Jacobi map followed by a $\sigma$-equivariant homomorphism
$J(C)\to B$ that realizes the primitive factor, and let
$a_B:C\times C\to B$ be its sum on the two coordinates.
On $T_\C$, the pullback $a_B^*$ is the injective antisymmetrization
\[
 \bigwedge^2 V_j\longrightarrow V_j\otimes V_j,
 \qquad v\wedge w\longmapsto v\otimes w-w\otimes v.
\]
The same-factor $H^2(C)$ terms vanish since $2j\not\equiv0\pmod m$.
The map $r$ acts trivially on this tensor, while $s$ acts by minus
tensor interchange on $H^1(C)\otimes H^1(C)$ and therefore also
fixes it. Quotient pushforward, composed with $a_B^*$, gives an
algebraic correspondence inducing an injection
$\alpha:T\hookrightarrow H^2(R_t,\Q)$.
For precision, resolve the quotient map equivariantly. The transfer
identity is $q^*q_*=\sum_g g^*$ modulo exceptional Tate summands,
so is $2m$ times the identity on $a_B^*T$; the Tate corrections
vanish because $T$ has no Tate subquotient.

There is an algebraic projector $p_T$ onto $T$ in $H^2(B)$.
To see this, choose a sufficiently general $k\in K$: the four
eigenvalues $\tau(k)^2$ can be separated from all mixed products
$\tau(k)\tau'(k)$, $\tau\ne\tau'$. Both sets are Galois stable,
so a rational polynomial in $k^*$ gives the projector. Clearing
denominators realizes $k$ by a rational algebraic correspondence;
the cohomological-degree projector is algebraic on an abelian variety.
Choose a $K$-compatible polarization, for example the one inherited
from the Jacobian. The norm-one elements of $K$ preserve its
polarization form. Their characters on $T$ are
$\pm2e_1,\pm2e_2$, distinct from the mixed-summand characters
$0,\pm e_1\pm e_2$; hence $T$ and its complement are orthogonal.
This argument remains valid at CM and nonsimple parameters.

Let $L_B$ be an ample Lefschetz operator. On primitive degree two,
the polarized adjoint of $\alpha$ is
\[
 \alpha^\dagger=p_T L_B^{-2}\,{}^t\!\alpha:
 H^2(R_t,\Q)\longrightarrow T.
\]
Here $L_B^{-2}:H^6(B)\to H^2(B)$ is algebraic by the Lefschetz
standard conjecture for abelian varieties; see
\cite[\S5]{MilneLefschetz}. The image of $T$ is
primitive on $R_t$, since its pairing with any ample class is a
map to a Tate structure and hence zero. Positivity of these
polarizations and injectivity of $\alpha$ make
$M=\alpha^\dagger\alpha|_T$ invertible. Its inverse is a rational
polynomial in $M$ by Cayley--Hamilton. Consequently
$\beta=M^{-1}\alpha^\dagger$ is an algebraic correspondence with
$\beta\alpha=\mathrm{id}_T$.
For a Hodge class $h\in T\otimes H^2(P)$, the class
$(\alpha\otimes1)h$ is algebraic on $R_t\times P$ by the preceding
paragraph. Applying $\beta\otimes1$ proves the assertion.
\end{proof}

\begin{theorem}[A complete cyclic-five-cover family in N2]\label{p:N2cyclic}
Let $B_t=J(C_{5,t})$ with $\End^0(B_t)=\Q(\zeta_5)$, and let
$S=J(v^2=z^5-1)$. Then the Hodge conjecture holds for $B_t\times S$.
Its Hodge dimensions are $(1,4,12,14,12,4,1)$, including four
exceptional degree-four classes. The endomorphism hypothesis holds
for very general $t$, so this is a nonconstant family of sixfolds
with a simple fourfold and a simple CM surface.
\end{theorem}
\begin{proof}
The surface has CM by $K=\Q(\zeta_5)$ with type $\{1,2\}$ and
$\Hg(S)=U_K$. Put $\gamma(\zeta_5)=\zeta_5^2$ and
$\nu(s)=s\gamma(s)$ on $U_K$. This is an isogeny of rational
tori. The actual joint group is
\[
 \{(g,s)\in U_K(H^1(B_t),h)\times U_K:
                  \det_K g=\nu(s)\}.
\]
Indeed the determinant Hodge cocharacter on $B_t$ is $(2,0)$,
whereas on $S$ it is $(1,1)$; $\nu$ sends the latter to the former.
The displayed rational group is connected, has the same dimension
as $\Hg(B_t)$, and contains the Hodge circle. Projection of the
actual group onto $\Hg(B_t)$ is surjective, proving equality.
On embeddings $1,2$ the map $\nu$ is $(u,v)\mapsto(uv,v/u)$.
Using the conjugate of the second $B_t$ block gives exactly the
representation of Lemma~\ref{p:quartic-mixed}. The four remaining
classes are algebraic by Lemma~\ref{p:cyclic-transfer}; divisors
and these classes generate the Hodge ring.

For nonemptiness, the branch data $(1,1,1,2)$ are equivalent to
case 11 of \cite[Table 1]{MoonenCyclic}. This family fills its
one-dimensional PEL component. Its generic group is the full
unitary group just used, so its generic endomorphism algebra is $K$.
This statement does not assert that every rational polarization
form with this CM field occurs in the family.
\end{proof}

\begin{theorem}[All weighted quartic-cover families]\label{p:weighted-families}
For any branch data in Lemma~\ref{p:cyclic-transfer}, put
$B_t=B_{m,\boldsymbol a,t}$ and choose
\[
 P=\begin{cases}
 J(v^2=z^5-1),&m=5,10,\\
 E_i\times E_{\sqrt{-2}},&m=8,\\
 E_i\times E_{\sqrt{-3}},&m=12,
 \end{cases}
\]
where $E_{\sqrt{-d}}$ has CM by $\Q(\sqrt{-d})$ and $E_i$ by
$\Q(i)$. Then the Hodge conjecture holds for $B_t\times P$ for every
$t\ne0,1$. Whenever $B_t$ is simple, it holds for every power of
$B_t\times P$. Very general parameters have simple $B_t$ with
$\End^0(B_t)=K$.
\end{theorem}
\begin{proof}
The determinant cocharacter has one pure and one balanced block.
For $m=5,10$, $K=\Q(\zeta_5)$ and the CM types of $P$ are conjugate
under field automorphisms. Orient the action at the pure block;
the isogeny $s\mapsto s\gamma(s)$ of Theorem~\ref{p:N2cyclic},
with a suitable generator $\gamma$, sends its cocharacter to the
fourfold determinant cocharacter. For $m=8,12$, orient both imaginary
quadratic actions at that block and use $(u,v)\mapsto uv$ in $U_K$.
Thus for non-CM $B_t$ the joint group is precisely the connected
mixed model of Lemma~\ref{p:quartic-mixed}, by the same containment,
connectedness, dimension, and surjectivity argument used there.

At a simple CM parameter the group is $U_L$, with $L$ octic:
for biquadratic $K$ this follows from
Proposition~\ref{p:N3residual}; for $K=\Q(\zeta_5)$ use the
rank-three exclusion in the proof of
Corollary~\ref{p:cyclic-allparameters}. That exclusion uses only
$K\subset L$ and simplicity, not the branch vector or covering degree.
The joint group is the corresponding CM mixed model, whose generators
are again divisors and the same $Q$. Lemma~\ref{p:cyclic-transfer}
makes $Q$ algebraic at every parameter, and
Proposition~\ref{p:mixed-allpowers} proves the all-powers claim for
simple $B_t$ in both cases.

If $B_t$ is not simple, each simple factor has dimension at most
three. After adjoining $P$, the resulting sixfold belongs to the
already proved rows N5--N10. This proves the sixfold assertion at
these parameters; no all-powers assertion for these additional
factorizations is needed here.
Finally the cyclic-cover monodromy theorem
\cite[Proposition 4.7(i)]{MoonenCyclic} gives $\SU(1,1)$ on the
balanced primitive block. The pure block excludes a diagonal
identification of the two derived factors, and the determinant
cocharacter spans $U_K$. The generic commutant is $K$; hence the
very general primitive factor is simple with precisely that
endomorphism algebra.
\end{proof}

\begin{lemma}[Specializing a fixed Hodge subvariation over a curve]
\label{p:curve-specialization}
Let $f:\mathcal X\to C$ be a smooth projective family over a smooth
connected projective complex curve, and let $\mathcal V\subset R^{2p}f_*\Q$
be a rational polarizable subvariation everywhere of type $(p,p)$.
If every class of $\mathcal V_c$ is algebraic for very general $c$,
the same holds at every $c\in C$.
\end{lemma}
\begin{proof}
After a finite etale cover, $\mathcal V$ is constant. Indeed its
polarization, after the Tate twist, is positive definite and preserved
by monodromy; intersection with an integral cohomology lattice gives
a monodromy-stable lattice. Its orthogonal automorphism group is finite.
Choose a very general base point outside the images of all
nondominating irreducible components of all relative Chow schemes.
There are only countably many such components, and each image is
closed because the fixed-degree relative Chow schemes are projective.
Represent a rational basis of $\mathcal V$ there by rational linear
combinations of effective cycles. Every component containing one of
these effective cycles dominates $C$. By projective hyperplane sections
through the chosen point, choose a curve in that component which still
dominates $C$, and normalize it. Pulling back the universal cycle gives
a cycle family over a finite cover of $C$; its cohomology class is a
flat section. Taking a common cover for the finitely many effective
cycles in a basis, their indicated rational combinations extend to
all fibers. On this connected cover the corresponding flat classes
are the given constant basis of $\mathcal V$, since they agree at the
chosen point. Restricting at a point over any desired $c$ proves the
assertion, and descent is harmless for rational cycles over $\C$.
Only classes in the specified subvariation are specialized; this does
not assert that all new Hodge classes in a special fiber extend.
\end{proof}

\begin{theorem}[Completion of the represented one-dimensional PEL types]
\label{p:PELcompletion}
Fix the marked rational polarized PEL datum of any family in
Theorem~\ref{p:weighted-families}, including its signatures and the
matching fixed surface $P$. Every simple fourfold $B$ with this datum
satisfies the Hodge conjecture on all powers of $B\times P$.
Likewise, every simple sixfold with the marked rational PEL datum of
Theorem~\ref{p:S10cyclic18} satisfies the Hodge conjecture on all powers.
Here equality of data means a $K$-linear identification of the rational
Hermitian spaces up to compatible positive conformal scaling, in the
same period-domain component. No new rational Hermitian form is added
by this completion step.
\end{theorem}
\begin{proof}
In each case $H^1$ has $K$-rank two and exactly one balanced embedding;
the other embeddings are pure. Hence the PEL domain has dimension one.
The $K/K^+$ Hermitian form is definite at a real place of $K^+$, so
it is anisotropic over $K^+$. Its special unitary group has no nontrivial
$K^+$-rational unipotent elements: the image at that definite place is
compact, where a unipotent unitary matrix must be the identity.
The restriction-of-scalars derived group is semisimple and has no
nontrivial rational unipotents. The arithmetic compactness criterion
\cite[Theorem 3.3(b)]{MilneShimura} makes each corresponding connected
PEL curve compact. At a sufficiently fine level it is a smooth
projective curve $S$ carrying its universal abelian scheme.

After a finite cover of the original parameter curve, its primitive
Jacobian factor has level structure and gives an algebraic morphism
to $S$. The monodromy calculation used in the covering theorems shows
that this morphism is nonconstant. It extends to the smooth projective
completion $C$ of the parameter curve because $S$ is projective.
The resulting nonconstant morphism $C\to S$ is surjective. Pull back
the universal abelian scheme to $C$; this is a smooth family also at
points missing from the original covering parameter space.
Changing a lattice only changes these realizations up to isogeny,
so it suffices to treat this one fine-level component of the rational
datum.

For the quartic datum, on the family times the fixed $P$, the space
$Q$ is a rank-four rational subvariation of $R^4$, everywhere of type
$(2,2)$. This follows directly from the four determinant eigenspaces:
their Hodge types and those of the fixed CM factor do not change with
the period. It is polarizable as a subvariation of the smooth projective
family. Lemma~\ref{p:cyclic-transfer} makes $Q$ algebraic over the
original parameter curve, so Lemma~\ref{p:curve-specialization} makes
it algebraic at every point of $C$. At any simple fourfold point,
the group dichotomy and the CM calculation in
Theorem~\ref{p:weighted-families} still apply: their inputs are the
field action and signatures, not the branch equation. Divisors and
$Q$ generate its sixfold Hodge ring, and
Proposition~\ref{p:mixed-allpowers} gives all powers.

For the sextic datum use the rank-two subvariation $W_k$ of $R^6$,
everywhere of type $(3,3)$. Theorem~\ref{p:S10cyclic18} proves its
algebraicity on the covering parameter curve, and the same
specialization lemma treats its completion. If a simple fiber is CM,
Theorem~\ref{p:CMcomplete} applies. Otherwise the rank-two-block
calculation in Theorem~\ref{p:S10cyclic18}, determined by signatures
$(2,0),(1,1),(0,2)$, gives the full derived group and norm-kernel
determinant image; Proposition~\ref{p:coupled-det} then gives all powers
from the algebraic Weil plane. This proves both assertions.
The argument specializes cycles from a dense set in a curve to its
missing points; it never deforms a cycle from a proper lower-dimensional
subfamily into a higher-dimensional PEL space.
\end{proof}

\begin{example}[Two additional marked rational Hermitian types]
\label{ex:weighted-newtypes}
The branch vectors
\[
 (m,\boldsymbol a)=(10,(1,2,2,5)),\qquad(12,(1,1,4,6))
\]
give the curves
$y^{10}=x(x-1)^2(x-t)^2$ and
$y^{12}=x(x-1)(x-t)^4$. Their genera are six and seven, respectively,
by $g=1+m-\frac12\sum_i\gcd(m,a_i)$. Their primitive factors both
have dimension four. The preceding theorem proves the claimed N2
and N3 products at every parameter, and all powers when that factor
is simple. At non-CM parameters their Hodge polynomial is
$1+4t+12t^2+14t^3+12t^4+4t^5+t^6$.

These extend the rational Hermitian types in the earlier examples.
For $x=\zeta_m^{a_1},y=\zeta_m^{a_2},z=\zeta_m^{a_3}$, the geometric
skew-Hermitian intersection matrix of \cite[Proposition 7.4]{YuZheng}
has determinant
\[
 D(\boldsymbol a)=\frac{1-xyz}{(1-x)(1-y)(1-z)}.
\]
Use the same imaginary scalar to turn each compared pair of skew
forms into Hermitian forms, so this scalar cancels in their ratio.
For the first comparison identify $\zeta_5=\zeta_{10}^2$; for the
second write $\sqrt3=2\zeta_{12}-\zeta_{12}^3$. Exact calculation gives
\[
 \frac{D_{10}(1,2,2,5)}{D_5(1,1,1,2)}=2,\qquad
 \frac{D_{12}(1,1,4,6)}{D_{12}(1,1,1,9)}
       =1-\frac{\sqrt3}{3}.
\]
At the prime over $2$ of $\Q(\sqrt5)$, the extension
$\Q(\zeta_5)/\Q(\sqrt5)$ is unramified quadratic and inert:
$2$ has order four modulo five, and its residue degree in the real
subfield is two. The first ratio has valuation one. At
$\pi=\sqrt3$ in $\Q(\sqrt3)$, the extension obtained by adjoining
$i$ is unramified and inert, since $X^2+1$ is irreducible over
$\mathbf F_3$. The second ratio equals $(\pi-1)/\pi$ and has
valuation $-1$. Local norms in an unramified quadratic extension
have even valuation, so neither ratio is a relative norm.
A basis change multiplies a Hermitian determinant by a relative norm;
conformal scaling in rank two multiplies it by a square, also a norm.
Thus each new marked rational type differs from the indicated old
one. Theorem~\ref{p:PELcompletion} covers all simple members of each
represented marked rational PEL type, not only the explicitly
parametrized fibers. This is a comparison within this manuscript,
not a claim of priority or of distinctness of every branch orbit.
\end{example}

\begin{remark}[Finite branch catalogue and its precise scope]
Modulo permutations of the four branch points and multiplication of
all exponents by a unit modulo $m$, the hypotheses of
Lemma~\ref{p:cyclic-transfer} give the following representatives:
\begin{center}\small
\begin{tabular}{cl}
\toprule
$m$ & Branch representatives\\\midrule
5 & $(1,1,1,2)$\\
8 & $(1,1,1,5),\ (1,1,2,4),\ (1,2,2,3)$\\
10 & $(1,1,2,6),\ (1,1,3,5),\ (1,2,2,5),\ (1,2,3,4)$\\
12 & $(1,1,1,9),\ (1,1,2,8),\ (1,1,3,7),\ (1,1,4,6),$\\
 & $(1,2,3,6),\ (1,2,4,5),\ (1,3,4,4),$\\
 & $(1,3,10,10),\ (2,3,3,4)$\\\bottomrule
\end{tabular}
\end{center}
The enumeration ranges over $1\le a_1,a_2,a_3<m$, sets
$a_4=-a_1-a_2-a_3\pmod m$, and applies exactly the nonzero, gcd,
and multiplicity tests in the lemma. Canonicalizing by all units and
sorting gives $1,3,4,9$ orbits, from $16,112,240,496$ ordered types.
The script \texttt{verify\_weighted\_quartic.py} independently checks
unordered quadruples and reconstructs all ordered orbits, and verifies
the two determinant ratios. All seventeen branch families are covered
by Theorem~\ref{p:weighted-families}; they are not claimed to represent
seventeen distinct rational PEL types. Exhaustion of these finite
covering data does not exhaust all forms in N2 or N3.
\end{remark}

\subsection*{Determinant stabilization and the residual N2 cycle}
We use one additional external assertion, denoted (Aux):
the auxiliary-spin proposition in the section
\emph{Degree-four tensors in spin families} of \cite{OpenAIK3Products}.
For a totally real Galois field $F/\Q$ and an $F$-quadratic space
$(V',q')$ of dimension $d'\geq5$, $d'\equiv5\pmod8$, signature
$(2,d'-2)$ at every real place and Witt index at least two,
choose a positive period plane at each real place. The complete
spin Hodge structure
\[
 W'=\Res_{F/\Q}C_F^+(V')=H^1(A',\Q)
\]
then has the following specified algebraic tensor. If $z'$ is a
volume element, sum
\[
 \mathcal P_\sigma=\sum_jL_{v_jz'}\otimes L_{v_j^\vee z'}
 \quad\text{over }\sigma:F\hookrightarrow\C,
                                                        \tag{N2.1}
\]
where the bases are $q'$-dual and each operator is supported on
the $\sigma$ summand. After polarization identifies the two
endomorphism factors with four weight-one slots, the resulting
rational tensor $\mathcal P$ is the class of a cycle in
$\CH^2((A')^4)_\Q$. The assertion holds for every choice of period
planes and every isogenous realization. We use precisely this
assertion as an external hypothesis; it is stronger than the
specified K3 Kuga--Satake correspondence alone. Its proof in the
new preprint is not independently verified here.

\begin{theorem}[Residual N2 from the auxiliary-spin tensor]
\label{p:N2auxiliary}
Assume (Aux) and the CM Hodge theorem
\cite[Theorem 1.1]{OpenAICM}. For every coupled non-CM case
$A=B\times P$ of Proposition~\ref{p:N2nonCMstructure}, there is
$Y\in\CH^2(A)_\Q$ with nonzero projection to its mixed space $Q$.
Consequently the rational Hodge conjecture holds for every $A^n$,
$n\geq1$. There is no restriction on the primitive quartic field
or on the rational Hermitian form.
\end{theorem}
\begin{proof}
\emph{Quadratic spaces with complementary period planes.}
Put $V=H^1(B,\Q)$, $K=\End^0(B)$ and $F=K^+$; this $F$ is the
real quadratic subfield, not the CM endomorphism field of $P$.
Denote the pure and mixed real places of $F$ by $c,a$.
Choose a rank-one $K$-CM weight-one structure $M$ with CM type
$\{\bar\sigma_c,\sigma_a\}$, realized by a CM abelian surface $C_M$.
Set
\[
 U=V\otimes_KM,\qquad R=\bigwedge\nolimits_K^2V.
\]
Both are weight-two structures, of $F$-ranks four and two.
The first has rational Hodge numbers $(1,6,1)$; its period plane
is at $a$, whereas that of $R$ is at $c$.
Compatible Hermitian polarizations and an $F$-scalar with prescribed
real signs give invariant $F$-quadratic forms with signatures
\[
 \operatorname{sig}(U_c)=(0,4),\quad
 \operatorname{sig}(U_a)=(2,2),\qquad
 \operatorname{sig}(R_c)=(2,0),\quad
 \operatorname{sig}(R_a)=(0,2).
                                                        \tag{N2.2}
\]
The positive planes are exactly the indicated period planes.
Here the norm form on the $K$-line $R$ may be scaled independently
at the two real places. All these forms are invariant under
$G=\Hg(B\times C_M)$.

Let $N_7=\langle-1,\ldots,-1\rangle_F$ have rank seven and trivial
$G$-action, and put $V'=U\perp R\perp N_7$.
Its $F$-dimension is thirteen and both signatures are $(2,11)$.
At finite places a quadratic form of dimension at least five is
isotropic. Hasse--Minkowski therefore splits one hyperbolic plane
over $F$; its complement has dimension eleven and signatures
$(1,10)$, so a second application splits another.
Thus the Witt index is at least two. Since $F/\Q$ is Galois,
all hypotheses of (Aux) hold.

\emph{A rational spin lift with the original central character.}
Fix a complex embedding of $F$. Let $V_\sigma^+$ be the
two-dimensional block at one extension to $K$, put
$d=\det V_\sigma^+$, and let $m$ denote the corresponding
character of $M$. A maximal isotropic subspace of
$(U\perp R)_\sigma$ is
\[
 L=(V_\sigma^+\otimes m)\oplus\C_d,\qquad
                  \det L=d^2m^2.
\]
The inverse image of $\GL(L)$ in $\operatorname{Spin}_6$ is
$\{(g,t):t^2=\det g\}$. Taking the actual character $t=dm$
gives a lift of $G_{\overline\Q}$ and its Fock module
$S=(dm)^{-1}\bigwedge^\bullet L$. On a maximal torus, write
$b_1,b_2$ for the weights of $V_\sigma^+$, so $d=b_1b_2$.
The complete character calculation is
\[
\begin{array}{c|c}
\text{Fock states}&\text{weights}\\ \hline
1&(dm)^{-1}\\
e_1,e_2&b_2^{-1},b_1^{-1}\\
e_3&m^{-1}\\
e_{12}&m\\
e_{13},e_{23}&b_1,b_2\\
e_{123}&dm .
\end{array}                                             \tag{N2.3}
\]
The two-dimensional blocks are representations of the full group,
not just identities of torus weights:
$V_\sigma^+\otimes d^{-1}\simeq(V_\sigma^+)^\vee$
and the $e_{13},e_{23}$ block is $V_\sigma^+$ itself.
In particular the center is correct.

The $G$-action on $V'$ is rational. Two lifts through the finite
central kernel of
$\Res_{F/\Q}\operatorname{Spin}(V')\to
 \Res_{F/\Q}\SO(V')$
agree, since $G$ is geometrically connected. The lifts just
constructed at the two embeddings, extended trivially across
$N_7$, therefore descend uniquely to $\Q$.
Their Hodge circle is the spin lift of the chosen period planes.
Explicitly, at $c$ the exponents of $d,m$ are $2,-1$, and at $a$
they are $0,1$; every weight in (N2.3) has exponent $1$ or $-1$.
Thus these are the actual weight-one constituents of $W'$.

Over $\C$, restriction of the Spin$_{13}$ spinor to
Spin$_6\times$Spin$_7$ is the full eight-dimensional Spin$_6$
spinor tensored with the eight-dimensional Spin$_7$ spinor.
The full even Clifford regular module contains this with nonzero
multiplicity. Formula (N2.3) at both $F$-embeddings consequently
contains all of $V_\C$. Base change for rational $G$-Hom spaces
and semisimplicity show that $V$ is a rational summand of $W'$.
The subspace $(W')^{G^{\rm der}}$ is a rational weight-one Hodge
summand with torus Hodge group; denote its CM abelian variety by
$D$. Thus $B$ and $D$ are isogeny factors of $A'$. The factor
maps and polarization adjoints are algebraic abelian homomorphisms.
No individual complex half-spinor is asserted to have a rational
model of the same dimension.

\emph{The nonzero mixed block.}
In (N2.3), use $e_1,e_2$ for the isotropic $U$ directions and
$e_3$ for the $R$ direction. Up to a nonzero normalization, the
$U$-part of (N2.1) acts on the Spin$_6$ factor as
\[
 C_U=\sum_{i=1}^2(\varepsilon_i\otimes\iota_i+
                         \iota_i\otimes\varepsilon_i),
\]
where $\varepsilon_i$ is exterior multiplication and $\iota_i$
contraction. Its CM-output projection is
\[
 C_U(e_{13}\otimes e_{23})
       =e_3\otimes e_{123}-e_{123}\otimes e_3,
 \qquad
 C_U(e_{13}\wedge e_{23})=2e_3\wedge e_{123}\ne0.
                                                        \tag{N2.4}
\]
The $R$ operators have zero projection to these CM states.
Operators from $N_7$ preserve the $B$/CM decomposition of the
Spin$_6$ factor, through its parity operator, and also give zero.
In odd dimension the volume element acts by a nonzero scalar on
the spinor, so $L_{vz'}$ has the same nonzero block as Clifford
multiplication by $v$. Multiplicity factors can be selected by
weight-one intertwiners. Hence no other summand cancels (N2.4).

By (Aux), rational weight-one factor maps applied to its cycle,
followed by the two diagonals, yield a rational cycle
$Z\in\CH^2(B\times D)_\Q$ whose component
\[
       0\ne z\in T_B\otimes H^2(D,\Q),\qquad
       T_B=\bigwedge\nolimits_K^2V,                     \tag{N2.5}
\]
is nonzero. To justify rationality of this assertion, the complex
selectors in (N2.4) lie in the complexifications of the rational
intertwiner spaces. If every tuple of rational selectors had
zero determinant projection after the diagonals, multilinearity
would force every complex tuple to have zero projection.
Formula (N2.4) contradicts this. Polarization converts the input
determinant line to its dual line, which is still a summand of
$T_B$; the nonzero projection is unaffected.

\emph{Transport to the given CM surface.}
The four determinant characters of $T_B$ form a single Galois
orbit and are distinct for the full unitary group. Thus $T_B$
is a simple rational Hodge structure. The tensor (N2.5) defines
an injection $T_B^\vee(-2)\to H^2(D,\Q)$, with image $C$.
The given coupled space $Q\ne0$ similarly supplies an isomorphic
summand $C_P\subset H^2(P,\Q)$.
Choose a Hodge projection onto $C$ and an isomorphism $C\to C_P$;
their composite is a rational Hodge map
$\phi:H^2(D,\Q)\to H^2(P,\Q)$.
Writing $d_D=\dim D$, its Poincare-dual class lies in
\[
 H^{2d_D-2}(D,\Q)\otimes H^2(P,\Q).
 \tag{N2.6}
\]
Since $D\times P$ is CM, the external CM theorem makes this
precise rational Hodge class algebraic:
choose $\Gamma\in\CH^{d_D}(D\times P)_\Q$ representing it.
Then
\[
 Y=(\operatorname{pr}_{BP})_*
   \bigl(\operatorname{pr}_{BD}^*Z\cdot
         \operatorname{pr}_{DP}^*\Gamma\bigr)
       \in\CH^2(B\times P)_\Q.
\]
Its codimension is $2+d_D-d_D=2$ and its mixed component is
$(1\otimes\phi)z\ne0$. The single-cycle criterion of
Proposition~\ref{p:N2nonCMstructure} and
Proposition~\ref{p:mixed-allpowers} finish the proof.
\end{proof}
\begin{example}
Choose $(B_1,\ldots,B_{2})$ jointly very generally in
$\prod_i\mathcal A_{d_i}(N)$ with $d_i=(4,2)$,
as constructed in Proposition~\ref{p:products}. Then each factor is
simple with $\End^0(B_i)=\Q$, and the factors are pairwise nonisogenous.
Their product is an actual sixfold of this partition, with
$\End^0(A)=\Q^{2}$ and the proved joint group above.
The Hodge-space dimensions in codimensions $0$ through $6$ are
$(1,2,3,3,3,2,1)$.
An explicit monomial basis over $\Q$ on $A$ is
$\theta_1^{a_1}\cdots\theta_{2}^{a_{2}}$ with
$0\leq a_i\leq d_i$; every such monomial is an intersection of divisors.
On powers add Poincare divisors between the copies of each individual
$B_i$. There are no cross-factor divisors or higher mixed invariants in
this verified jointly generic example.
\end{example}
\section{N3: The \texorpdfstring{$4+1+1$}{4+1+1} case (\provedstatus)}\label{sec:N3}

\status{The independent cases, the imaginary-quadratic endomorphism
stratum, and every repeated elliptic pair are proved. The new external
CM theorem also settles all remaining cases with a CM fourfold.
Theorem~\ref{p:N3KS} covers the residual non-CM fourfold with
biquadratic CM endomorphism field and two distinct matching CM
elliptic fields, using the specified complete KS correspondence.
It supplies the required nonzero mixed projection for every rational
Hermitian type and proves all powers of each residual model.}

\begin{corollary}[All remaining CM-fourfold cases in N3]\label{p:N3CM-new}
Assume the CM input in Theorem~\ref{p:openai-inputs}.
If $B,E_1,E_2$ are CM, then the Hodge conjecture holds on all
powers of $B\times E_1\times E_2$, whether or not the elliptic
fields coincide or the joint torus splits as a product.
Together with Proposition~\ref{p:N3repeated}, this settles every
sixfold in N3 whose fourfold factor is CM.
\end{corollary}
\begin{proof}
The all-CM product and each of its powers are covered by
Lemma~\ref{p:CMproducts-new}. Proposition~\ref{p:N3repeated}
leaves only two nonisogenous CM elliptic factors as a possible
residual case. Hence, when $B$ is CM, the preceding argument
settles every case that proposition leaves open.
\end{proof}
\begin{example}[Distinct elliptic fields with a CM fourfold]
Take the simple CM fourfold $B$ of Example~\ref{ex:N2CM-new},
and elliptic curves $E_1,E_2$ with CM by
$k_1=\Q(\sqrt{-3})$ and $k_2=\Q(\sqrt{-15})$, respectively.
Both fields are contained in $\Q(\zeta_{15})$ and are distinct.
For the type $\{1,2,4,7\}$, the multiplicities for the four
embeddings of $k_1k_2$, indexed by signs, are
$(n_{++},n_{+-},n_{-+},n_{--})=(2,1,1,0)$ after orientation.
Thus the fourfold has the pure--mixed quartic pattern rather than
two balanced quartic signatures. Nevertheless all powers of
$B\times E_1\times E_2$ satisfy the Hodge conjecture by
Corollary~\ref{p:N3CM-new}, including their mixed classes.
\end{example}

\begin{proposition}
Let $A\sim B_1\times\cdots\times B_{3}$ with dimensions
$(4,1,1)$. Assume the actual joint Hodge group is the
independent product of the full symplectic groups on these factors.
Then all Hodge classes on all powers of $A$ are divisor-generated,
and $P_A(t)=(1+t+\cdots+t^{4})(1+t)(1+t)$.
\end{proposition}
\begin{proof}
Use Lemma~\ref{p:isogeny} to replace the isogenous variety by the product.
Under the stated group equality,
\[
 \left(\bigwedge^\bullet\bigoplus_i H^1(B_i)^{\oplus n}\right)^{\Hg(A)}
 =\bigotimes_i\left(\bigwedge^\bullet H^1(B_i)^{\oplus n}\right)^{\Sp(H^1(B_i))}.
\]
Each factor is generated by degree-two contractions by Lemma~\ref{p:fft}.
These are divisor classes on the product, so every invariant is algebraic.
On one copy the generators are the polarization powers through the
dimension of each factor, giving the displayed polynomial.
\end{proof}
\begin{proposition}[Quadratic-center and independent strata]\label{p:N3strata}
Let $A=B\times E_1\times E_2$, with $B$ a simple fourfold.
The Hodge conjecture holds if no CM field of an $E_i$ embeds in the
center of $\End^0(B)$. It also holds whenever $\End^0(B)$ itself is
an imaginary quadratic field $k$, including both $E_i$ with CM by $k$.
\end{proposition}
\begin{proof}
Collect repeated elliptic isogeny classes. In the first case apply
Theorem~\ref{p:lowdim}(2) successively to the distinct classes: their
Hodge groups separate from $B$, and distinct elliptic classes separate
from one another. The fourfold theorem and divisor generation on
elliptic products then prove algebraicity.

Suppose $\End^0(B)=k$. Over $\C$ write $H^1(B)=W\oplus W^*$ with
$\dim W=4$. The two constituents are irreducible and nonisomorphic.
The faithful irreducible minuscule representations of dimension four
are the standard $\SL_4$, standard $\Sp_4$, and $2\otimes2$ of
$\SL_2\times\SL_2$. In the latter two cases the representation is
self-dual. For a nontrivial derived Hodge circle its weights are
symmetric; translating them to the two weights $1,-1$ forces central
weight zero and multiplicities $(2,2)$. A trivial derived circle would
force the entire group to be a torus and is impossible here.
Thus in these self-dual cases the determinant circle is trivial; all
its conjugates are trivial, the connected center vanishes, and
$W\cong W^*$ contradicts the commutant. Hence the derived group is
$\SL_4$. Pure signatures are excluded by simplicity.
For signature $(2,2)$ the group is $\SU_4$; for $(3,1)$ or $(1,3)$ it
is the full unitary group, with determinant-circle exponent $\pm2$.

In the balanced case $\SU_4$ has no elliptic quotient, so all elliptic
factors separate. In the unbalanced case only elliptic curves with CM
by $k$ can interact. Let their multiplicity be $m\leq2$, and orient
their $k$-action opposite to the larger signature on $B$. The joint
group is the connected kernel of $\det(v)u^2=1$ in $U(W)\times U_k$:
it contains the Hodge circle and has the same dimension as $U(W)$,
onto which the actual group projects. The remaining elliptic factors
separate. On the single $B$ factor, $\SL_4$-invariants are pairings and
the degree-four volumes, whose central weights are $\det^{\pm1}$.
To cancel a volume requires elliptic weight $\mp2$. This is impossible
for $m<2$ and, for $m=2$, uses exactly two same-sign elliptic coordinates.
The only new classes are therefore the two conjugate lines of
$W_k(B\times E^2)$ in degree six. This product has signature $(3,3)$
and a rank-one $k$-elliptic summand, so Lemma~\ref{p:oddscale} makes
its Weil plane algebraic. All other invariants are divisor products.
\end{proof}
\begin{proposition}[All repeated elliptic factors]\label{p:N3repeated}
For every simple fourfold $B$ and elliptic curve $E$, the Hodge
conjecture holds for $B\times E^2$, without a restriction on
$\End^0(B)$. More generally, for $B\times E_1\times E_2$ the
only cases not covered here have $E_1,E_2$ nonisogenous CM curves
with both CM fields embedding in the center of $\End^0(B)$.
\end{proposition}
\begin{proof}
Moonen--Zarhin's fivefold classification
\cite[Theorem (0.2) and \S5.11]{MZ} shows that $\Hg(B\times E)$
is the independent product unless $E$ has CM by $k$ acting on $B$
with signature $(3,1)$, up to interchange. Repeating $E$ does not
change this group. In the independent case the Hodge ring on
$B\times E^2$ is a tensor product, and its classes are algebraic
by the fourfold theorem and elliptic divisor generation.

In the exceptional case the same classification gives a CM field
$F=\End^0(B)$ containing $k$, the full unitary Lefschetz group
$G=U_F(H^1(B),h)$, and the connected joint group
\[
 H=\{(g,u)\in G\times U_k:\det_k(g)u^2=1\}.
\]
Here the action on $E$ is opposite to the larger $k$-signature of $B$.
Write $[F:\Q]=2e$ and $q=4/e$, so $e=1,2,4$ and
$G_\C=\prod_{j=1}^e\GL(W_j)$ with $\dim W_j=q$.
Orient the embeddings over $k$ consistently. The kernel of
$\det_k$ has derived factors $\SL(W_j)$ and determinant relation
$\prod_j d_j=1$. Its exterior invariants on the single $B$ factor
are pairings and the two simultaneous volumes, of degree $eq=4$.
Indeed central neutrality forces all signed determinant exponents
to agree; on one copy they are $0,1$, or $-1$.

The volumes acquire elliptic weights $\mp2$. In the exterior
algebra of two copies of $H^1(E)$, weight $2$ or $-2$ consists of
the unique product of the two same-sign coordinates. Hence the only
additional invariants are the two lines of $W_k(B\times E^2)$ in
degree six; the other invariants are divisor products. The rank-one
elliptic summand allows Lemma~\ref{p:oddscale}, proving their
algebraicity.

Finally, if the two elliptic curves are nonisogenous and one has
no CM field embedding in the center of $\End^0(B)$ (including
a non-CM curve), Proposition (3.8) of \cite{MZ} separates that curve
from $B$ times the other curve. In the CM case its field also differs
from that of the other elliptic curve. The remaining fivefold has
algebraic Hodge classes by Theorem (0.2) of \cite{MZ} and the
fourfold theorem. Their exterior products prove the claim.
\end{proof}
\Needspace{12\baselineskip}
\begin{proposition}[The full residual structure in N3]\label{p:N3residual}
Let $A=B\times E_1\times E_2$ be an N3 case not already covered
by Proposition~\ref{p:N3repeated}. Write $k_i=\End^0(E_i)$.
Then $k_1\ne k_2$ are imaginary quadratic fields in the center
of $\End^0(B)$, and $K=k_1k_2$ acts on $H^1(B)$ with signatures
$(2,0),(1,1)$, after orienting the two actions. Put
\[
 T=\bigwedge\nolimits_K^2 H^1(B,\Q),\qquad
 Q=\bigl(T\otimes H^1(E_1,\Q)\otimes H^1(E_2,\Q)\bigr)
       \cap H^{2,2}(A).
\]
Here $T$ has the same summand interpretation as in
Lemma~\ref{p:cyclic-transfer}. There are exactly two alternatives:
\begin{enumerate}
\item $\End^0(B)=K$, $\Hg(B)_\C=\GL_2^2$, and
$P_A(t)=1+4t+12t^2+14t^3+12t^4+4t^5+t^6$;
\item $B$ is simple CM with field $L\supset K$ of degree eight,
$\Hg(B)=U_L$, and
\[
 P_A(t)=(1+t)^6+4t^2(1+t)^2
 =1+6t+19t^2+28t^3+19t^4+6t^5+t^6.
\]
\end{enumerate}
In both cases the actual joint group is the connected group
\[
 H=\{(g,u,v)\in\Hg(B)\times U_{k_1}\times U_{k_2}:
                         \det_K g=uv\text{ in }U_K\}.
\]
The rational space $Q$ has dimension four and, together with
divisors, generates the entire Hodge ring of $A$.
\end{proposition}
\begin{proof}
The field and embedding restrictions follow from
Proposition~\ref{p:N3repeated}. Since the center contains $K$,
the dimension formula for the simple algebra $\End^0(B)$ forces
it to be a CM field of degree four or eight.
Index the four embeddings of $K$ by their signs on $k_1,k_2$.
Their $H^{1,0}$ multiplicities have the form
\[
 (n_{++},n_{+-},n_{-+},n_{--})=(a,b,2-b,2-a),
 \qquad 0\leq a,b\leq2.
\]
The signatures of $k_1,k_2$ are given by $a+b$ and $a+2-b$.
Neither can be $0$ or $4$: a pure imaginary-quadratic action
would make $B$ isogenous to a fourth power of an elliptic curve.
If $a=b=1$, both $K$ determinants are trivial on the Hodge circle
and hence on its rational Hodge group. That group lies in
$\SL_2^2$, so each standard block is isomorphic to its dual.
The resulting off-diagonal intertwiners in the complex commutant
do not commute with $K\otimes\C$, contradicting centrality of
$K$ in $\End^0(B)$. Enumerating the remaining pairs $(a,b)$,
and conjugating the two actions if necessary, gives
$(a,b)=(2,1)$. In particular both $k_i$ have signature $(3,1)$.

If the endomorphism field is $K$, the derived-factor and determinant
argument in Lemma~\ref{p:cyclic-transfer} gives the full unitary
group. If it is $L$ of degree eight, choose the four embeddings
of $L$ over the positive embedding of $k_1$. The Hodge cocharacter
has three entries $1$ and one entry $-1$. Galois elements fixing
$k_1$ act transitively on these four embeddings. Their orbit
contains all four such vectors, the rows of $J_4-2I_4$, whose
determinant is $-16$. Thus the Hodge group has rank four and is
the full torus $U_L$.

Orient the elliptic actions at the pure $K$ embedding. The map
$\nu:U_{k_1}\times U_{k_2}\to U_K$, $(u,v)\mapsto uv$,
is an isogeny with complex coordinates $(uv,u/v)$. It sends the
elliptic Hodge cocharacter $(1,1)$ to $(2,0)$, the determinant
cocharacter of $B$. Hence the displayed rational group contains
the joint Hodge circle. The map $\det_K$ has connected kernel:
over $\C$ it is either the two ordinary determinants on
$\GL_2^2$, or $(a_1,a_2,b_1,b_2)\mapsto(a_1a_2,b_1b_2)$
on $\mathbf G_m^4$. Thus $H$ is connected. Its dimension equals
$\dim\Hg(B)$, and the actual joint group projects onto
$\Hg(B)$, proving equality. Using only the two squared norm
relations would leave an extraneous connected component; the
primitive relation $\det_K g=uv$ specifies the actual group.

The non-CM invariant calculation is Lemma~\ref{p:quartic-mixed}.
In the CM case, each two-dimensional block splits into two torus
characters. Its determinant-kernel invariants are the two
individual pairings, their product, and the two volumes. Thus its
weight-zero polynomial is $(1+t)^2$, while its volume weights
and degrees are unchanged. The same neutrality equations again
allow only one active volume block. The divisor polynomial is
$(1+t)^6$ and the additional contribution is
$4t^2(1+t)^2$. These four degree-four lines descend to $Q$;
multiplication by divisors generates all remaining degrees.
\end{proof}

\Needspace{7\baselineskip}
\begin{corollary}[A single-cycle criterion for all residual N3 cases]
\label{p:N3onecycle}
In Proposition~\ref{p:N3residual}, the Hodge conjecture holds for
$A$ if and only if there is a cycle $Z\in\CH^2(A)_\Q$ whose
cohomology class has nonzero projection to $Q$.
\end{corollary}
\begin{proof}
The projector onto $T$ is a rational polynomial in a general
$K$ endomorphism on $H^2(B)$, as in
Lemma~\ref{p:cyclic-transfer}; it is algebraic. Tensoring it with
the degree-one projectors of the two elliptic curves gives an
algebraic projector which sends a Hodge class to its $Q$ component.
The field $k_1\otimes_\Q k_2=K$ acts on
$H^1(E_1)\otimes H^1(E_2)$ by algebraic correspondences.
Consequently $Q$ is a $K$-module. Its rational dimension is four,
so it is one-dimensional over $K$; every nonzero vector generates
it. A cycle with nonzero projection therefore makes all of $Q$
algebraic. Proposition~\ref{p:N3residual} proves sufficiency, and
necessity is immediate. This is a criterion for each fixed $A$;
it does not assert existence or deformability of the required cycle.
\end{proof}

\begin{theorem}[Two coupled cyclic-cover families in N3]\label{p:N3cyclic}
For $m=8$ or $12$, let $B_{m,t}$ be as in
Lemma~\ref{p:cyclic-transfer}, with
$\End^0(B_{m,t})=\Q(\zeta_m)$. Let $E_1$ have CM by $k_1=\Q(i)$,
and let $E_2$ have CM by
\[
 k_2=\Q(\sqrt{-2})\quad(m=8),\qquad
 k_2=\Q(\sqrt{-3})\quad(m=12).
\]
Then the Hodge conjecture holds for $B_{m,t}\times E_1\times E_2$,
with Hodge dimensions $(1,4,12,14,12,4,1)$. The endomorphism
hypothesis holds for very general $t$ in each family.
\end{theorem}
\begin{proof}
Here $K=k_1k_2$. Orient both elliptic actions positively at the
pure-signature embedding of $B_{m,t}$. The rational torus isogeny
\[
 \nu:U_{k_1}\times U_{k_2}\longrightarrow U_K,
 \qquad (u,v)\longmapsto uv
\]
has complex coordinates $(uv,u/v)$ after choosing the other
embedding over the same embedding of $k_1$. The determinant Hodge
cocharacter of the fourfold is $(2,0)$, the image under $\nu$ of
the elliptic cocharacter $(1,1)$. Since the two elliptic curves
are nonisogenous, their joint group is $U_{k_1}\times U_{k_2}$.
As in Theorem~\ref{p:N2cyclic}, connectedness and dimension give
the actual joint group as $\det_K(g)=\nu(u,v)$.
Lemma~\ref{p:quartic-mixed} identifies all extra classes in
$T\otimes H^2(E_1\times E_2)$, and
Lemma~\ref{p:cyclic-transfer} algebraizes them.

Finally, the primitive eigenspace multiplicities computed in that
lemma are $(2,1,1,0)$. The cyclic-cover monodromy result
\cite[Proposition 4.7(i)]{MoonenCyclic} gives full $\SU(1,1)$
projection on the balanced block. Rational Galois conjugation gives
both complex $\SL_2$ projections. They cannot be diagonally
identified, since one real unitary factor is compact and the other
is not. The determinant cocharacter spans both central directions.
Thus the generic Hodge group is full $\GL_2^2$ after
complexification and its commutant is precisely $K$.
The primitive factor, not the whole Jacobian, is the fourfold here.
\end{proof}

\Needspace{10\baselineskip}
\begin{corollary}[All parameters in the three cyclic-cover families]
\label{p:cyclic-allparameters}
The Hodge conjecture holds for all $t\in\C\setminus\{0,1\}$
in the products of Theorems~\ref{p:N2cyclic} and
\ref{p:N3cyclic}, without the assumption
$\End^0(B_{m,t})=\Q(\zeta_m)$. If the fourfold factor remains
simple but becomes CM, the Hodge dimensions are
$(1,6,19,28,19,6,1)$ instead of $(1,4,12,14,12,4,1)$.
\end{corollary}
\begin{proof}
If $B_{m,t}$ is not simple, every simple factor has dimension at
most three. Its product with the specified surface or elliptic
pair belongs to the already proved rows N5--N10. Assume it is
simple. The dichotomy in Lemma~\ref{p:cyclic-transfer} shows
that it either has the stated quartic endomorphism field or is CM
with an octic field $L$. The former case was proved above.
For $m=8,12$, Proposition~\ref{p:N3residual} supplies the CM
generators, and the now unrestricted
Lemma~\ref{p:cyclic-transfer} algebraizes $Q$ at this parameter.

For $m=5$, write $K=\Q(\zeta_5)\subset L$. We first prove
$\Hg(B_{5,t})=U_L$. The eight distinct CM weights of a simple
fourfold force its Hodge rank to be at least three: a rank-$r$
orbit basis of signed Hodge cocharacters gives at most $2^r$
distinct weight columns. If its rank were three,
\cite[Lemma (3.7)]{MZ} would give an imaginary quadratic field
$k\subset L$ with $\Hg(B_{5,t})=\ker N_{L/k}$ and
$k$-signature $(2,2)$. Since $K$ has no imaginary quadratic
subfield, $L=Kk$ is Galois with group $C_4\times C_2$.
Let $\gamma$ generate $\operatorname{Gal}(L/k)$, and let $\tau$
fix $K$ and conjugate $k$. Complex conjugation is
$c=\gamma^2\tau$. A balanced CM type, indexed over $k$ by
$1,\gamma,\gamma^2,\gamma^3$, has two plus and two minus signs.
If the plus signs are adjacent, $\gamma^2$ negates the sign vector
and $\tau$ fixes the CM type; if they are opposite, $\gamma^2$
itself fixes the type. In either case the type is induced from a
proper subfield, contrary to simplicity. Rank three is impossible.

Thus $\Hg(B_{5,t})=U_L$. Replacing $\GL_2^2$ in the group of
Theorem~\ref{p:N2cyclic} by its maximal torus $U_L$ gives the
actual joint group: $\det_K=N_{L/K}$ has connected kernel, and
the same connectedness, Hodge-circle and dimension argument
applies. Its weights are those in the CM calculation of
Proposition~\ref{p:N3residual}. Divisors and the same four-dimensional
$Q\subset T\otimes H^2(S)$ generate its Hodge ring, and
Lemma~\ref{p:cyclic-transfer} proves their algebraicity. This
argument treats each special fibre directly; it does not infer
algebraicity by an unsupported deformation or specialization step.
\end{proof}
\subsection*{The complete Kuga--Satake construction for residual N3}
We use the following explicit external input. For a projective K3 surface
$X$, let $C_X$ be the \emph{complete} even-Clifford Kuga--Satake abelian
variety of $T(X)$. The specified standard map
\[
 \kappa_X:T(X)\longrightarrow H^1(C_X,\Q)^{\otimes2}
\]
is induced by a rational cycle in $\CH^2(X\times C_X^2)_\Q$.
This is the assertion of \cite[Theorem 1.1]{OpenAIKS}, used as an
external hypothesis, not independently proved here. We only need it for
the Picard-number-fourteen surfaces constructed below.

\begin{lemma}[Connected Spin descent]\label{p:N3spin-descent}
Let $G$ be a geometrically connected algebraic group over a field $F$
of characteristic zero and $r:G\to\SO(W,q)$ an $F$-morphism. If
$r_{\overline F}$ lifts to $\operatorname{Spin}(W,q)$, then that lift is
unique and descends to $F$.
\end{lemma}
\begin{proof}
Two lifts differ by a morphism from the connected group $G_{\overline F}$
to $\mu_2$, necessarily the identity character. Every Galois conjugate
of a lift is therefore the same lift, proving descent.
\end{proof}

\begin{theorem}[Residual N3 under the specified KS correspondence]
\label{p:N3KS}
Assume the external Kuga--Satake assertion above. Let
$A=B\times E_1\times E_2$ be in the non-CM alternative of
Proposition~\ref{p:N3residual}: $B$ is simple,
$\End^0(B)=K=k_1k_2$ for distinct imaginary quadratic fields,
$E_i$ has CM by $k_i$, and the $K$-multiplicities are $(2,1,1,0)$.
Then there is a cycle $Z\in\CH^2(A)_\Q$ with nonzero projection to
the four-dimensional space $Q$ of that proposition. In particular
the rational Hodge conjecture holds for every $A^n$, $n\geq1$.
No Hermitian discriminant or covering-family restriction is imposed.
\end{theorem}
\begin{proof}
\emph{A geometric K3-type structure.}
Write $k_1=\Q(a)$, $a^2=-d$ with $d>0$, and
$V=H^1(B,\Q)$. The rational Hodge projector
\[
 p=\tfrac12\bigl(1+(a_B\otimes a_{E_1})/d\bigr)
 \quad\hbox{on }V\otimes H^1(E_1,\Q)
\]
has an eight-dimensional image $W$. Orient the elliptic characters
$u,v$ as in the proof of Proposition~\ref{p:N3residual}, and write
$V_\C=V_1\oplus V_1^\vee\oplus V_2\oplus V_2^\vee$,
where $\det V_1=uv$ and $\det V_2=u/v$. Here $\vee$ denotes the
dual \emph{Hodge-group representation}; the Hodge weights are assigned
separately. Then
\[
 W_\C=V_1u^{-1}\oplus V_2u^{-1}
       \oplus V_1^\vee u\oplus V_2^\vee u,
 \qquad h(W)=(1,6,1).
\]
The derived group $\SL_2^2$ has no fixed vector in $W$, so $W$ has
no rational $(1,1)$ class. The restricted geometric polarization gives
a rational quadratic form $q$ of signature $(2,6)$, using the K3
cup-product sign. Scale it to an even integral lattice $\Lambda_W$.
Its discriminant group has length at most eight. The inequalities
$2<3$, $6<19$ and $\ell(A_{\Lambda_W})+2\leq10<22-8$
give a primitive embedding in the K3 lattice by Nikulin's criterion
\cite[Theorem 14.1.12]{HuybrechtsK3}. Period-map surjectivity
\cite[Theorem 6.3.1]{HuybrechtsK3} then gives a projective K3 surface
$X$ with $\rho(X)=14$ and a rational polarized Hodge isometry
\[
                 T(X)\simeq(W,q).                         \tag{N3.1}
\]
The isometry (N3.1) is not assumed algebraic.

\emph{The rational Spin lift and KS factors.}
Over $\C$, set $L_i=V_i u^{-1}$ and $L=L_1\oplus L_2$.
The quadratic pairing identifies $W_\C=L\oplus L^\vee$, and
\[
 \det L=(\det V_1)(\det V_2)u^{-4}=u^{-2}.
\]
For the standard inclusion $\GL(L)\subset\SO(L\oplus L^\vee)$,
its Spin inverse image is $\{(g,t):t^2=\det g\}$, acting on
$\bigwedge^\bullet L$ by $t^{-1}\bigwedge^\bullet g$.
Taking $t=u^{-1}$ gives a lift of the actual connected joint
Hodge group of Proposition~\ref{p:N3residual} over $\overline\Q$.
The action on $(W,q)$ is rational, so Lemma~\ref{p:N3spin-descent}
gives a rational group lift $H\to\operatorname{Spin}(W,q)$.
Its Hodge circle is the KS lift of the period circle; the induced
action on the complete even Clifford module
$U=H^1(C_X,\Q)=\operatorname{Cl}^+(W,q)$ is its actual weight-one
Hodge structure.

The two complex half-spin representations are
\[
 S^+=u\bigwedge^{\rm even}L
  =u\oplus v\oplus v^{-1}\oplus u^{-1}
      \oplus(V_1\otimes V_2\otimes u^{-1}),
\qquad
 S^-=u\bigwedge^{\rm odd}L
  =V_1\oplus V_2\oplus V_1^\vee\oplus V_2^\vee.
\]
The equalities are of $H_\C$-modules; they do not posit rational
models of individual half-spinors. The \emph{complete} regular module
satisfies
\[
 U_\C\simeq(S^+)^{\oplus8}\oplus(S^-)^{\oplus8}.
\]
Since $H$ is reductive, scalar extension of rational $H$-Hom
spaces to $\C$ identifies them with their complex counterparts.
Thus $V,H^1(E_1),H^1(E_2)$ are rational Hodge summands of $U$;
indeed one obtains an isogeny
\[
               C_X\sim B^8\times E_1^8\times E_2^8\times D,
               \qquad\dim D=16.                       \tag{N3.2}
\]
Only the first three types of factors are needed. Their rational
weight-one inclusions, projections and polarization adjoints are
induced by rational abelian homomorphisms.

\emph{Algebraic extraction of the needed matrix coefficients.}
Choose $t_0\in W_\Q$ with $q(t_0)\ne0$. In the endomorphism
realization, the specified KS map is
\[
 \kappa_0:W(1)\longrightarrow\End(U),\qquad
 \kappa_0(w)(x)=wxt_0;
                                                        \tag{N3.3}
\]
see \cite[Chapter 4, Proposition 2.6 and Remark 2.7]{HuybrechtsK3}.
The polarization of $U$ converts (N3.3) into the tensor map in the
external KS hypothesis. Right multiplication $R_r:x\mapsto xr$ by
$r\in\operatorname{Cl}^+(W,q)$ commutes with the left KS Hodge
action and is therefore an algebraic rational endomorphism of $C_X$.
Postcomposing (N3.3) with $R_r$ replaces $t_0$ by $t_0r$.
Since $t_0$ is invertible, $t_0\operatorname{Cl}^+=\operatorname{Cl}^-$:
every odd right matrix coefficient can be selected by algebraic
right multiplication.

Over $\C$ write
\[
 \operatorname{Cl}^+
   =\End(S^+)\oplus\End(S^-),\qquad
 \operatorname{Cl}^-
   =\Hom(S^-,S^+)\oplus\Hom(S^+,S^-).
\]
The right matrix index is a trivial $H$-multiplicity space. If
$e_1,e_2$ is a basis of $L_1$, the vectors $1,e_1\wedge e_2\in S^+$
have characters $u,v$, respectively. In the regular summand
$S^+\otimes(S^+)^\vee$, choose a multiplicity covector
$\lambda\in(S^+)^\vee$. In $S^-\otimes(S^-)^\vee$, choose a
right coefficient $\tau\in S^-$ and an odd matrix
$b:S^-\to S^+$ with $\lambda(b\tau)=1$; such $b$ is among the
complex linear combinations of $t_0r$ above. Evaluate the output
at $\tau$ and project its left index to $V_1\subset S^-$.
The resulting complex matrix components
of (N3.3) are
\[
 \alpha_1(w)=\operatorname{pr}_{V_1}c(w)1,
 \qquad
 \alpha_2(w)=\operatorname{pr}_{V_1}c(w)(e_1\wedge e_2),
                                                        \tag{N3.4}
\]
viewed respectively in
$H^1(B)\otimes H^1(E_1)$ and
$H^1(B)\otimes H^1(E_2)$ by elliptic polarization and the
corresponding Tate twist. The weights on both sides are two.
These component selectors are complex $H$-equivariant maps on
weight-one factors of $U$. Rational $H$-Hom spaces base-change to
the complex Hom spaces; hence (N3.4) belongs to the complex span of
maps induced by finitely many rational algebraic correspondences
\[
 \mathcal G_i\subset\CH^2(X\times B\times E_i)_\Q.
                                                        \tag{N3.5}
\]
Concretely, apply the rational factor maps in (N3.2), their
polarization adjoints, and right multiplication to the two abelian
slots of the KS cycle. No algebraicity of the abstract embedding
$W\subset H^2(B\times E_1)$ is used.

\emph{Nonzero contraction and rational cycle.}
Let $f_1,f_2\in L_1^\vee$ be dual to $e_1,e_2$. In the exterior
Clifford model $c(e_i)=e_i\wedge(-)$ and $c(f_i)=\iota_i$.
If $b_1,b_2$ denote the corresponding basis of $V_1$, then
\[
 \begin{array}{c|cc}&e_1&e_2\\\hline
 \alpha_1&b_1u^{-1}&b_2u^{-1}
 \end{array},\qquad
 \begin{array}{c|cc}&f_1&f_2\\\hline
 \alpha_2&b_2v^{-1}&-b_1v^{-1}
 \end{array}.
\]
The opposite entries and the $L_2$ terms vanish after projection
to $V_1$. Contracting with $q^{-1}$ and pulling the two $B$ slots
back along the diagonal gives a nonzero scalar multiple of
\[
        (b_1\wedge b_2)\otimes u^{-1}\otimes v^{-1}.
                                                        \tag{N3.6}
\]
Indeed the $L_1\otimes L_1^\vee$ terms give
$b_1\otimes b_2-b_2\otimes b_1$; the normalization of $q$ changes
only the nonzero scalar. Since $\det V_1=uv$, (N3.6) is an
$H$-invariant in the target $Q_\C$.

To realize this contraction by rational cycles, choose a rational
point $x_0\in X$ and a basis $D_a$ of $\NS(X)_\Q$, with
$G_{ab}=D_a\cdot D_b$. The algebraic transcendental projector
\[
 \pi_X^{\rm tr}=\Delta_X-x_0\times X-X\times x_0
   -\sum_{a,b}(G^{-1})_{ab}D_a\times D_b
   \in\CH^2(X\times X)_\Q                          \tag{N3.7}
\]
has cohomology class equal to the inverse pairing on $T(X)$.
For $\Gamma_i\in\mathcal G_i$, form their exterior product,
intersect it with the pullback of $\pi_X^{\rm tr}$ on the two
$X$ slots, push forward along $X\times X$, and pull back by
\[
 \delta:A\longrightarrow(B\times E_1)\times(B\times E_2),
 \quad(b,e_1,e_2)\longmapsto((b,e_1),(b,e_2)).
\]
Call the resulting rational cycle $F(\Gamma_1,\Gamma_2)$.
Its codimension is $4+2-4=2$ before $\delta^*$ and remains two
after it. By (N3.6), the complexification of this bilinear
construction has a pair of inputs with nonzero $Q$-projection.
If every pair of rational correspondences in (N3.5) had zero
projection, all complex linear combinations would too, a
contradiction. Thus some
$Z=F(\Gamma_1,\Gamma_2)\in\CH^2(A)_\Q$ has nonzero
$Q$-projection. Corollary~\ref{p:N3onecycle} makes all of $Q$
algebraic; Proposition~\ref{p:mixed-allpowers} gives the assertion
for all powers.\qedhere
\end{proof}

\begin{remark}[Logical scope of Theorem~\ref{p:N3KS}]
The proof uses the \emph{specified} full even-Clifford KS map as an
external algebraicity input. It does not turn a complex half-spinor
into a rational abelian factor, assume the Hodge isometry (N3.1)
algebraic, or claim that an individual complex line in (N3.6) is a
rational cycle. The rational descent occurs in the full regular module
and the finite rational correspondence spaces (N3.5).
\end{remark}
\begin{example}
Choose $(B_1,\ldots,B_{3})$ jointly very generally in
$\prod_i\mathcal A_{d_i}(N)$ with $d_i=(4,1,1)$,
as constructed in Proposition~\ref{p:products}. Then each factor is
simple with $\End^0(B_i)=\Q$, and the factors are pairwise nonisogenous.
Their product is an actual sixfold of this partition, with
$\End^0(A)=\Q^{3}$ and the proved joint group above.
The Hodge-space dimensions in codimensions $0$ through $6$ are
$(1,3,4,4,4,3,1)$.
An explicit monomial basis over $\Q$ on $A$ is
$\theta_1^{a_1}\cdots\theta_{3}^{a_{3}}$ with
$0\leq a_i\leq d_i$; every such monomial is an intersection of divisors.
On powers add Poincare divisors between the copies of each individual
$B_i$. There are no cross-factor divisors or higher mixed invariants in
this verified jointly generic example.
\end{example}
\begin{example}[A mixed Weil plane requiring both elliptic factors]
Take a very general simple $k=\Q(i)$ fourfold with Hermitian matrix
$\operatorname{diag}(1,1,1,-1)$, and let $E:y^2=x^3-x$ have the opposite
$k$-action. The PEL argument gives $\Hg(B)=U(3,1)$ and $\End^0(B)=k$.
On $A=B\times E^2$ the total signature is $(3,3)$ and the determinant
is $-1$. Proposition~\ref{p:N3strata} yields
\[
 P_A(t)=(1+t+\cdots+t^4)(1+4t+t^2)+2t^3,
 \qquad (1,5,6,8,6,5,1).
\]
The five divisor classes consist of one on $B$ and four on $E^2$.
They generate the ordinary part; the two additional middle classes
are the algebraic Weil plane. With only one elliptic factor the weight
$2$ of the fourfold volume could not be cancelled.
\end{example}
\section{N4: The \texorpdfstring{$3+3$}{3+3} case (\provedstatus)}\label{sec:N4}

\status{Proved for all simple threefold pairs, including nonisogenous unitary
and CM factors. The actual Hodge-circle calculation below rules out
extra central couplings; the remaining Weil plane is algebraic.}
\begin{proposition}
Let $A\sim B_1\times\cdots\times B_{2}$ with dimensions
$(3,3)$. Assume the actual joint Hodge group is the
independent product of the full symplectic groups on these factors.
Then all Hodge classes on all powers of $A$ are divisor-generated,
and $P_A(t)=(1+t+\cdots+t^{3})(1+t+\cdots+t^{3})$.
\end{proposition}
\begin{proof}
Use Lemma~\ref{p:isogeny} to replace the isogenous variety by the product.
Under the stated group equality,
\[
 \left(\bigwedge^\bullet\bigoplus_i H^1(B_i)^{\oplus n}\right)^{\Hg(A)}
 =\bigotimes_i\left(\bigwedge^\bullet H^1(B_i)^{\oplus n}\right)^{\Sp(H^1(B_i))}.
\]
Each factor is generated by degree-two contractions by Lemma~\ref{p:fft}.
These are divisor classes on the product, so every invariant is algebraic.
On one copy the generators are the polarization powers through the
dimension of each factor, giving the displayed polynomial.
\end{proof}
\begin{lemma}[Two sextic CM factors: the actual Hodge-circle test]
\label{p:CMthreepair}
Let $B,C$ be nonisogenous simple CM threefolds with endomorphism fields
$F_B,F_C$. Their joint Hodge group is $U_{F_B}\times U_{F_C}$ unless
the fields contain a common imaginary quadratic field $k$.
In that case orient their $k$-signatures oppositely; the joint group is
\[
 \ker\bigl(N_{F_B/k}N_{F_C/k}:U_{F_B}\times U_{F_C}\longrightarrow U_k\bigr).
\]
\end{lemma}
\begin{proof}
We work with rational cocharacter modules, on which the absolute Galois
group acts through a finite quotient. The joint Hodge torus is determined
by the rational span of the Galois orbit of its Hodge cocharacter.
For either sextic field this cocharacter has coordinates
$f\in\{1,-1\}^3$ on three conjugate embedding pairs, and its orbit spans
the full three-dimensional module, by \cite[\S2.3]{MZ}.
All finite-group modules below are semisimple over $\Q$.
Thus a failure of a product must identify isomorphic irreducible
constituents, and the corresponding projections of the two Hodge
cocharacters must satisfy the proposed graph relation.

First suppose a sextic CM field $F$ contains no imaginary quadratic
subfield. Its module is irreducible by \cite[Proposition (2.4)]{MZ}.
Let $G$ be its faithful signed permutation group and $P\subset S_3$
its transitive permutation quotient, so $P=C_3$ or $S_3$.
The sign subgroup $N=G\cap\{\pm1\}^3$ contains complex conjugation
$c=(-1,-1,-1)$. As an $\mathbf F_2[P]$-submodule it is either
$\langle c\rangle$ or all of $\mathbf F_2^3$: the permutation module
is the direct sum of the constant line and the irreducible
sum-zero plane for $C_3$.
If $N=\langle c\rangle$, the product of the three signs is a
character $\epsilon:G\to\{\pm1\}$ with $\epsilon(c)=-1$.
It is trivial on the stabilizer of a positive embedding coordinate.
To check this, that stabilizer projects either to the identity or to
a transposition fixing that coordinate. In the latter case a negative
product of the other two signs would square to a nonconstant sign
vector, contradicting $N=\langle c\rangle$.
The quadratic field cut out by $\epsilon$ would therefore lie in $F$
and be imaginary, a contradiction. Hence $N=\{\pm1\}^3$ and
$G=N\rtimes P$.

The subgroup $N$ is characteristic: it is the largest normal
$2$-subgroup, since $P$ has none. Its three eigenlines in $\Q^3$ are
precisely the coordinate lines. If two such irreducible modules were
identified, their faithful Galois quotients would agree, and an
intertwiner would permute these lines. Transitivity forces its three
nonzero coefficients to have the same absolute value. Compatibility
with the two vectors in $\{\pm1\}^3$ forces that value to be one.
Consequently it identifies the actual weight sets $\{\pm e_1,\pm
 e_2,\pm e_3\}$ of $H^1(B)$ and $H^1(C)$ and their Hodge circles.
Their complex representations on the joint torus would then have
nonzero intertwiners. Formation of this intertwiner space commutes with
extension from $\Q$ to $\C$, yielding a nonzero rational Hodge morphism
$H^1(B)\to H^1(C)$. It would give an isogeny, a contradiction.
There is therefore no common quotient in this case. Nor can this
irreducible three-dimensional module share a quotient with the
$1+2$ decomposition considered next.

Now let $F=kF_0$, with $F_0$ real cubic. The field $k$ is unique,
since two imaginary quadratic subfields would yield a real quadratic
subfield of $F_0$. If $\Omega$ is the normal closure of $F_0$, then
$\operatorname{Gal}(k\Omega/\Q)=C_2\times P$, where
$P=C_3$ or $S_3$. Write $V=\{(x_1,x_2,x_3):\sum x_i=0\}$ for the
standard rational plane, and $\chi_k$ for the quadratic character.
The cocharacter module is
\[
 \chi_k\ \oplus\ (\chi_k\otimes V).
\]
The first constituent can be shared only when the quadratic fields
agree. A simple CM type is not constant, since a constant type is
induced from $k$. Reverse the $k$-orientation to arrange $\sum f_i=1$.
Its projection to $V$ is then $\frac23 v_j$, where
$v_1=(-2,1,1)$ and $v_2,v_3$ are its coordinate permutations.

The two-dimensional constituent is faithful on $C_2\times P$.
For $P=C_3$ its endomorphism field is cyclic sextic; all nonconstant
CM types are conjugate under its field automorphisms, up to complex
conjugation. Two factors sharing this constituent would have the same
normal closure and hence the same field, and would be isogenous.
For $P=S_3$, equality of faithful constituents gives the same normal
closure $L$ and the same maximal real subfield $\Omega=L^{\langle c\rangle}$.
The real cubic subfields of $\Omega$ are conjugate, so identify their
permutation modules. There are exactly two imaginary quadratic
subfields of $L$, with characters differing by $\operatorname{sgn}_P$.
If the two $k$'s coincide, any intertwiner on $V$ is scalar because
$\End_{\Q[P]}(V)=\Q$. Two vectors among $v_1,v_2,v_3$ are proportional
only when equal. A compatible graph therefore identifies the CM types
under a field isomorphism and makes the factors isogenous.

If the two $k$'s differ, an intertwiner is, up to a nonzero rational
scalar, $J=\rho-\rho^{-1}$, with $\rho=(123)$ on $V$.
Indeed $J$ commutes with $\rho$ and anticommutes with each
transposition, so it intertwines $V$ and $V\otimes\operatorname{sgn}_P$;
these irreducible modules have one-dimensional intertwiner space.
Each $Jv_j$ has one zero coordinate, whereas every $v_i$ has no zero
coordinate. Hence no such intertwiner carries one projected Hodge
cocharacter to the other. The orbit spans both copies, and there is
no two-dimensional coupling. This check is necessary even though the
abstract quotient modules are isomorphic.

Only the one-dimensional constituent for a common $k$ can remain
coupled. Its Hodge-circle exponents are $1$ and $-1$ after opposite
orientation. Thus its relation is precisely the displayed product of
norms. Over $\C$ this is a primitive relation on six torus coordinates,
so its kernel is connected. The orbit-span calculation gives equality
of Lie algebras, hence equality of the connected groups.
\end{proof}

\begin{proposition}[Threefold strata]\label{p:N4strata}
Let $B,C$ be simple threefolds. The Hodge conjecture holds for $B\times C$
in each of the following cases: $B,C$ are isogenous; at least one has
Type I; or one has imaginary quadratic endomorphism algebra and the
other is CM. It also holds if their joint group is a product, or is the
full determinant kernel
\[
 \{(g_1,g_2)\in U_k(B)\times U_k(C):\det_k(g_1)\det_k(g_2)=1\},
\]
where their $k$-signatures have been oriented as $(2,1)$ and $(1,2)$.
The last displayed group equality is a hypothesis of this intermediate
proposition; Theorem~\ref{p:N4complete} below removes the need to
assume it in the classification of the whole row.
\end{proposition}
\begin{proof}
For isogenous factors the prime-dimensional theorem applies to $B^2$.
For nonisogenous factors with a Type I factor,
\cite[Proposition (2.4), Lemma (3.4)]{MZ} gives the product group;
its invariants are products of divisor classes.

Suppose $\End^0(B)=k$ is imaginary quadratic and $\End^0(C)=F$ is
sextic CM. Their individual groups are $U_k(B)$ and $U_F$.
Any common quotient must factor through $\det_k:U_k(B)\to U_k$.
A rank-one quotient of $U_F$ is isogenous to $U_{k'}$ for an imaginary
quadratic subfield $k'\subset F$, by \cite[Lemma (3.7)]{MZ} applied
to its connected kernel. Nontrivial coupling therefore requires
$k\subset F$. If this fails, the group is a product.
If $k\subset F$, both determinant Hodge circles have exponents $\pm1$:
a pure $k$-signature on $C$ would contradict simplicity.
After reversing one action the joint group equals
\[
 H=\{(g,t):\det_k(g)N_{F/k}(t)=1\}.
\]
Indeed it contains the Hodge circle. Its kernel over $U_F$ is $\SU_3$,
which also lies in the actual joint derived group, since a torus has
trivial derived group and projection onto $\SU_3$ is surjective.
Consequently both groups have dimension $3+8=11$ and coincide.
The kernel is connected, as is seen from the primitive determinant
relation over $\C$.

For invariants, first take $\SL_3$ on the $B$ coordinates and
$\ker N_{F/k}$ on the $C$ coordinates. In each single factor the only
invariants besides pairings are its degree-three volumes, by the
calculation in Lemma~\ref{p:prime-ell}. Central neutrality allows the
product of one volume from each factor, with matching determinant
weights. These two complex lines form $W_k(B\times C)$; all other
invariants are divisor products. Its $k$-rank-three summand $B$ is odd,
so Lemma~\ref{p:oddscale} proves algebraicity.
For the displayed determinant-kernel hypothesis the same calculation
uses $\SL_3\times\SL_3$ instead and gives the identical conclusion.
\end{proof}
\begin{theorem}[The complete $3+3$ row]\label{p:N4complete}
For any two simple abelian threefolds $B,C$, all Hodge classes on
$B\times C$ are algebraic. If they are nonisogenous, their joint group
is a product except possibly for the single determinant relation
associated with a common imaginary quadratic field. In that exceptional
case the Hodge ring on the sixfold is generated by divisors and
$W_k(B\times C)$.
\end{theorem}
\begin{proof}
Proposition~\ref{p:N4strata} treats isogenous factors, Type I factors,
and a unitary factor paired with a CM factor. Lemma~\ref{p:CMthreepair}
treats two CM factors. It remains to justify the group assertion for
two nonisogenous non-CM factors with imaginary quadratic endomorphism
fields $k_B,k_C$.

Their individual groups are $U_3$ and their derived groups are
$\Q$-simple forms of $\SL_3$, by \cite[\S2.3]{MZ}.
The quadratic field recovered from the Galois action on the two
vertices of their $A_2$ Dynkin diagram is $k_B$, respectively $k_C$.
If these fields differ, the derived groups have no common simple
quotient, and their determinant tori have distinct quadratic Galois
characters. Goursat's lemma gives the product group.

Suppose $k_B=k_C=k$ and orient both signatures as $(2,1)$ for this
paragraph. The joint group lies in $\det_k g_B=\det_k g_C$ because
the two determinant Hodge circles have exponent one. A failure of
independence of the derived groups would identify their $\SL_3$
factors. On the chosen $k$-eigenspaces their standard three-dimensional
modules would then be either equivalent or dual. The normalized Hodge
circle has weights $(1,1,-1)$; its traceless part has weights
$(2/3,2/3,-4/3)$. Duality negates this multiset and cannot preserve it.
In the equivalent case the central scalar actions also agree, since
the determinant relation makes their traces equal. Thus the two
modules are equivalent for the full connected joint group. Complex
conjugation gives the same conclusion for the other $k$-eigenspaces.
As above, base change for $\Hom$ gives a nonzero rational Hodge
morphism, contradicting nonisogeny. The derived group is consequently
the full product. The determinant-circle image generates the diagonal
norm-one torus. The actual group equals the connected determinant
kernel by dimension; reversing the action on $C$ writes its relation
as $\det_k g_B\det_k g_C=1$.

For the remaining CM--CM case, Lemma~\ref{p:CMthreepair} gives the
product or the same norm-determinant kernel. In a product group the
prime-dimensional theorem makes every invariant divisor-generated.
In a determinant kernel, first take the kernel of the individual
determinant on each factor. Over $\C$ this is $\SL_3$ in a unitary
factor and $\{(z_1,z_2,z_3):z_1z_2z_3=1\}$ in a CM factor.
On one copy of either factor its exterior invariants are generated by
degree-two pairings and the two degree-three volumes. In the CM case
this follows because an exponent vector in $\{-1,0,1\}^3$ is trivial
on the norm kernel only when its three coordinates agree; the unitary
case is the standard $\SL_3$ invariant calculation.
The residual determinant torus permits only the opposite-weight
product of a volume from each factor, besides products of pairings.
The two resulting degree-six lines form the rational plane
$W_k(B\times C)$. Its algebraicity follows from
Lemma~\ref{p:oddscale}, since each factor has $k$-rank three.
All other classes are products of divisors. Rational descent is as in
Lemma~\ref{p:fft}; this proves the theorem on the entire sixfold.
\end{proof}
\begin{example}
Choose $(B_1,\ldots,B_{2})$ jointly very generally in
$\prod_i\mathcal A_{d_i}(N)$ with $d_i=(3,3)$,
as constructed in Proposition~\ref{p:products}. Then each factor is
simple with $\End^0(B_i)=\Q$, and the factors are pairwise nonisogenous.
Their product is an actual sixfold of this partition, with
$\End^0(A)=\Q^{2}$ and the proved joint group above.
The Hodge-space dimensions in codimensions $0$ through $6$ are
$(1,2,3,4,3,2,1)$.
An explicit monomial basis over $\Q$ on $A$ is
$\theta_1^{a_1}\cdots\theta_{2}^{a_{2}}$ with
$0\leq a_i\leq d_i$; every such monomial is an intersection of divisors.
On powers add Poincare divisors between the copies of each individual
$B_i$. There are no cross-factor divisors or higher mixed invariants in
this verified jointly generic example.
\end{example}
\begin{example}[Two nonisogenous unitary threefolds]
Over $k=\Q(i)$ use Hermitian matrices
$h_B=\operatorname{diag}(1,1,-1)$ and
$h_C=\operatorname{diag}(1,-1,-1)$. Choose a jointly very general point
of the product of their two-dimensional PEL domains.
Derived monodromy is the independent product $\SU(h_B)\times\SU(h_C)$.
The determinant circle has coordinates $(1,-1)$, so the joint group
is exactly the determinant kernel in Proposition~\ref{p:N4strata}.
Its commutant is $k\times k$, proving that the factors are simple
and nonisogenous. The sixfold has split determinant $-1$, and
\[
 P_A(t)=(1+t+t^2+t^3)^2+2t^3,
 \qquad (1,2,3,6,3,2,1).
\]
The two polarizations generate the ordinary part and $W_k(A)$ supplies
the two algebraic middle classes. The group hypothesis has been verified
for this family; it is not inferred merely from the partition $3+3$.
\end{example}
\begin{example}[Two CM threefolds with a mixed Weil plane]
Let $F_1=\Q(\zeta_7+\zeta_7^{-1})$,
$F_2=\Q(\zeta_9+\zeta_9^{-1})$, and $K_j=\Q(i)F_j$.
Each $K_j$ is cyclic sextic. Choose signed CM types $(1,1,-1)$ and
$(-1,-1,1)$ relative to $\Q(i)$, respectively, and let $B,C$ be the
corresponding CM threefolds. These nonconstant types are primitive:
the only proper CM subfield is $\Q(i)$ and an induced type would be
constant. Thus the factors are simple with endomorphism fields $K_j$.
They are nonisogenous, since these fields are nonisomorphic (their
ramification at $7$ differs). Lemma~\ref{p:CMthreepair} gives the joint
rank-five norm kernel. The six opposite-coordinate pairings and the
two mixed volumes give
\[
 P_{B\times C}(t)=(1+t)^6+2t^3,
 \qquad (1,6,15,22,15,6,1).
\]
The twenty ordinary middle classes are divisor products; the two
additional classes form the algebraic plane $W_{\Q(i)}(B\times C)$.
This is a coupled CM example of the complete theorem, rather than
an independence assumption.
\end{example}
\section{N5: The \texorpdfstring{$3+2+1$}{3+2+1} case (\provedstatus)}\label{sec:N5}

\status{Proved for every sixfold in the $3+2+1$ row, including CM
factors and the exceptional threefold--elliptic interaction.}
\begin{theorem}
Let $A\sim B\times S\times E$, where $B,S$ are simple of dimensions
three and two and $E$ is elliptic. Then the Hodge conjecture holds for $A$.
Its Hodge ring is generated by divisors and the pullback of the Weil
plane on $B\times E$, when the matching imaginary-quadratic case occurs.
\end{theorem}
\begin{proof}
Lemma~\ref{p:surface-three} separates $\Hg(S)$ from $\Hg(B\times E)$.
The Kunneth invariant ring is therefore the tensor product of the two
factor Hodge rings. The surface ring is divisor-generated by the
prime-dimensional theorem. Lemma~\ref{p:prime-ell} with $g=3$ gives
exactly the stated generators on the fourfold $B\times E$, and proves
their algebraicity using the fourfold Weil theorem. Their exterior
products with surface divisor classes are algebraic, proving the result.
\end{proof}
\begin{example}
Choose $(B_1,\ldots,B_{3})$ jointly very generally in
$\prod_i\mathcal A_{d_i}(N)$ with $d_i=(3,2,1)$,
as constructed in Proposition~\ref{p:products}. Then each factor is
simple with $\End^0(B_i)=\Q$, and the factors are pairwise nonisogenous.
Their product is an actual sixfold of this partition, with
$\End^0(A)=\Q^{3}$ and the proved joint group above.
The Hodge-space dimensions in codimensions $0$ through $6$ are
$(1,3,5,6,5,3,1)$.
An explicit monomial basis over $\Q$ on $A$ is
$\theta_1^{a_1}\cdots\theta_{3}^{a_{3}}$ with
$0\leq a_i\leq d_i$; every such monomial is an intersection of divisors.
On powers add Poincare divisors between the copies of each individual
$B_i$. There are no cross-factor divisors or higher mixed invariants in
this verified jointly generic example.
\end{example}
\begin{example}[A fourfold Weil plane inside the $3+2+1$ row]\label{ex:N5weil}
Take a very general $k=\Q(i)$ threefold $B$ with Hermitian datum
$\operatorname{diag}(1,1,-1)$, the curve $E:y^2=x^3-x$ with opposite
$k$-action, and a very general simple surface $S$ with $\End^0(S)=\Q$.
Then $\Hg(B)=U(2,1)$, and Lemma~\ref{p:surface-three} separates $S$.
The Hodge ring is generated by divisors and the pullback of the
algebraic plane $W_k(B\times E)$. The dimension polynomial is
\[
 P_A(t)=(1+2t+4t^2+2t^3+t^4)(1+t+t^2),
 \qquad (1,3,7,8,7,3,1).
\]
Thus this example includes exceptional classes and their products
with divisors on the surface.
\end{example}
\section{N6: The \texorpdfstring{$3+1+1+1$}{3+1+1+1} case (\provedstatus)}\label{sec:N6}

\status{Proved for every sixfold in the $3+1+1+1$ row. Elliptic
factors may be isogenous or have CM.}
\begin{theorem}
If $A\sim B\times E_1\times E_2\times E_3$, where $B$ is a simple
threefold and the $E_i$ are elliptic, the Hodge conjecture holds for $A$.
The generators are divisors and pullbacks of fourfold Weil classes from
$B$ times a matching CM elliptic factor, as in Lemma~\ref{p:three-ell}.
\end{theorem}
\begin{proof}
Apply that lemma to $P=E_1\times E_2\times E_3$ and use isogeny
invariance. Grouping repeated elliptic classes is included in its proof,
so no independence assumption on the three elliptic curves is imposed.
\end{proof}
\begin{example}
Choose $(B_1,\ldots,B_{4})$ jointly very generally in
$\prod_i\mathcal A_{d_i}(N)$ with $d_i=(3,1,1,1)$,
as constructed in Proposition~\ref{p:products}. Then each factor is
simple with $\End^0(B_i)=\Q$, and the factors are pairwise nonisogenous.
Their product is an actual sixfold of this partition, with
$\End^0(A)=\Q^{4}$ and the proved joint group above.
The Hodge-space dimensions in codimensions $0$ through $6$ are
$(1,4,7,8,7,4,1)$.
An explicit monomial basis over $\Q$ on $A$ is
$\theta_1^{a_1}\cdots\theta_{4}^{a_{4}}$ with
$0\leq a_i\leq d_i$; every such monomial is an intersection of divisors.
On powers add Poincare divisors between the copies of each individual
$B_i$. There are no cross-factor divisors or higher mixed invariants in
this verified jointly generic example.
\end{example}
\begin{example}[Repeated matching elliptic factors]
Use the $k=\Q(i)$ threefold $B$ and matching elliptic curve $E$
from Example~\ref{ex:N5weil}, and put $A=B\times E^3$. Pairings contribute
$(1+t+t^2+t^3)(1+9t+9t^2+t^3)$.
A threefold volume must be paired with elliptic weight $\mp1$.
Choosing $a$ positive and $a+1$ negative coordinates in the three
elliptic copies gives $2\binom3a\binom3{a+1}$ such invariants,
for $a=0,1,2$. Hence
\[
 P_A(t)=(1+t+t^2+t^3)(1+9t+9t^2+t^3)+6t^2+18t^3+6t^4,
\]
with dimensions $(1,10,25,38,25,10,1)$.
Lemma~\ref{p:three-ell} supplies algebraic generators: divisors and
pullbacks of $W_k(B\times E)$ under the $k$-linear maps $E^3\to E$.
\end{example}
\section{N7: The \texorpdfstring{$2+2+2$}{2+2+2} case (\provedstatus)}\label{sec:N7}

\status{The Hodge conjecture is proved for the full $2+2+2$ row by
an established product-of-surfaces theorem. A uniform explicit list of
all joint CM character relations is not claimed here.}
\begin{theorem}
Every $A\sim S_1\times S_2\times S_3$, with $S_i$ simple abelian
surfaces, satisfies the rational Hodge conjecture, as does every power.
\end{theorem}
\begin{proof}
Ram\'on Mar\'i's theorem in Theorem~\ref{p:lowdim}(3) applies to the
three abelian surfaces. A power of this product is again a product of
abelian surfaces, so the same theorem applies. Finally use isogeny
invariance. The theorem permits arbitrary CM fields, isogenies among
factors, and common Hodge-group quotients; pairwise group independence
is not used as a substitute for its hypothesis.
\end{proof}
\begin{example}[An explicit repeated CM surface]\label{ex:CMsurface}
Let $S=J(C)$ for $C:y^2=x^5-1$. The curve has genus two, and
$(x,y)\mapsto(\zeta_5x,y)$ gives CM by $F=\Q(\zeta_5)$ with type
$\{1,2\}$. Its signed Galois orbit contains $(1,1)$ and $(1,-1)$,
so $\Hg(S)=U_F$ has rank two. The four distinct eigencharacters have
commutant $F$, proving $\End^0(S)=F$ and simplicity.
For $A=S^3$, each of its two conjugate character pairs has multiplicity
three. Neutral monomials choose the same number of each sign, so
\[
 P_A(t)=(1+9t+9t^2+t^3)^2,
 \qquad (1,18,99,164,99,18,1).
\]
They factor into degree-two pairings. Rational descent and Lefschetz
$(1,1)$ show that the eighteen rational divisor classes generate the
Hodge ring, including the Poincare classes between repeated factors.
Here $\End^0(A)=M_3(F)$.
\end{example}
\section{N8: The \texorpdfstring{$2+2+1+1$}{2+2+1+1} case (\provedstatus)}\label{sec:N8}

\status{Proved for every sixfold in the $2+2+1+1$ row, without
restrictions on endomorphisms or repetition of factors.}
\begin{theorem}
Every $A\sim S_1\times S_2\times E_1\times E_2$ in this row satisfies
the rational Hodge conjecture, as does every power.
\end{theorem}
\begin{proof}
Apply Theorem~\ref{p:lowdim}(3), since all four factors have dimension
at most two. This includes both isogenous and nonisogenous factors;
$A^n$ has the same dimension bound on its simple factors. Isogeny
invariance transfers the assertion to $A$.
\end{proof}
\begin{example}[A surface square and an elliptic square]
Let $S$ be the explicit CM surface of Example~\ref{ex:CMsurface} and
$E:y^2=x^3-x$. The field $\Q(i)$ does not embed in $\Q(\zeta_5)$,
whose only quadratic subfield is real. Thus
$\Hg(S\times E)=U_{\Q(\zeta_5)}\times U_{\Q(i)}$, by
Theorem~\ref{p:lowdim}(2). For $A=S^2\times E^2$, the three independent
conjugate character pairs each have multiplicity two, and
\[
 P_A(t)=(1+4t+t^2)^3,
 \qquad (1,12,51,88,51,12,1).
\]
Opposite-character pairings give all invariants and therefore twelve
rational divisor generators. This verifies a repeated CM case directly,
as well as by the general theorem.
\end{example}
\section{N9: The \texorpdfstring{$2+1+1+1+1$}{2+1+1+1+1} case (\provedstatus)}\label{sec:N9}

\status{Proved for every sixfold in the $2+1+1+1+1$ row, including
all CM and repeated-factor cases.}
\begin{theorem}
Every $A\sim S\times E_1\times E_2\times E_3\times E_4$ in this row
satisfies the rational Hodge conjecture, as does every power.
\end{theorem}
\begin{proof}
All simple factors have dimension at most two, so apply
Theorem~\ref{p:lowdim}(3) and isogeny invariance, also to every power.
Alternatively, the simple surface has no imaginary quadratic subfield
in the center of its endomorphism algebra; \cite[Proposition (3.8)]{MZ}
separates it from the elliptic isogeny classes. Divisor generation on
surface powers and on elliptic products then gives the same conclusion.
\end{proof}
\begin{example}[One CM surface and four equal elliptic factors]
Use $S$ and $E$ from the preceding example and put $A=S\times E^4$.
Their independent Hodge groups give
\[
 P_A(t)=(1+t)^2(1+16t+36t^2+16t^3+t^4),
 \qquad (1,18,69,104,69,18,1).
\]
Two divisors from $S$ and sixteen from $E^4$ generate the Hodge ring;
the latter include graph and Poincare classes associated with the
$\Q(i)$-linear homomorphisms between elliptic copies.
All classes are algebraic intersections of these divisors.
\end{example}
\section{N10: The \texorpdfstring{$1+1+1+1+1+1$}{1+1+1+1+1+1} case (\provedstatus)}\label{sec:N10}

\status{Full row: every product of six elliptic curves, with arbitrary
CM behavior and isogenous repetitions, on every power.}
\begin{theorem}
If a sixfold $A$ is isogenous to a product of six elliptic curves, every
Hodge class on every power of $A$ is generated by divisor classes.
\end{theorem}
\begin{proof}
Reduce by isogeny and collect equal isogeny classes. On powers one has
$A^n\sim\prod_i E_i^{m_i n}$ with the $E_i$ pairwise nonisogenous.
A non-CM elliptic curve has Hodge group $\SL_2$.
A subdirect product coupling two non-CM $\SL_2$ factors would be a twisted
diagonal by Goursat's lemma. Automorphisms of $\mathfrak{sl}_2$ are inner,
so their standard Hodge representations would be isomorphic. The
corresponding Hodge isomorphism of $H^1$ comes from a rational
homomorphism of elliptic curves and would make them isogenous.
This is excluded. The non-CM joint group is therefore the full product
of its $\SL_2$ factors.
For a CM elliptic curve with imaginary quadratic field $K_i$, its Hodge
group is the one-dimensional norm-one torus of $K_i$.
Distinct isogeny classes of CM elliptic curves have distinct imaginary
quadratic fields. In the anti-conjugate character space the one-dimensional
rational Galois modules for these tori are the distinct quadratic characters
of these fields. Their direct sum is multiplicity-free. The signed Hodge
cocharacter has a nonzero component in every summand. Applying the
rational character projectors to its Galois orbit therefore spans every
summand. The joint CM Hodge group is the product of these tori.
Finally no connected common quotient of a semisimple product of $\SL_2$
groups and a torus is nontrivial: a quotient of the former has semisimple
Lie algebra and a quotient of the latter is abelian. Goursat therefore
separates the CM and non-CM contributions.
Repeated copies are acted on diagonally by their isogeny-class group.
For non-CM copies, symplectic invariant theory gives degree-two pairings.
For CM copies, independent norm-one characters give exactly opposite-weight
pairings, also in degree two. The rational degree-two invariant space is
$\NS(A^n)_\Q$. Rational descent as in Lemma~\ref{p:fft} proves the conclusion.
Graphs of homomorphisms between isogenous copies and the corresponding
Poincare divisors provide its geometric generators.
\end{proof}
\begin{example}
Let $E$ be $y^2=x^3-x$; then $\End^0(E)=\Q(i)$ and take $A=E^6$.
The diagonal norm-one torus has characters $\chi,-\chi$ on two
six-dimensional multiplicity spaces. In degree $2p$, an invariant selects
$p$ vectors from each space, so
\[
 \dim_\Q\B^p(A)=\binom6p^2,\qquad
 P_A(t)=1+36t+225t^2+400t^3+225t^4+36t^5+t^6.
\]
The divisor space is the $36$-dimensional space of Hermitian $6\times6$
matrices over $\Q(i)$, represented by factor polarizations and Poincare
classes associated with the endomorphisms $1$ and $i$ between factors.
Its products generate every Hodge class. Here
$\End^0(A)=M_6(\Q(i))$; this differs from the independent elliptic example
but belongs to the same numerical partition. The distinction illustrates
why a dimension partition alone cannot fix Hodge-space dimensions.
\end{example}

\section*{Conditional conclusion and external dependencies}
\addcontentsline{toc}{section}{Conditional conclusion and external dependencies}
\begin{pretheorem}[Conditional sixfold conclusion]
\label{p:conditional-sixfold}
Assume the split-eightfold and CM assertions stated in
Theorem~\ref{p:openai-inputs}, the specified full Kuga--Satake
algebraicity assertion used in Theorem~\ref{p:N3KS}, and the
auxiliary-spin assertion (Aux) stated before
Theorem~\ref{p:N2auxiliary}, together with the classical external
results cited in the corresponding rows. Then every complex
abelian sixfold satisfies the rational Hodge conjecture.
\end{pretheorem}
\begin{proof}
Poincare complete reducibility and the endomorphism classification
reduce, up to isogeny, to the twenty-one rows S1--S11 and N1--N10.
The split-eightfold input completes S7, S8 and S10 through
Theorem~\ref{p:all-sixfold-weil}; the other simple rows and N1,
N4--N10 are treated in their sections.
Corollaries~\ref{p:N2CM-new} and~\ref{p:N3CM-new} handle the
remaining CM-fourfold strata. The non-CM residual models of
N2 and N3 are covered by Theorems~\ref{p:N2auxiliary}
and~\ref{p:N3KS}, respectively. The independent and other
nonresidual strata of those rows were already proved.
Rational algebraicity is invariant under isogeny, so these
arguments cover every sixfold.
\end{proof}

For clarity, the new external dependencies are as follows.
\begin{center}
\begin{tabular}{>{\raggedright\arraybackslash}p{0.44\textwidth}
                >{\raggedright\arraybackslash}p{0.46\textwidth}}
\toprule
External assertion & Application in this manuscript\\
\midrule
Split-eightfold Weil algebraicity \cite{OpenAIEightfold}
 & S7, S8, S10, via stabilization and contraction\\
CM abelian Hodge theorem \cite{OpenAICM}
 & CM-fourfold strata; CM transfer in residual N2\\
Specified complete KS correspondence \cite{OpenAIKS}
 & Residual N3, via algebraic matrix coefficients\\
Auxiliary degree-four spin tensor \cite{OpenAIK3Products}
 & Residual N2, via determinant stabilization\\
\bottomrule
\end{tabular}
\end{center}

These inputs have different logical roles. In particular the CM
theorem does not apply directly to a product with a non-CM fourfold.
Theorem~\ref{p:N2auxiliary} first obtains a cycle on $B\times D$
from (Aux), and only then invokes the CM theorem on $D\times P$.
Theorem~\ref{p:N3KS} uses the specified KS map through algebraic
weight-one factor maps; no abstract K3 Hodge isometry is declared
algebraic. The auxiliary-spin input is not deduced here from the
ordinary KS input.

All twenty-one sixfold rows are now covered relative to the stated
inputs. The all-powers conclusions are those explicitly proved in
the individual sections, including the residual N2 and N3 models;
the sixfold conclusion alone does not assert an unrestricted
all-powers theorem. The new external preprints and this manuscript
still require independent mathematical review. In particular,
the conditional conclusion is not an independent proof of the
external assertions on which it depends.

\clearpage
\begin{thebibliography}{99}
\addcontentsline{toc}{section}{References}
\bibitem{Borselli}
A. P. Borselli, \emph{Galois Groups of CM Fields in Degrees 24, 28,
and 30}, Ph.D. thesis, Lehigh University, 2017, \S2.2.2.
\href{https://preserve.lehigh.edu/lehigh-scholarship/graduate-publications-theses-dissertations/theses-dissertations/galois-groups-cm}{Lehigh Preserve: official thesis record}.
\bibitem{Moonen}
B. Moonen, \emph{Notes on Mumford--Tate groups}, Centre Emile Borel,
Paris, March 1999, especially \S\S2--3.
\url{https://www.math.ru.nl/~bmoonen/Lecturenotes/CEBnotesMT.pdf}.
\bibitem{MZ}
B. Moonen and Yu. G. Zarhin, \emph{Hodge classes on abelian varieties of
low dimension}, Math. Ann. \textbf{315} (1999), 711--733.
\url{https://arxiv.org/abs/math/9901113}.
\bibitem{MilneShimura}
J. S. Milne, \emph{Introduction to Shimura Varieties}, revised notes,
Theorem 3.3(b) and the PEL moduli interpretation.
\url{https://www.jmilne.org/math/xnotes/svi.pdf}.
\bibitem{MoonenCyclic}
B. Moonen, \emph{Special subvarieties arising from families of cyclic
covers of the projective line}, Doc. Math. \textbf{15} (2010), 793--819,
Lemma 2.7, Table 1 and Proposition 4.7.
\url{https://arxiv.org/abs/1006.3400}.
\bibitem{RamonMari}
J. J. Ram\'on Mar\'i, \emph{On the Hodge conjecture for products of
certain surfaces}, Collect. Math. \textbf{59} (2008), 1--26.
The product-of-abelian-surfaces result is Proposition 2.18 in the
2005 preprint (Theorem 3.15 in the published numbering).
\url{https://arxiv.org/abs/math/0505357}.
\bibitem{Vial}
C. Vial, \emph{Generic cycles, Lefschetz representations, and the generalized
Hodge and Bloch conjectures for abelian varieties}, especially \S3.2.
\url{https://arxiv.org/abs/1803.00857}.
\url{https://www.math.uni-bielefeld.de/~vial/SymAbelian.pdf}.
\bibitem{MilneLefschetz}
J. S. Milne, \emph{Lefschetz classes on abelian varieties},
Duke Math. J. \textbf{96} (1999), 639--675, \S5;
see also the author's erratum concerning Proposition 5.9.
\url{https://www.jmilne.org/math/articles/1999a.html}.
\bibitem{MilneTateStandard}
J. S. Milne, \emph{The Tate and standard conjectures for certain
abelian varieties}, arXiv:2112.12815, Corollary 1.11.
\url{https://www.jmilne.org/math/articles/tcav2.pdf}.
\bibitem{YuZheng}
C. Yu and Z. Zheng, \emph{Commensurability among Deligne--Mostow
monodromy groups}, arXiv:2411.09682v2 (2025), Proposition 7.4.
\url{https://arxiv.org/abs/2411.09682v2}.
\bibitem{Markman}
E. Markman, \emph{Cycles on abelian $2n$-folds of Weil type from secant
sheaves on abelian $n$-folds}, arXiv:2502.03415v2 (2025),
Theorem 1.5.1 and Corollary 1.6.1.
\url{https://arxiv.org/abs/2502.03415v2}.
\bibitem{Yanai}
H. Yanai, \emph{On the index of degeneracy of a CM abelian variety},
J. Th\'eor. Nombres Bordeaux \textbf{27} (2015), 815--820, \S1.
\url{https://www.numdam.org/articles/10.5802/jtnb.925/}.
\bibitem{MarkmanSurvey}
E. Markman, \emph{Secant sheaves and Weil classes on abelian varieties},
arXiv:2509.23403v2 (2026), Theorems 1.2, 1.4 and Corollary 1.3.
\url{https://arxiv.org/abs/2509.23403v2}.
\bibitem{OpenAIEightfold}
OpenAI, \emph{Algebraicity of Weil classes on split abelian eightfolds},
OpenAI Math Release preprint, September 18, 2026, Theorem 1.1;
public release October 6, 2026.
\href{https://github.com/openai/math/blob/main/preprints/Algebraicity-of-Weil-classes-on-split-abelian-eightfolds-September-18-2026/README.md}{Official manuscript and release record}.
\bibitem{OpenAICM}
OpenAI, \emph{The rational Hodge conjecture for CM abelian varieties},
OpenAI Math Release preprint, September 30, 2026, Theorem 1.1;
public release October 6, 2026.
\href{https://github.com/openai/math/blob/main/preprints/The-rational-Hodge-conjecture-for-CM-abelian-varieties-September-30-2026/README.md}{Official manuscript and release record}.
\bibitem{OpenAIWeilPowers}
OpenAI, \emph{Weil classes and Hodge classes on abelian powers},
OpenAI Math Release preprint, September 30, 2026, Lemma 3.2;
public release October 6, 2026.
\href{https://github.com/openai/math/blob/main/preprints/Weil-classes-and-Hodge-classes-on-abelian-powers-September-30-2026/README.md}{Official manuscript and release record}.
\bibitem{OpenAIKS}
OpenAI, \emph{Algebraicity of Kuga--Satake Correspondences for K3 Surfaces},
OpenAI Math Release preprint, October 3, 2026, Theorem 1.1;
public release October 6, 2026. Used as the specified complete
even-Clifford correspondence input in Theorem~\ref{p:N3KS}.
\href{https://github.com/openai/math/blob/main/preprints/Algebraicity-of-Kuga-Satake-Correspondences-for-K3-Surfaces-October-3-2026/README.md}{Official manuscript and release record}.
\bibitem{HuybrechtsK3}
D. Huybrechts, \emph{Lectures on K3 Surfaces}, Cambridge University
Press, 2016; author lecture-note version, Chapters 4, 6 and 14.
\url{https://www.math.uni-bonn.de/people/huybrech/K3Global.pdf}.
\bibitem{OpenAIK3Products}
OpenAI, \emph{The rational Hodge conjecture for products of K3 surfaces},
preprint, October 4, 2026, section \emph{Degree-four tensors in spin
families}, auxiliary-spin proposition.
Source pinned at commit
\texttt{adc7f1241b42e322a6451854ab7e4b4c146bf78a};
the exact assertion used is (Aux) preceding
Theorem~\ref{p:N2auxiliary}.
\href{https://github.com/openai/math/blob/adc7f1241b42e322a6451854ab7e4b4c146bf78a/preprints/The-rational-Hodge-conjecture-for-products-of-K3-surfaces-October-4-2026/build/manuscript/auxiliary.tex}{Pinned source of the auxiliary-spin proposition}.
\end{thebibliography}
\end{document}
